% Champ magnétique, force de Lorentz et force de Laplace — Physique Tle D — Cours signé OnBuch+
% Programme officiel MINESEC. Compiler avec : tectonic P04-champ-magnetique-forces-magnetiques.tex
\newcommand{\DOCMATIERE}{Physique}
\newcommand{\DOCNIVEAU}{Tle D}
\newcommand{\DOCMODULE}{Module 2 : Mouvements et interactions — évolutions temporelles des systèmes mécaniques}
\newcommand{\DOCLECON}{P04}
\newcommand{\DOCTITRE}{Champ magnétique, force de Lorentz et force de Laplace}
% OnBuch+ — préambule « cours signés » ENRICHI (compilé avec tectonic / XeLaTeX)
% Système visuel « Pop » d'OnBuch+ : fond crème (jamais blanc), bordures encre
% épaisses, ombres « dures » décalées, titres Archivo Black en capitales.
% Version MATHÉMATIQUES : graphes (pgfplots/tikz), boîtes pédagogiques
% (définition, propriété, méthode, attention, le savais-tu, exercice résolu,
% exercice, TP, bilan), page de garde et signature OnBuch+ ancrée en bas.
%
% Macros attendues dans main.tex AVANT \input{preamble} :
%   \DOCMATIERE  \DOCNIVEAU  \DOCMODULE  \DOCLECON (numéro)  \DOCTITRE
\documentclass[11pt]{article}

\usepackage{fontspec}
\usepackage[french]{babel}
\usepackage[a4paper,margin=2cm,top=3.4cm,bottom=2.8cm,headheight=1.4cm,footskip=1.3cm]{geometry}
\usepackage{amsmath,amssymb,amsfonts}
\usepackage{mathtools} % \xmapsto, \coloneqq... souvent utilisés par les modèles
\usepackage{cancel} % simplifications barrées (bilans de forces)
% \abs{z} (module, valeur absolue) et \norm{u} (norme), utilisés par les modèles
\providecommand{\abs}{}\let\abs\relax\DeclarePairedDelimiter{\abs}{\lvert}{\rvert}
\providecommand{\norm}{}\let\norm\relax\DeclarePairedDelimiter{\norm}{\lVert}{\rVert}
\usepackage[table]{xcolor} % [table] : fournit aussi \rowcolors (alternance de lignes)
\usepackage{pagecolor}
\usepackage{tikz}
\usetikzlibrary{arrows.meta,positioning,calc,shapes.geometric,shapes.misc,shapes.arrows,decorations.pathmorphing,decorations.markings,patterns,shadows,mindmap,trees,fit,backgrounds,matrix,intersections}
\usepackage{pgfplots}
\usepackage[version=4]{mhchem} % équations nucléaires \ce{^{235}_{92}U}
\usepackage[european,siunitx]{circuitikz} % schémas électriques
\pgfplotsset{compat=1.18}
\usepgfplotslibrary{fillbetween} % aires entre courbes (calcul intégral)
\usepackage{enumitem}
\usepackage{fancyhdr}
% [most] charge toutes les bibliothèques tcolorbox (listings, minted, raster,
% documentation...) et gonfle inutilement la mémoire TeX sur un document long
% (~300 boîtes) : on ne charge que ce qui est réellement utilisé.
\usepackage[skins,breakable]{tcolorbox}
\usepackage{graphicx}
\usepackage{colortbl}
\usepackage{booktabs}
\usepackage{tabularx}
\usepackage{diagbox} % case d'en-tête diagonale des tableaux à double entrée
% Colonnes extensibles centrée / gauche / droite (C, L, R), souvent utilisées par les modèles
\newcolumntype{C}{>{\centering\arraybackslash}X}
\newcolumntype{L}{>{\raggedright\arraybackslash}X}
\newcolumntype{R}{>{\raggedleft\arraybackslash}X}
\usepackage{array}
\usepackage{multicol}
\usepackage{siunitx}
% parse-numbers=false : accepte aussi \SI{2,5 \times 10^{-3}}{...}
\sisetup{locale=FR,output-decimal-marker={,},group-minimum-digits=4,parse-numbers=false}
\DeclareSIUnit\division{div} % réglages d'oscilloscope (ms/div, V/div)
\usepackage{qrcode}
% \degree : commande fréquemment utilisée par les modèles pour un angle ou une
% température en dehors de \SI{}{} (non fournie par les paquets chargés).
\providecommand{\degree}{^{\circ}}
% \qed (fin de démonstration), utilisé par les modèles sans amsthm
\providecommand{\qed}{\hfill$\square$}
% \barre : double barre des valeurs interdites dans un tableau de variations (tkz-tab)
\providecommand{\barre}{\ensuremath{\|}}
% environnement proof (amsthm non chargé) utilisé par les modèles
\ifdefined\proof\else\newenvironment{proof}[1][Démonstration]{\par\noindent\textit{#1.}\ }{\qed\par}\fi
\usepackage{lastpage}
\usepackage{hyperref}
\hypersetup{hidelinks}

% ============================================================
% Palette « Pop » OnBuch+ — mêmes hex que la classe P (plus_ui.dart)
% ============================================================
\definecolor{poporange}{HTML}{F59321}
\definecolor{popdark}{HTML}{E07A0C}
\definecolor{popcream}{HTML}{FFF6E8}
\definecolor{popink}{HTML}{1C1714}
\definecolor{poppurple}{HTML}{7A5AE0}
\definecolor{popgreen}{HTML}{2E9E6A}
\definecolor{poppink}{HTML}{E0457B}
\definecolor{popblue}{HTML}{2F7DE1}
\definecolor{popgold}{HTML}{C98A12}
\definecolor{popmuted}{HTML}{8A7F76}
\definecolor{popline}{HTML}{EADFD2}
\definecolor{popgray}{HTML}{8A7F76}
\colorlet{popgrey}{popgray}
% Alias tolérants (le modèle nomme parfois les couleurs différemment)
\colorlet{popwhite}{white}
\colorlet{popblack}{popink}
\colorlet{popred}{poppink}
\colorlet{popyellow}{popgold}
% Teintes claires pour fonds de tableaux / zones
\colorlet{popblueL}{popblue!10!white}
\colorlet{poporangeL}{poporange!14!white}
\colorlet{popgreenL}{popgreen!12!white}
\colorlet{poppurpleL}{poppurple!12!white}
\colorlet{poppinkL}{poppink!12!white}
\colorlet{popgoldL}{popgold!14!white}

\pagecolor{popcream}

% ============================================================
% Typographie
% ============================================================
\defaultfontfeatures{Ligatures=TeX}
\setmainfont{PlusJakartaSans-Medium.ttf}[Path=../../../fonts/,BoldFont=PlusJakartaSans-Bold.ttf]
% Maths en sans-serif (Fira Math) avec chiffres et lettres droites de Plus
% Jakarta, pour que les formules se fondent dans le texte.
\usepackage[math-style=ISO,bold-style=ISO]{unicode-math}
\setmathfont{FiraMath-Regular.otf}[Path=../../../fonts/]
\setmathfont{PlusJakartaSans-Medium.ttf}[Path=../../../fonts/,range={up/num,up/{Latin,latin},\mathup}]
\usepackage{icomma} % virgule décimale sans espace en mode math
% Caractères mathématiques Unicode absents des polices de texte (Plus Jakarta,
% Archivo Black) : rendus par leur équivalent mathématique, en texte comme en
% formule — sinon ils s'affichent en carré vide (ex. « Arithmétique dans ℤ »).
\usepackage{newunicodechar}
\newunicodechar{ℝ}{\ensuremath{\mathbb{R}}}
\newunicodechar{ℤ}{\ensuremath{\mathbb{Z}}}
\newunicodechar{ℂ}{\ensuremath{\mathbb{C}}}
\newunicodechar{ℕ}{\ensuremath{\mathbb{N}}}
\newunicodechar{ℚ}{\ensuremath{\mathbb{Q}}}
\newunicodechar{𝔻}{\ensuremath{\mathbb{D}}}
\newunicodechar{α}{\ensuremath{\alpha}}
\newunicodechar{β}{\ensuremath{\beta}}
\newunicodechar{γ}{\ensuremath{\gamma}}
\newunicodechar{δ}{\ensuremath{\delta}}
\newunicodechar{ε}{\ensuremath{\varepsilon}}
\newunicodechar{θ}{\ensuremath{\theta}}
\newunicodechar{λ}{\ensuremath{\lambda}}
\newunicodechar{μ}{\ensuremath{\mu}}
\newunicodechar{σ}{\ensuremath{\sigma}}
\newunicodechar{τ}{\ensuremath{\tau}}
\newunicodechar{φ}{\ensuremath{\varphi}}
\newunicodechar{ψ}{\ensuremath{\psi}}
\newunicodechar{ω}{\ensuremath{\omega}}
\newunicodechar{Σ}{\ensuremath{\Sigma}}
% majuscules produites par \MakeUppercase dans les titres (α-aminés -> Α)
\newunicodechar{Α}{\ensuremath{\alpha}}
\newunicodechar{Β}{\ensuremath{\beta}}
\newunicodechar{Γ}{\ensuremath{\Gamma}}
\newunicodechar{Δ}{\ensuremath{\Delta}}
\newunicodechar{Ε}{\ensuremath{\varepsilon}}
\newunicodechar{Θ}{\ensuremath{\Theta}}
\newunicodechar{Λ}{\ensuremath{\Lambda}}
\newunicodechar{Μ}{\ensuremath{\mu}}
\newunicodechar{Τ}{\ensuremath{\tau}}
\newunicodechar{Φ}{\ensuremath{\Phi}}
\newunicodechar{Ψ}{\ensuremath{\Psi}}
\newunicodechar{Ω}{\ensuremath{\Omega}}
\newunicodechar{↦}{\ensuremath{\mapsto}}
\newunicodechar{∈}{\ensuremath{\in}}
\newunicodechar{∉}{\ensuremath{\notin}}
\newunicodechar{∧}{\ensuremath{\wedge}}
\newunicodechar{⇔}{\ensuremath{\Leftrightarrow}}
\newunicodechar{⟺}{\ensuremath{\Longleftrightarrow}}
\newunicodechar{⇒}{\ensuremath{\Rightarrow}}
\newunicodechar{⟹}{\ensuremath{\Longrightarrow}}
\newunicodechar{∘}{\ensuremath{\circ}}
\newunicodechar{∪}{\ensuremath{\cup}}
\newunicodechar{∩}{\ensuremath{\cap}}
\newunicodechar{≡}{\ensuremath{\equiv}}
\newunicodechar{⊂}{\ensuremath{\subset}}
\newunicodechar{⊥}{\ensuremath{\perp}}
\newunicodechar{∃}{\ensuremath{\exists}}
\newunicodechar{∀}{\ensuremath{\forall}}
\newunicodechar{∛}{\ensuremath{\sqrt[3]{\,}}}
\newunicodechar{∜}{\ensuremath{\sqrt[4]{\,}}}
\newunicodechar{⃗}{\ensuremath{{}^{\to}}}
\newunicodechar{ⁿ}{\ifmmode^{n}\else\textsuperscript{\ensuremath{n}}\fi}
\newunicodechar{ˣ}{\ifmmode^{x}\else\textsuperscript{\ensuremath{x}}\fi}
\newunicodechar{ᵇ}{\ifmmode^{b}\else\textsuperscript{\ensuremath{b}}\fi}
\newunicodechar{ʳ}{\ifmmode^{r}\else\textsuperscript{\ensuremath{r}}\fi}
\newunicodechar{ᵘ}{\ifmmode^{u}\else\textsuperscript{\ensuremath{u}}\fi}
\newunicodechar{ʸ}{\ifmmode^{y}\else\textsuperscript{\ensuremath{y}}\fi}
\newunicodechar{ᵃ}{\ifmmode^{a}\else\textsuperscript{\ensuremath{a}}\fi}
\newunicodechar{ᵉ}{\ifmmode^{e}\else\textsuperscript{\ensuremath{e}}\fi}
\newunicodechar{ᵏ}{\ifmmode^{k}\else\textsuperscript{\ensuremath{k}}\fi}
\newunicodechar{ᵖ}{\ifmmode^{p}\else\textsuperscript{\ensuremath{p}}\fi}
\newunicodechar{ᴺ}{\ifmmode^{N}\else\textsuperscript{\ensuremath{N}}\fi}
\newunicodechar{ᶜ}{\ifmmode^{c}\else\textsuperscript{\ensuremath{c}}\fi}
\newunicodechar{ʰ}{\ifmmode^{h}\else\textsuperscript{\ensuremath{h}}\fi}
\newunicodechar{ᵠ}{\ifmmode^{q}\else\textsuperscript{\ensuremath{q}}\fi}
\newunicodechar{ₙ}{\ifmmode_{n}\else\textsubscript{\ensuremath{n}}\fi}
\newunicodechar{ₐ}{\ifmmode_{a}\else\textsubscript{\ensuremath{a}}\fi}
\newunicodechar{ᵢ}{\ifmmode_{i}\else\textsubscript{\ensuremath{i}}\fi}
\newunicodechar{ᵣ}{\ifmmode_{r}\else\textsubscript{\ensuremath{r}}\fi}
\newunicodechar{ₖ}{\ifmmode_{k}\else\textsubscript{\ensuremath{k}}\fi}
\newunicodechar{ᵦ}{\ifmmode_{\beta}\else\textsubscript{\ensuremath{\beta}}\fi}
\newunicodechar{ₓ}{\ifmmode_{x}\else\textsubscript{\ensuremath{x}}\fi}
\newfontfamily\popdisplay{ArchivoBlack-Regular.ttf}[Path=../../../fonts/]
\newfontfamily\poplogo{SpaceGrotesk-Bold.ttf}[Path=../../../fonts/]
\newfontfamily\popsemi{PlusJakartaSans-SemiBold.ttf}[Path=../../../fonts/]
\newfontfamily\popxbold{PlusJakartaSans-ExtraBold.ttf}[Path=../../../fonts/]
\newfontfamily\popxboldls{PlusJakartaSans-ExtraBold.ttf}[Path=../../../fonts/,LetterSpace=3.0]

\setlist[enumerate]{leftmargin=1.7em,itemsep=5pt,topsep=4pt}
\setlist[itemize]{leftmargin=1.7em,itemsep=4pt,topsep=4pt,label={\color{poporange}\textbullet}}
\setlength{\parindent}{0pt}
\setlength{\parskip}{5pt}
\renewcommand{\arraystretch}{1.25}

% pgfplots : style Pop par défaut
\pgfplotsset{
  every axis/.append style={axis line style={popink,line width=1pt},tick style={popink},
    grid style={popline,line width=0.5pt},label style={font=\small},tick label style={font=\footnotesize}},
  popcurve/.style={poporange,line width=1.8pt,no markers,smooth},
}
% Même style disponible dans un \draw TikZ simple (hors pgfplots)
\tikzset{popcurve/.style={poporange,line width=1.8pt}}
\tikzset{popgrid/.style={popline,very thin}}

% ============================================================
% Pill / squiggle
% ============================================================
\newcommand{\poppill}[3]{%
  \tikz[baseline=-0.55ex]\node[draw=popink,line width=1.2pt,rounded corners=2.6pt,%
    fill=#2,inner xsep=6.5pt,inner ysep=3pt]{%
    \popxboldls\fontsize{7}{7}\selectfont\color{#3}\MakeUppercase{#1}};%
}
\newcommand{\popsquiggle}[1]{%
  \tikz[baseline=-0.4ex]{
    \draw[#1,line width=2.6pt,line cap=round]
      plot [smooth] coordinates {(0,0.09) (0.34,0.01) (0.68,0.21) (1.02,0.01) (1.36,0.10)};
  }%
}
% Mot-clé surligné (marqueur orange sous le texte)
\newcommand{\cle}[1]{{\popxbold\color{popdark}#1}}

% ============================================================
% En-tête / pied
% ============================================================
\pagestyle{fancy}
\fancyhf{}
\renewcommand{\headrulewidth}{0pt}
\renewcommand{\footrulewidth}{0pt}
\lhead{\raisebox{-0.15ex}{\poplogo\fontsize{13}{13}\selectfont\color{popink}OnBuch\color{poporange}+}}
\rhead{\poppill{\DOCMATIERE\ \textbullet\ \DOCNIVEAU\ \textbullet\ Leçon \DOCLECON}{white}{popink}}
\lfoot{\popxbold\fontsize{7.6}{7.6}\selectfont\color{popmuted}\MakeUppercase{Cours signé OnBuch+}\ \textbullet\ {\popsemi\DOCTITRE}}
\rfoot{\tikz[baseline=-0.6ex]\node[rounded corners=8.5pt,fill=popink,minimum height=17pt,minimum width=17pt,inner xsep=5pt,inner sep=0]{\popxbold\fontsize{8}{8}\selectfont\color{popcream}\thepage\,/\,\pageref*{LastPage}};}

% ============================================================
% Titres
% ============================================================
\newcommand{\chaptitre}[2]{%
  \begin{tcolorbox}[enhanced,colback=poporange,colframe=popink,boxrule=2.6pt,arc=11pt,
      drop shadow=popink,left=15pt,right=15pt,top=12pt,bottom=11pt,before skip=3pt,after skip=18pt]
    \poppill{#2}{popink}{popcream}\par\vspace{9pt}
    {\raggedright\hyphenpenalty=10000\exhyphenpenalty=10000\popdisplay\fontsize{22}{24}\selectfont\color{popcream}\MakeUppercase{#1}\par}
  \end{tcolorbox}
}
% Section numérotée : pastille orange avec le numéro + titre Archivo
\newcounter{popsec}
\newcommand{\coursec}[1]{%
  \stepcounter{popsec}\setcounter{popsub}{0}%
  \par\vspace{16pt}\noindent
  \begin{minipage}{\linewidth}
  \tikz[baseline=-0.7ex]\node[draw=popink,line width=1.6pt,fill=poporange,rounded corners=4pt,minimum size=22pt,inner sep=0]{\popdisplay\fontsize{12}{12}\selectfont\color{popcream}\thepopsec};%
  \hspace{8pt}{\popdisplay\fontsize{15}{17}\selectfont\color{popink}\MakeUppercase{#1}}\par
  \vspace{4pt}\noindent{\color{poporange}\rule{36pt}{3.6pt}}
  \end{minipage}\par\nopagebreak\vspace{8pt}
}
\newcounter{popsub}
\newcommand{\courssub}[1]{%
  \stepcounter{popsub}%
  \par\vspace{9pt}\noindent{\popxbold\fontsize{11.5}{13}\selectfont\color{popink}{\color{poporange}\thepopsec.\thepopsub}\hspace{6pt}#1}\par\nopagebreak\vspace{4pt}
}

% ============================================================
% Boîtes pédagogiques — même recette : fond blanc, bordure épaisse colorée,
% ombre dure encre, badge plein. Titre optionnel [#1].
% ============================================================
\tcbset{popbase/.style={breakable,enhanced,colback=white,boxrule=2.3pt,arc=9pt,
  shadow={3pt}{-3pt}{0pt}{popink},left=14pt,right=14pt,top=11pt,bottom=11pt,before skip=11pt,after skip=15pt,
  fontupper=\popsemi\fontsize{10.6}{14.5}\selectfont\color{popink},
  fontlower=\popsemi\fontsize{10.6}{14.5}\selectfont\color{popink}}}
% Coche dessinée (le glyphe ✓ n'existe pas dans les polices du document)
\newcommand{\popcheck}[1]{\tikz[baseline=-0.5ex]\draw[#1,line width=1.6pt,line cap=round,line join=round](-0.13,0.0)--(-0.04,-0.09)--(0.13,0.1);}
\providecommand{\coche}{\popcheck{popgreen}}
\newcommand{\popboxhead}[3]{\poppill{#1}{#2}{white}\ifx\relax#3\relax\else\hspace{6pt}{\popxbold\fontsize{10.6}{12}\selectfont\color{popink}#3}\fi\par\vspace{7pt}}

\newtcolorbox{aretenir}[1][]{popbase,colframe=popgreen,before upper={\popboxhead{À retenir}{popgreen}{#1}}}
\newtcolorbox{exemplebox}[1][]{popbase,colframe=poporange,before upper={\popboxhead{Exemple}{poporange}{#1}}}
\newtcolorbox{definition}[1][]{popbase,colframe=popblue,before upper={\popboxhead{Définition}{popblue}{#1}}}
\newtcolorbox{propriete}[1][]{popbase,colframe=poppurple,before upper={\popboxhead{Propriété}{poppurple}{#1}}}
\newtcolorbox{methode}[1][]{popbase,colframe=popgold,before upper={\popboxhead{Méthode}{popgold}{#1}}}
\newtcolorbox{attention}[1][]{popbase,colframe=poppink,before upper={\popboxhead{Attention}{poppink}{#1}}}
\newtcolorbox{experience}[1][]{popbase,colframe=popblue,colback=popblueL,before upper={\popboxhead{Expérience}{popblue}{#1}}}
% « Le savais-tu ? » : carte violette pleine (contexte, culture, Cameroun)
\newtcolorbox{savaistu}[1][]{popbase,colback=poppurple,colframe=popink,
  fontupper=\popsemi\fontsize{10.4}{14.2}\selectfont\color{white},
  before upper={\poppill{Le savais-tu\,?}{white}{poppurple}\ifx\relax#1\relax\else\hspace{6pt}{\popxbold\color{white}#1}\fi\par\vspace{7pt}}}
% Exercice résolu : énoncé en haut, \tcblower puis la solution sur fond crème
\newcounter{popexo}\newcounter{popexr}
\newtcolorbox{exoresolu}[1][]{popbase,colframe=poporange,colbacklower=popcream,
  segmentation style={popink,line width=1.2pt,dashed},
  before upper={\stepcounter{popexr}\popboxhead{Exercice résolu \thepopexr}{poporange}{#1}},
  before lower={\poppill{Solution}{popink}{popcream}\par\vspace{6pt}}}
% Exercice (non résolu en place) — la correction est dans la partie Corrigés
\newtcolorbox{exercice}[1][]{popbase,colframe=popink,
  before upper={\stepcounter{popexo}\popboxhead{Exercice \thepopexo}{popink}{#1}}}
% Corrigé
\newtcolorbox{corrige}[1][]{popbase,colframe=popgreen,colback=popgreenL,
  before upper={\popboxhead{Corrigé}{popgreen}{#1}}}
% Objectifs / prérequis / bilan
\newtcolorbox{objectifs}{popbase,colframe=popink,colback=poporangeL,
  before upper={\popboxhead{Objectifs de la leçon}{popink}{}}}
\newtcolorbox{prerequis}{popbase,colframe=popink,colback=popblueL,
  before upper={\popboxhead{Prérequis}{popblue}{}}}
\newtcolorbox{bilan}{popbase,colframe=popink,colback=white,boxrule=2.8pt,
  before upper={\popboxhead{Fiche bilan}{poporange}{L'essentiel à connaître pour l'examen}}}
% Figure encadrée avec légende
\newtcolorbox{popfigure}[1][]{enhanced,breakable=false,colback=white,colframe=popline,boxrule=1.4pt,arc=8pt,
  left=10pt,right=10pt,top=10pt,bottom=8pt,before skip=10pt,after skip=12pt,halign=center,
  lower separated=false,halign lower=center,
  fontlower=\popsemi\fontsize{9}{11.5}\selectfont\color{popmuted},
  #1}
\newcommand{\legende}[1]{\par\vspace{6pt}{\popsemi\fontsize{9}{11.5}\selectfont\color{popmuted}#1}}

% ============================================================
% Signature OnBuch+
% ============================================================
\newcommand{\poplogoinline}[1]{{\poplogo\fontsize{#1}{#1}\selectfont\color{popink}OnBuch\color{poporange}+}}
\newcommand{\signatureonbuch}{%
  \begin{tcolorbox}[enhanced,colback=popink,colframe=popink,boxrule=2.2pt,arc=10pt,
      drop shadow=poporange,left=14pt,right=14pt,top=10pt,bottom=10pt,before skip=0pt,after skip=0pt]
    \tikz[baseline=-0.7ex]\node[circle,fill=popgreen,draw=popcream,line width=1.3pt,minimum size=22pt,inner sep=0]{\popcheck{white}};%
    \hspace{9pt}\begin{minipage}[c]{0.62\linewidth}
      {\popxbold\fontsize{12}{13}\selectfont\color{popcream}Signé\ \textbullet\ {\poplogo OnBuch\color{poporange}+}}\par\vspace{2pt}
      {\raggedright\popsemi\fontsize{8}{10.5}\selectfont\color{popline}\MakeUppercase{Équipe pédagogique OnBuch+}\\\MakeUppercase{Conforme au programme officiel MINESEC}\par}
    \end{minipage}\hfill
    {\poplogo\fontsize{22}{22}\selectfont\color{popcream}OnBuch\color{poporange}+}
  \end{tcolorbox}%
}

% ============================================================
% Page de garde
% #1 titre · #2 matière · #3 niveau · #4 module · #5 n° leçon · #6 durée ·
% #7 description / objectifs (texte ou itemize)
% ============================================================
\newcommand{\pagegarde}[7]{%
  \thispagestyle{empty}
  \newgeometry{a4paper,margin=1.7cm,top=1.6cm,bottom=1.4cm}%
  \begin{tikzpicture}[remember picture,overlay]
    \fill[poporange!18] ([xshift=-3.2cm,yshift=-3.6cm]current page.north east) circle (4.6cm);
    \fill[poppurple!14] ([xshift=2.4cm,yshift=7.6cm]current page.south west) circle (3.8cm);
  \end{tikzpicture}%
  {\poplogo\fontsize{30}{30}\selectfont\color{popink}OnBuch\color{poporange}+}\hfill
  \raisebox{0.6ex}{\poppill{Cours signé \textbullet\ Premium}{popink}{popcream}}\par
  \vspace{1pt}\popsquiggle{poppurple}\par
  \vspace{8pt}
  \begin{tcolorbox}[enhanced,colback=poporange,colframe=popink,boxrule=2.8pt,arc=13pt,
      drop shadow=popink,left=18pt,right=18pt,top=12pt,bottom=13pt,after skip=10pt]
    \poppill{#2 \textbullet\ #3}{popink}{popcream}\hspace{5pt}\poppill{#4}{white}{popink}\par\vspace{8pt}
    {\popdisplay\fontsize{14}{14}\selectfont\color{popink}LEÇON #5}\par\vspace{4pt}
    {\raggedright\hyphenpenalty=10000\exhyphenpenalty=10000\popdisplay\fontsize{25}{27}\selectfont\color{popcream}\MakeUppercase{#1}\par}
    \vspace{8pt}
    {\popxbold\fontsize{9.5}{9.5}\selectfont\color{popink}\MakeUppercase{Durée conseillée : #6}}
  \end{tcolorbox}
  \begin{tcolorbox}[enhanced,colback=white,colframe=popink,boxrule=2.2pt,arc=10pt,
      drop shadow=popink,left=16pt,right=16pt,top=10pt,bottom=10pt,after skip=8pt]
    \poppill{Dans cette leçon}{poppurple}{white}\par\vspace{5pt}
    \popsemi\fontsize{9.3}{12.6}\selectfont\color{popink}#7
  \end{tcolorbox}
  \noindent\begin{minipage}[t]{0.487\linewidth}
    \begin{tcolorbox}[enhanced,colback=popgreen,colframe=popink,boxrule=2.2pt,arc=10pt,
        drop shadow=popink,left=11pt,right=11pt,top=10pt,bottom=10pt,height=3.1cm,valign=center]
      \tcbox[colback=white,colframe=popink,boxrule=1.3pt,arc=5pt,boxsep=0pt,nobeforeafter,
        left=3pt,right=3pt,top=3pt,bottom=3pt,valign=center]{\qrcode[height=1.9cm]{https://play.google.com/store/apps/details?id=com.luvvix.onbuch&hl=fr}}%
      \hspace{8pt}\begin{minipage}[c]{0.5\linewidth}
        {\popdisplay\fontsize{11}{12.5}\selectfont\color{white}\MakeUppercase{Télécharge OnBuch}}\par\vspace{4pt}
        {\raggedright\popsemi\fontsize{8.4}{11}\selectfont\color{white}Gratuit \textbullet\ Google Play\\Tuteur IA, annales, concours, résultats\par}
      \end{minipage}
    \end{tcolorbox}
  \end{minipage}\hfill
  \begin{minipage}[t]{0.487\linewidth}
    \begin{tcolorbox}[enhanced,colback=poppurple,colframe=popink,boxrule=2.2pt,arc=10pt,
        drop shadow=popink,left=11pt,right=11pt,top=10pt,bottom=10pt,height=3.1cm,valign=center]
      \tikz[baseline=-0.7ex]\node[circle,fill=white,draw=popink,line width=1.3pt,minimum size=1.6cm,inner sep=0]{\popdisplay\fontsize{24}{24}\selectfont\color{poppurple}+};%
      \hspace{8pt}\begin{minipage}[c]{0.6\linewidth}
        {\popdisplay\fontsize{11}{12.5}\selectfont\color{white}\MakeUppercase{Passe à OnBuch+}}\par\vspace{4pt}
        {\raggedright\popsemi\fontsize{8.4}{11}\selectfont\color{white}Tous les cours signés, Tuteur IA illimité, fiches PDF — onglet OnBuch+ de l'app\par}
      \end{minipage}
    \end{tcolorbox}
  \end{minipage}
  \vspace{8pt}
  \popcontenu
  \vfill
  \signatureonbuch
  \restoregeometry
  \newpage
}

% Rangée « Ce cours contient » (page de garde)
\newcommand{\poptile}[3]{% #1 couleur #2 gros texte #3 légende
  \begin{tcolorbox}[enhanced,colback=#1,colframe=popink,boxrule=2pt,arc=9pt,shadow={2.5pt}{-2.5pt}{0pt}{popink},
      left=6pt,right=6pt,top=9pt,bottom=9pt,halign=center,valign=center,height=1.75cm,width=\linewidth,nobeforeafter]
    {\popdisplay\fontsize{12.5}{13}\selectfont\color{white}\MakeUppercase{#2}}\par\vspace{3pt}
    {\centering\popsemi\fontsize{7.8}{9.5}\selectfont\color{white}#3\par}
  \end{tcolorbox}}
\newcommand{\popcontenu}{%
  \par\noindent{\popxboldls\fontsize{7.5}{7.5}\selectfont\color{popmuted}CE COURS SIGNÉ CONTIENT}\par\vspace{7pt}
  \noindent\begin{minipage}[t]{0.235\linewidth}\poptile{popblue}{Cours}{complet, illustré, pas à pas}\end{minipage}\hfill
  \begin{minipage}[t]{0.235\linewidth}\poptile{poporange}{Méthodes}{et exercices résolus}\end{minipage}\hfill
  \begin{minipage}[t]{0.235\linewidth}\poptile{poppink}{Exercices}{type examen + corrigés}\end{minipage}\hfill
  \begin{minipage}[t]{0.235\linewidth}\poptile{popgreen}{Fiche bilan}{l'essentiel à retenir}\end{minipage}\par}

% Sommaire
\newtcolorbox{sommaire}{enhanced,colback=white,colframe=popink,boxrule=2.2pt,arc=10pt,
  drop shadow=popink,left=18pt,right=18pt,top=13pt,bottom=13pt,before skip=6pt,after skip=20pt,
  before upper={\poppill{Sommaire}{poppurple}{white}\par\vspace{9pt}\popsemi\fontsize{10.6}{14.5}\selectfont\color{popink}}}

% Signature de fin de cours — poussée tout en bas de la dernière page
\newcommand{\findecours}{%
  \par\vfill\nopagebreak
  \begin{center}\popsquiggle{poporange}\end{center}\vspace{4pt}
  \signatureonbuch
}

%% FIGFIX-BEGIN v4 — correctifs globaux des figures (docs/onbuch-generation/tools/figfix)
\usepackage{adjustbox}\usepackage{etoolbox}
% toute figure TikZ/pgfplots plus large que la ligne est réduite à la largeur disponible
\BeforeBeginEnvironment{tikzpicture}{\begin{adjustbox}{max width=\linewidth}}
\AfterEndEnvironment{tikzpicture}{\end{adjustbox}}
% légendes pgfplots lisibles par-dessus les courbes
\pgfplotsset{every axis/.append style={legend style={fill=white,fill opacity=0.92,text opacity=1,draw=black!15,inner sep=3pt}}}
% étiquettes d'arêtes (syntaxe edge["..."]) avec fond blanc
\tikzset{every edge quotes/.append style={fill=white,inner sep=1.2pt}}
%% FIGFIX-END
\begin{document}

\pagegarde{Champ magnétique, force de Lorentz et force de Laplace}{Physique}{Tle D}{Module 2 : Mouvements et interactions — évolutions temporelles des systèmes mécaniques}{P04}{5 h}{Cette leçon introduit le champ magnétique, ses sources (aimants et courants) et ses effets sur les charges en mouvement (force de Lorentz) et sur les conducteurs parcourus par un courant (force de Laplace). Elle aboutit aux applications technologiques majeures : haut-parleur, moteur électrique et définition de l'ampère.
\vspace{4pt}
\begin{multicols}{2}\small
\begin{itemize}
\item À la fin de cette leçon, je sais décrire les interactions magnétiques entre aimants et définir le vecteur champ magnétique
\item À la fin de cette leçon, je sais représenter le vecteur champ magnétique en un point et tracer les lignes de champ d'un aimant ou d'un courant
\item À la fin de cette leçon, je sais donner les caractéristiques du champ magnétique créé par un fil rectiligne, une spire, une bobine plate et un solénoïde
\item À la fin de cette leçon, je sais exprimer la force de Lorentz F = q v ∧ B et déterminer sa direction, son sens et sa norme
\item À la fin de cette leçon, je sais exprimer la force de Laplace F = I ℓ ∧ B et appliquer la règle de la main droite
\item À la fin de cette leçon, je sais expliquer le fonctionnement d'un haut-parleur et le principe du moteur électrique à courant continu
\end{itemize}\end{multicols}}

\begin{sommaire}
\begin{multicols}{2}
\begin{enumerate}
\item Aimants et champ magnétique
\item Champ magnétique créé par un courant
\item Force de Lorentz sur une charge en mouvement
\item Force de Laplace sur un conducteur
\item Applications : haut-parleur et moteur électrique
\item Interaction entre deux conducteurs parallèles et définition de l'ampère
\item Activité d'intégration
\item Méthodes et exercices résolus
\item Exercices
\item Corrigés des exercices
\item Fiche bilan
\end{enumerate}
\end{multicols}
\end{sommaire}

\begin{objectifs}
\begin{itemize}
\item À la fin de cette leçon, je sais décrire les interactions magnétiques entre aimants et définir le vecteur champ magnétique
\item À la fin de cette leçon, je sais représenter le vecteur champ magnétique en un point et tracer les lignes de champ d'un aimant ou d'un courant
\item À la fin de cette leçon, je sais donner les caractéristiques du champ magnétique créé par un fil rectiligne, une spire, une bobine plate et un solénoïde
\item À la fin de cette leçon, je sais exprimer la force de Lorentz F = q v ∧ B et déterminer sa direction, son sens et sa norme
\item À la fin de cette leçon, je sais exprimer la force de Laplace F = I ℓ ∧ B et appliquer la règle de la main droite
\item À la fin de cette leçon, je sais expliquer le fonctionnement d'un haut-parleur et le principe du moteur électrique à courant continu
\item À la fin de cette leçon, je sais analyser l'interaction entre deux conducteurs parallèles et relier ce phénomène à la définition de l'ampère
\item À la fin de cette leçon, je sais résoudre un problème type Bac combinant champ magnétique et forces magnétiques
\end{itemize}
\end{objectifs}

\begin{prerequis}
\begin{itemize}
\item Produit vectoriel de deux vecteurs : direction, sens, norme (mathématiques de Première/Terminale)
\item Notion de vecteur, représentation et addition vectorielle
\item Intensité du courant électrique, tension, loi d'Ohm (Seconde/Première)
\item Notion de force et caractéristiques d'une force (point d'application, direction, sens, norme)
\item Charge électrique élémentaire e = 1,6 × 10⁻¹⁹ C et structure de l'atome
\end{itemize}
\end{prerequis}

\newpage

\coursec{Aimants et champ magnétique}

À Douala, un technicien répare le haut-parleur d'un bar de quartier : pourquoi la membrane vibre-t-elle quand le courant passe dans la bobine plongée dans l'aimant ? Au lycée, une boussole placée près d'un fil électrique parcouru par un courant dévie brutalement : le courant crée-t-il un champ magnétique ? Dans les centrales hydroélectriques de Song Loulou et d'Édéa, d'immenses forces magnétiques font tourner les alternateurs qui alimentent tout le Cameroun. Comprendre ces phénomènes, c'est comprendre le fonctionnement des moteurs, des haut-parleurs et des générateurs qui nous entourent.

Tous ces dispositifs reposent sur une même notion fondamentale : le \cle{champ magnétique}. Avant de comprendre comment un courant crée un champ magnétique ou comment une charge en mouvement subit une force, nous devons d'abord répondre à des questions plus simples : qu'est-ce qu'un aimant ? Comment détecter et décrire un champ magnétique ? Quelles sont les propriétés de ce champ autour d'un aimant, et qu'en est-il du champ magnétique terrestre ? C'est l'objet de cette première section.

\courssub{Les aimants : pôles Nord et Sud}

\textbf{Un peu d'histoire et de vocabulaire.}\par

Depuis l'Antiquité, on connaît la \cle{magnétite}, un minerai de fer (oxyde de fer $\mathrm{Fe_3O_4}$) qui attire naturellement les objets en fer : c'est l'\emph{aimant naturel}. Son nom vient de la région de Magnésie, en Asie Mineure, où l'on extrayait ce minerai. Aujourd'hui, on fabrique des \emph{aimants artificiels} bien plus puissants : aimants en ferrite (ceux des haut-parleurs), aimants en alliage d'aluminium-nickel-cobalt (alnico), et aimants au néodyme-fer-bore, capables de soulever des masses des centaines de fois supérieures à la leur.

Un aimant possède toujours deux régions où son action magnétique est maximale : ses \cle{pôles}. On les distingue ainsi :

\begin{itemize}
\item le \cle{pôle Nord} (noté N) : c'est l'extrémité qui, lorsque l'aimant est libre de tourner autour d'un axe vertical (aimant suspendu à un fil ou posé sur un flotteur), s'oriente vers le Nord géographique ;
\item le \cle{pôle Sud} (noté S) : c'est l'extrémité qui s'oriente vers le Sud géographique.
\end{itemize}

\begin{propriete}[Inséparabilité des pôles]
Les pôles d'un aimant sont \textbf{inséparables} : si l'on casse un aimant en deux, on n'obtient pas un pôle Nord isolé et un pôle Sud isolé, mais \textbf{deux aimants complets}, chacun doté d'un pôle Nord et d'un pôle Sud. Il n'existe pas de « monopôle magnétique » dans la nature, contrairement aux charges électriques qui, elles, peuvent exister séparément.
\end{propriete}

Cette différence avec l'électrostatique est fondamentale : une charge positive peut exister sans charge négative, mais un pôle Nord magnétique n'existe jamais sans son pôle Sud associé.

\courssub{Interactions magnétiques : attraction et répulsion}

\begin{experience}[Interaction entre deux aimants]
\textbf{Matériel :} deux aimants droits dont les pôles sont repérés (N en rouge, S en bleu par exemple), un fil de suspension.

\textbf{Protocole :} on suspend l'un des aimants horizontalement par un fil fin de sorte qu'il puisse tourner librement. On approche successivement chaque pôle du second aimant de chaque pôle de l'aimant suspendu.

\textbf{Observations :}
\begin{itemize}
\item deux pôles de \emph{même nom} (N--N ou S--S) se \textbf{repoussent} ;
\item deux pôles de \emph{noms contraires} (N--S) s'\textbf{attirent}.
\end{itemize}

\textbf{Interprétation :} les aimants exercent l'un sur l'autre des actions mécaniques à distance, appelées \emph{interactions magnétiques}. Ces interactions vérifient le principe des actions réciproques (troisième loi de Newton) : si l'aimant 1 attire l'aimant 2 avec une force $\vec{F}_{1/2}$, alors l'aimant 2 attire l'aimant 1 avec la force $\vec{F}_{2/1} = -\vec{F}_{1/2}$.
\end{experience}

\begin{aretenir}[Règle des pôles]
Deux pôles de même nom se repoussent ; deux pôles de noms contraires s'attirent. Un aimant attire aussi les objets en fer, en acier, en nickel ou en cobalt (matériaux dits \emph{ferromagnétiques}), quel que soit le pôle présenté.
\end{aretenir}

\courssub{Détection du champ magnétique}

Comment mettre en évidence l'existence d'une action magnétique en un point de l'espace, sans y placer un second aimant ? Deux outils simples suffisent.

\textbf{La boussole (aiguille aimantée).}\par

Une \cle{aiguille aimantée} est une petite aiguille d'acier aimantée, mobile autour d'un axe vertical. Placée en un point $M$ de l'espace, elle tourne puis s'immobilise dans une direction bien déterminée. Si l'on déplace la boussole autour d'un aimant, la direction prise par l'aiguille change d'un point à l'autre : l'aimant crée autour de lui quelque chose qui oriente l'aiguille. Ce « quelque chose » est le champ magnétique.

\textbf{La limaille de fer.}\par

Chaque grain de limaille de fer se comporte, dans un champ magnétique, comme une minuscule aiguille aimantée. En saupoudrant de la limaille sur une plaque de verre posée au-dessus d'un aimant, on voit les grains s'orienter et dessiner des courbes régulières : c'est le \cle{spectre magnétique} de l'aimant. La limaille révèle ainsi la géométrie du champ dans tout le plan.

\begin{definition}[Vecteur champ magnétique]
En tout point $M$ de l'espace, le champ magnétique est représenté par un vecteur $\vec{B}(M)$ appelé \textbf{vecteur champ magnétique} :
\begin{itemize}
\item sa \textbf{direction} et son \textbf{sens} sont ceux de l'aiguille aimantée placée en $M$, orientée \textbf{du pôle Sud vers le pôle Nord} de l'aiguille ;
\item sa \textbf{norme} $B$ se mesure avec un teslamètre (muni d'une sonde de Hall) et s'exprime en \textbf{tesla}, de symbole $\mathrm{T}$, en hommage à Nikola Tesla.
\end{itemize}
\end{definition}

\begin{exemplebox}[Ordres de grandeur du champ magnétique]
\begin{center}
\begin{tabularx}{\linewidth}{|X|c|}
\hline
\rowcolor{popblueL} \textbf{Source du champ} & \textbf{Valeur typique de $B$} \\
\hline
Champ magnétique terrestre (au Cameroun) & $\approx \SI{5e-5}{T}$ \\
\hline
Aimant usuel (aimant droit de laboratoire) & $\SI{1e-2}{T}$ à $\SI{1e-1}{T}$ \\
\hline
Aimant en U de laboratoire (entre les branches) & $\approx \SI{0,05}{T}$ \\
\hline
Électroaimant puissant & quelques teslas (\SI{1}{T} à \SI{10}{T}) \\
\hline
IRM hospitalière & \SI{1,5}{T} à \SI{3}{T} \\
\hline
\end{tabularx}
\end{center}
Le tesla est une unité très « grande » : le champ terrestre, qui suffit pourtant à orienter les boussoles, est cent mille fois plus faible que le champ d'un électroaimant de laboratoire.
\end{exemplebox}

\courssub{Lignes de champ magnétique}

Pour visualiser le champ magnétique dans tout l'espace, on utilise une représentation géométrique : les lignes de champ.

\begin{definition}[Ligne de champ magnétique]
Une \textbf{ligne de champ magnétique} est une courbe de l'espace qui, en chacun de ses points $M$, est \textbf{tangente au vecteur} $\vec{B}(M)$ et orientée dans le sens de $\vec{B}$. L'ensemble des lignes de champ constitue le \textbf{spectre magnétique}.
\end{definition}

\begin{propriete}[Propriétés des lignes de champ (admise)]
\begin{itemize}
\item Les lignes de champ magnétique sont des \textbf{courbes fermées} : elles sortent du pôle Nord de l'aimant, entrent par le pôle Sud, et se referment \emph{à l'intérieur} de l'aimant, du Sud vers le Nord.
\item À l'\textbf{extérieur} de l'aimant, elles sont donc orientées \textbf{du pôle Nord vers le pôle Sud}.
\item Deux lignes de champ \textbf{ne se croisent jamais} : en un point donné, le vecteur $\vec{B}$ a une direction unique.
\item Là où les lignes se resserrent (près des pôles), le champ est plus intense ; là où elles s'écartent, il est plus faible.
\end{itemize}
\end{propriete}

\begin{attention}[Ligne de champ ou trajectoire ?]
Une ligne de champ n'est \textbf{pas} la trajectoire d'une charge ou d'un pôle : c'est une courbe purement géométrique qui indique, en chaque point, la direction et le sens de $\vec{B}$. Ne dis jamais qu'une aiguille « suit une ligne de champ » : elle s'oriente \emph{tangentiellement} à la ligne de champ au point où elle se trouve.
\end{attention}

\begin{popfigure}
\begin{center}
\begin{tikzpicture}[scale=1.05, >=Stealth]
% Aimant droit : Sud (gauche, bleu), Nord (droite, orange)
\fill[popblue!70] (-1.6,-0.45) rectangle (0,0.45);
\fill[poporange!80] (0,-0.45) rectangle (1.6,0.45);
\draw[thick, popdark] (-1.6,-0.45) rectangle (1.6,0.45);
\node[white, font=\bfseries\large] at (-0.8,0) {S};
\node[white, font=\bfseries\large] at (0.8,0) {N};
% Lignes de champ (du N vers le S à l'extérieur) : flèche à la fin du chemin (->)
\foreach \y/\h in {0.9/1.15, 1.5/1.7, 2.1/2.25}{
  \draw[popink, thick, ->]
    (1.6,0.1) .. controls (\h+1.2,\y+0.4) and (-\h-1.2,\y+0.4) .. (-1.6,0.1);
  \draw[popink, thick, ->]
    (1.6,-0.1) .. controls (\h+1.2,-\y-0.4) and (-\h-1.2,-\y-0.4) .. (-1.6,-0.1);
}
% Vecteurs B en quelques points
% À droite de l'aimant (x>0) : B vers la gauche (vers S)
\draw[->, very thick, popgreen!80!black] (3.4,0) -- (2.6,0) node[midway, above]{$\vec{B}$};
% À gauche de l'aimant (x<0) : B vers la droite (vers S qui est à gauche)
\draw[->, very thick, popgreen!80!black] (-3.4,0) -- (-2.6,0) node[midway, above]{$\vec{B}$};
% Au-dessus de l'aimant (y>0) : B vers la gauche (N->S)
\draw[->, very thick, popgreen!80!black] (0.5,2.6) -- (-0.5,2.6) node[midway, right]{$\vec{B}$};
% Boussole en bas à droite : pointe vers le pôle Sud (haut-gauche, angle 135°)
\draw[popdark] (3.0,-1.6) circle (0.28);
\draw[->, thick, poppurple] (3.0,-1.6) -- ++(135:0.25);
\node[popdark, font=\small] at (3.0,-2.1) {boussole};
\end{tikzpicture}
\end{center}
\legende{Spectre magnétique d'un aimant droit : à l'extérieur, les lignes de champ vont du pôle Nord vers le pôle Sud. Le vecteur $\vec{B}$ est tangent aux lignes de champ en chaque point.}
\end{popfigure}

\courssub{Champ magnétique uniforme}

Dans certaines régions de l'espace, le champ magnétique garde les mêmes caractéristiques partout : c'est le cas, par exemple, entre les branches d'un aimant en U.

\begin{definition}[Champ magnétique uniforme]
Un champ magnétique est dit \textbf{uniforme} dans une région de l'espace si le vecteur $\vec{B}$ y est \textbf{constant en direction, en sens et en norme}. Les lignes de champ y sont alors des \textbf{droites parallèles équidistantes}.
\end{definition}

C'est précisément la situation réalisée dans l'entrefer (l'espace) d'un aimant en U : entre les deux branches, les lignes de champ sont quasi rectilignes et parallèles, allant du pôle Nord vers le pôle Sud. Cette zone de champ uniforme est très utilisée au laboratoire : c'est là que l'on place les conducteurs pour étudier la force de Laplace (section 4) ou les bobines des haut-parleurs (section 5).

\begin{popfigure}
\begin{center}
\begin{tikzpicture}[scale=1.0, >=Stealth]
% Aimant en U : Nord en haut (orange), Sud en bas (bleu)
\fill[poporange!80] (-2.6,0.6) rectangle (-1.4,2.4);
\fill[popblue!70] (-2.6,-2.4) rectangle (-1.4,-0.6);
\fill[popmuted!60] (-2.6,-0.6) rectangle (-1.4,0.6);
\draw[thick, popdark] (-2.6,-2.4) rectangle (-1.4,2.4);
\draw[thick, popdark] (-1.4,-0.6) -- (-1.4,0.6);
\node[white, font=\bfseries\large] at (-2,1.5) {N};
\node[white, font=\bfseries\large] at (-2,-1.5) {S};
% Lignes de champ uniformes entre les branches (du N vers le S -> vers le bas)
\foreach \x in {-1.0,-0.2,0.6,1.4}{
  \draw[popink, thick, ->] (\x,0.55) -- (\x,-0.55);
}
% Vecteurs B égaux (vers le bas)
\foreach \x in {-1.0,-0.2,0.6,1.4}{
  \draw[->, very thick, popgreen!80!black] (\x,1.3) -- (\x,0.75);
}
\node[popgreen!60!black] at (2.2,1.05) {$\vec{B}$ uniforme};
\node[popdark, font=\small, align=center] at (0.2,-1.5) {zone de champ\\uniforme (entrefer)};
\end{tikzpicture}
\end{center}
\legende{Aimant en U : entre les branches, les lignes de champ sont parallèles et équidistantes, orientées du pôle Nord vers le pôle Sud ; le champ $\vec{B}$ y est uniforme.}
\end{popfigure}

\begin{exoresolu}[Lecture d'un spectre magnétique]
\textbf{Énoncé.} On place une petite boussole en un point $M$ situé à $\SI{3}{cm}$ au-dessus du milieu d'un aimant droit dont le pôle Nord est à droite. (a) Quelle orientation prend l'aiguille de la boussole ? (b) Un teslamètre indique $B = \SI{2,0e-3}{T}$ en $M$. Comparer cette valeur au champ magnétique terrestre $B_T \approx \SI{5e-5}{T}$ et conclure.
\tcblower
\textbf{Solution.} (a) En $M$, la ligne de champ de l'aimant est tangente à une courbe allant du pôle Nord (à droite) vers le pôle Sud (à gauche) : au sommet de la courbe, la tangente est horizontale et dirigée \emph{vers la gauche}. L'aiguille de la boussole s'oriente donc horizontalement, son pôle Nord pointant vers le pôle Sud de l'aimant (vers la gauche).

(b) Formons le rapport :
\[ \frac{B}{B_T} = \frac{\num{2,0e-3}}{\num{5e-5}} = 40. \]
Le champ de l'aimant en $M$ est \textbf{40 fois plus intense} que le champ terrestre : c'est donc lui qui impose l'orientation de la boussole, l'effet du champ terrestre étant négligeable à cette distance. Si l'on éloignait suffisamment la boussole de l'aimant, $B$ diminuerait et la boussole finirait par s'aligner sur le champ terrestre (direction Nord--Sud géographique).
\end{exoresolu}

\courssub{Le champ magnétique terrestre}

\textbf{La Terre, un aimant géant.}\par

La Terre elle-même crée un champ magnétique, appelé \cle{champ magnétique terrestre}. Tout se passe, en première approximation, comme si la Terre contenait en son centre un immense aimant droit, légèrement incliné par rapport à l'axe de rotation. Ce champ trouve son origine dans les mouvements du fer liquide du noyau externe terrestre (effet « dynamo »).

\begin{attention}[Le piège classique : pôle Nord géographique et pôle Nord magnétique]
L'aiguille de la boussole pointe son pôle \emph{Nord} vers le Nord géographique. Or deux pôles de noms contraires s'attirent ! Cela signifie que le pôle magnétique situé \textbf{près du pôle Nord géographique} est en réalité un \textbf{pôle Sud magnétique}. Dit autrement : l'« aimant terrestre » a son pôle Sud magnétique vers le Nord géographique, et son pôle Nord magnétique vers le Sud géographique. C'est l'une des erreurs les plus fréquentes au Baccalauréat : retiens que \emph{le Nord géographique attire le Nord de la boussole, donc c'est un Sud magnétique}.
\end{attention}

\textbf{Composantes horizontale et verticale.}\par

En un point de la surface terrestre, le vecteur champ magnétique $\vec{B}_T$ n'est généralement pas horizontal : il « plonge » vers le sol dans l'hémisphère Nord. On le décompose en :
\begin{itemize}
\item une \cle{composante horizontale} $\vec{B}_H$, dirigée vers le Nord magnétique : c'est elle qui oriente les boussoles de navigation ;
\item une \cle{composante verticale} $\vec{B}_V$, dirigée vers le bas dans l'hémisphère Nord.
\end{itemize}
L'angle $I$ que fait $\vec{B}_T$ avec le plan horizontal s'appelle l'\cle{inclinaison} magnétique. On a alors les relations :
\[ B_H = B_T \cos I \qquad \text{et} \qquad B_V = B_T \sin I. \]
Aux pôles magnétiques, $I = 90°$ : le champ est vertical et une boussole classique ne sert à rien. À l'équateur magnétique (qui passe non loin de l'Afrique centrale), $I \approx 0$ : le champ est quasi horizontal, et $B_T \approx B_H \approx \SI{5e-5}{T}$. C'est pourquoi, au Cameroun, les boussoles fonctionnent particulièrement bien.

\begin{popfigure}
\begin{center}
\begin{tikzpicture}[scale=1.0, >=Stealth]
% Terre
\fill[popblue!15] (0,0) circle (2);
\draw[thick, popblue] (0,0) circle (2);
% Axe géographique
\draw[dashed, popdark] (0,-2.6) -- (0,2.6);
\node[popdark, font=\small] at (0.15,2.75) {N géo};
\node[popdark, font=\small] at (0.15,-2.75) {S géo};
% Aimant intérieur incliné (Nord magnétique vers le bas à gauche, Sud magnétique vers le haut à droite)
\begin{scope}[rotate=12]
\fill[poporange!80] (-0.35,-1.5) rectangle (0.35,0); % Nord magnétique (blanc N) en bas
\fill[popblue!70] (-0.35,0) rectangle (0.35,1.5);   % Sud magnétique (blanc S) en haut
\draw[thick, popdark] (-0.35,-1.5) rectangle (0.35,1.5);
\node[white, font=\bfseries] at (0,0.75) {S}; % Sud magnétique près du Nord géo
\node[white, font=\bfseries] at (0,-0.75) {N}; % Nord magnétique près du Sud géo
\end{scope}
% Lignes de champ extérieures (du Nord magnétique vers le Sud magnétique)
% Nord magnétique vers angle -78° (ou 282°), Sud magnétique vers 102°.
% On trace des lignes sortantes (Nord magnétique) et entrantes (Sud magnétique) de façon schématique.
\foreach \a in {30, 90, 150}{
  \draw[popink, thick, ->]
    (\a+180:2.9) .. controls (\a+180:4.6) and (\a:4.6) .. (\a:2.9);
}
\node[popink, font=\small] at (4.4,1.2) {lignes de champ};
\node[popdark, font=\small, align=center] at (-4.1,-1.6) {pôle Sud magnétique\\(proche du Nord géographique)};
\draw[->, popdark, thin] (-3.2,-1.3) -- (-0.6,1.9);
\end{tikzpicture}
\end{center}
\legende{La Terre assimilée à un aimant droit : le pôle Sud magnétique de l'« aimant terrestre » est proche du pôle Nord géographique, ce qui explique l'orientation des boussoles.}
\end{popfigure}

\textbf{La boussole et la navigation.}\par

Pendant des siècles, la boussole a été l'instrument de navigation par excellence : en s'alignant sur la composante horizontale $\vec{B}_H$ du champ terrestre, elle indique le Nord magnétique, proche du Nord géographique (l'écart entre les deux, appelé \emph{déclinaison}, dépend du lieu et varie lentement au fil des ans). Marins, explorateurs puis aviateurs s'en sont servis pour s'orienter bien avant l'invention du GPS.

\begin{savaistu}[Des oiseaux navigateurs au-dessus du Cameroun]
Certains oiseaux migrateurs qui traversent le Cameroun lors de leurs voyages entre l'Europe et l'Afrique australe — comme les cigognes ou les guêpiers — possèdent dans leur organisme des récepteurs sensibles au champ magnétique terrestre (on parle de \emph{magnétoréception}). Ils utilisent la direction et l'inclinaison de $\vec{B}_T$ comme une véritable boussole interne pour retrouver leur route sur des milliers de kilomètres. Certains animaux, comme les tortues marines et les homards, partagent cette étonnante faculté !
\end{savaistu}

\begin{exoresolu}[Composantes du champ magnétique terrestre]
\textbf{Énoncé.} En un lieu, le champ magnétique terrestre a pour norme $B_T = \SI{4,7e-5}{T}$ et l'inclinaison vaut $I = 20°$ (le champ plonge vers le sol). Calculer les composantes horizontale $B_H$ et verticale $B_V$. On donne $\cos 20° \approx \num{0,94}$ et $\sin 20° \approx \num{0,34}$.
\tcblower
\textbf{Solution.} On applique directement les relations de projection :
\begin{align*}
B_H &= B_T \cos I = \num{4,7e-5} \times \num{0,94} \approx \SI{4,4e-5}{T}, \\
B_V &= B_T \sin I = \num{4,7e-5} \times \num{0,34} \approx \SI{1,6e-5}{T}.
\end{align*}
Vérification par le théorème de Pythagore :
\[ \sqrt{B_H^2 + B_V^2} = \sqrt{(\num{4,4e-5})^2 + (\num{1,6e-5})^2} \approx \SI{4,7e-5}{T} = B_T. \]
La composante horizontale, seule utile pour une boussole de navigation, représente ici environ $94\,\%$ du champ total : le lieu est proche de l'équateur magnétique.
\end{exoresolu}

\begin{aretenir}[L'essentiel de la section]
\begin{itemize}
\item Un aimant possède deux pôles inséparables, Nord et Sud ; deux pôles de même nom se repoussent, deux pôles de noms contraires s'attirent.
\item En tout point $M$, le vecteur champ magnétique $\vec{B}(M)$ a la direction et le sens de l'aiguille aimantée placée en $M$ (du Sud vers le Nord de l'aiguille) ; sa norme s'exprime en tesla ($\mathrm{T}$).
\item Les lignes de champ sont des courbes fermées, tangentes à $\vec{B}$, orientées du Nord vers le Sud à l'extérieur de l'aimant ; elles ne se croisent jamais.
\item Un champ est uniforme si $\vec{B}$ y est constant : c'est le cas entre les branches d'un aimant en U.
\item La Terre se comporte comme un aimant droit dont le pôle Sud magnétique est proche du Nord géographique ; le champ terrestre ($\approx \SI{5e-5}{T}$) se décompose en composantes horizontale et verticale, reliées par l'inclinaison $I$.
\end{itemize}
\end{aretenir}

Maintenant que nous savons décrire le champ magnétique d'un aimant, une question s'impose naturellement, celle de la boussole qui dévie près d'un fil parcouru par un courant : \emph{un courant électrique crée-t-il, lui aussi, un champ magnétique ?} C'est ce que nous allons découvrir dans la section suivante.

\coursec{Champ magnétique créé par un courant}

Dans la section précédente, nous avons découvert que les aimants créent autour d'eux un champ magnétique $\vec{B}$, matérialisé par des lignes de champ allant du pôle Nord vers le pôle Sud. Une question surgit alors naturellement : les aimants sont-ils les \emph{seules} sources de champ magnétique ? En 1820, le physicien danois Hans Christian \cle{Œrsted} apporte une réponse qui va révolutionner la physique : \textbf{un simple courant électrique crée un champ magnétique}. Cette découverte, fondatrice de l'électromagnétisme, est à l'origine de tous nos moteurs électriques, des relais des centrales hydroélectriques d'Edéa et de Nachtigal aux petits moteurs des ventilateurs de nos salles de classe. Étudions-la en détail.

\courssub{L'expérience d'Œrsted : le courant dévie l'aiguille aimantée}

\begin{experience}[Expérience d'Œrsted (1820)]
\textbf{Matériel :} un fil de cuivre rectiligne, un générateur de courant continu, un interrupteur, une boussole (aiguille aimantée) placée au-dessous du fil.

\textbf{Protocole :}
\begin{enumerate}
\item Disposer le fil horizontalement, au-dessus de la boussole, parallèlement à l'aiguille au repos (direction Nord--Sud).
\item Interrupteur ouvert (pas de courant) : observer l'aiguille.
\item Fermer l'interrupteur et faire circuler un courant $I$ : observer à nouveau.
\item Inverser le sens du courant et observer.
\end{enumerate}

\textbf{Observations :}
\begin{itemize}
\item Sans courant, l'aiguille s'oriente selon le champ magnétique terrestre (Nord--Sud).
\item Avec le courant, l'aiguille \textbf{dévie} et prend une nouvelle direction d'équilibre, quasiment perpendiculaire au fil si le courant est assez intense.
\item En inversant le sens du courant, l'aiguille tourne d'environ $180\degres$ : la déviation change de sens.
\end{itemize}

\textbf{Interprétation :} l'aiguille aimantée ne peut tourner que sous l'action d'un champ magnétique. Le courant qui traverse le fil a donc créé, dans l'espace environnant, un champ magnétique qui se superpose au champ terrestre. Le sens de ce champ dépend du sens du courant.
\end{experience}

\begin{aretenir}[Conclusion fondamentale]
Tout conducteur parcouru par un courant électrique crée dans l'espace qui l'entoure un \cle{champ magnétique} $\vec{B}$. La direction et le sens de $\vec{B}$ dépendent de la géométrie du circuit et du sens du courant ; sa norme dépend de l'intensité $I$ et de la position du point d'observation.
\end{aretenir}

\begin{savaistu}[Une découverte par hasard... ou presque]
Œrsted fit sa découverte lors d'une démonstration publique à Copenhague : en fermant un circuit, il remarqua que l'aiguille d'une boussole voisine bougeait. Quelques semaines plus tard, Ampère reprenait ces résultats et construisait la théorie complète de l'électromagnétisme. Aujourd'hui, sans cette expérience, pas d'électro-aimants des relais de la SONATREL, pas de moteurs, pas de haut-parleurs dans nos téléphones !
\end{savaistu}

\courssub{Règles d'orientation : relier le sens du courant au sens de $\vec{B}$}

Le champ magnétique créé par un courant s'enroule autour du conducteur. Pour trouver son sens, on dispose de règles équivalentes.

\begin{methode}[Trouver le sens de $\vec{B}$ créé par un courant]
\textbf{Règle du tire-bouchon de Maxwell (universelle) :} on place un tire-bouchon le long du conducteur et on le fait \emph{tourner} dans le sens du courant (vu en coupe, le courant tourne autour du fil) ; le tire-bouchon \emph{progresse} dans le sens de $\vec{B}$.

\textbf{Règle de la main droite (pour un fil rectiligne) :}
\begin{enumerate}
\item Pointer le \textbf{pouce} de la main droite dans le sens du courant $I$.
\item Enrouler les quatre autres doigts autour du fil : ils indiquent le sens des lignes de champ, c'est-à-dire le sens de $\vec{B}$.
\end{enumerate}

\textbf{Règle de la main droite (pour une spire ou un solénoïde) :}
\begin{enumerate}
\item Enrouler les doigts de la main droite dans le sens du courant $I$ (dans les spires).
\item Le \textbf{pouce} tendu indique alors le sens de $\vec{B}$ \emph{à l'intérieur} de la spire (du pôle Sud vers le pôle Nord).
\end{enumerate}
\end{methode}

\courssub{Champ créé par un fil rectiligne infini}

\courssub{Description des lignes de champ}

L'expérience de la limaille de fer (ou de petites boussoles disposées autour d'un fil vertical traversant une plaque horizontale) montre que les lignes de champ d'un fil rectiligne sont des \textbf{cercles concentriques centrés sur le fil}, situés dans des plans perpendiculaires au fil. Plus on s'éloigne du fil, plus le champ s'affaiblit : les lignes s'écartent.

\begin{popfigure}
\begin{center}
\begin{tikzpicture}[scale=1.0]
% Cercles = lignes de champ (sens trigonométrique = courant sortant)
\draw[popblue, thick, -{Stealth}] (0,0) circle (0.6);
\draw[popblue, thick, -{Stealth}] (0,0) circle (1.1);
\draw[popblue, thick, -{Stealth}] (0,0) circle (1.6);
\draw[popblue, thick, -{Stealth}] (0,0) circle (2.1);
% Fil en coupe : courant sortant (point)
\fill[poporange] (0,0) circle (0.18);
\draw[popdark, thick] (0,0) circle (0.18);
\fill[popdark] (0,0) circle (0.03); % Point central
% Vecteur B en un point M à l'Est
\fill[popdark] (1.6,0) circle (0.05) node[below right, popdark]{$M$};
\draw[-{Stealth}, poppink, very thick] (1.6,0) -- (1.6,1.0) node[right, poppink]{$\vec{B}(M)$};
% Distance d
\draw[{Stealth}-{Stealth}, popdark] (0,-0.25) -- (1.6,-0.25) node[midway, below]{$d$};
% Légende courant
\node[popdark] at (0,-2.6) {$\odot$ : courant sortant ($I$ vers nous)};
\end{tikzpicture}
\end{center}
\legende{Fil rectiligne vu en coupe (le courant $I$ sort du plan de la figure, symbole $\odot$). Les lignes de champ sont des cercles centrés sur le fil ; $\vec{B}(M)$ est tangent au cercle passant par $M$, de norme $B = \dfrac{\mu_0 I}{2\pi d}$.}
\end{popfigure}

\begin{propriete}[Champ créé par un fil rectiligne infini (admise)]
Un fil rectiligne de grande longueur, parcouru par un courant d'intensité $I$, crée en un point $M$ situé à la distance $d$ du fil un champ magnétique :
\begin{itemize}
\item \textbf{direction :} tangent au cercle de centre le fil, passant par $M$, dans le plan perpendiculaire au fil ;
\item \textbf{sens :} donné par la règle du tire-bouchon ou de la main droite ;
\item \textbf{norme :}
\[ B = \frac{\mu_0\, I}{2\pi d} \]
où $\mu_0 = 4\pi \times 10^{-7}\ \text{T}\cdot\text{m}\cdot\text{A}^{-1}$ est la \cle{perméabilité magnétique du vide}.
\end{itemize}
Cette formule est admise en Terminale ; elle résulte de la loi de Biot et Savart, étudiée dans le supérieur.
\end{propriete}

\textbf{Vérification dimensionnelle.} $[\mu_0] = \text{T}\cdot\text{m}\cdot\text{A}^{-1}$, donc $\left[\dfrac{\mu_0 I}{2\pi d}\right] = \dfrac{\text{T}\cdot\text{m}\cdot\text{A}^{-1}\times\text{A}}{\text{m}} = \text{T}$ : on obtient bien des teslas.

\begin{exemplebox}[Champ créé par une ligne électrique]
Une ligne de transport d'électricité est parcourue par un courant $I = 500\,\text{A}$. Calculons le champ magnétique créé à $d = 10\,\text{m}$ de la ligne :
\[ B = \frac{\mu_0 I}{2\pi d} = \frac{4\pi \times 10^{-7} \times 500}{2\pi \times 10} = \frac{2 \times 10^{-7} \times 500}{10} = 1{,}0 \times 10^{-5}\,\text{T}. \]
Ce champ vaut environ le quart de la composante horizontale du champ magnétique terrestre ($\approx 2 \times 10^{-5}\,\text{T}$ sous nos latitudes) : il reste modeste, mais parfaitement détectable avec une boussole.
\end{exemplebox}

\courssub{Champ créé par une spire circulaire}

Considérons maintenant un fil enroulé en une spire circulaire de rayon $R$, parcourue par un courant $I$. Les lignes de champ se resserrent à l'intérieur de la spire : au centre $O$, le champ est \textbf{perpendiculaire au plan de la spire}.

\begin{propriete}[Champ au centre d'une spire circulaire (admise)]
Une spire circulaire de rayon $R$ parcourue par un courant $I$ crée en son centre un champ magnétique :
\begin{itemize}
\item \textbf{direction :} perpendiculaire au plan de la spire ;
\item \textbf{sens :} donné par la règle de la main droite (doigts dans le sens du courant, pouce = sens de $\vec{B}$ au centre) ou du tire-bouchon ;
\item \textbf{norme :} $B = \dfrac{\mu_0 I}{2R}$.
\end{itemize}
\end{propriete}

\begin{definition}[Faces Nord et Sud d'une spire]
Une spire parcourue par un courant se comporte comme un aimant : elle présente une \cle{face Nord} et une \cle{face Sud}.
\begin{itemize}
\item La face par laquelle les lignes de champ \textbf{sortent} est la face \textbf{Nord} : vue de cette face, le courant circule dans le sens \emph{trigonométrique} (anti-horaire).
\item La face par laquelle les lignes de champ \textbf{entrent} est la face \textbf{Sud} : vue de cette face, le courant circule dans le sens \emph{horaire}.
\end{itemize}
Moyen mnémotechnique : \textbf{N}ord = courant \textbf{N}on horaire (trigonométrique) ; \textbf{S}ud = courant horaire (\textbf{S}ens des aiguilles).
\end{definition}

\begin{popfigure}
\begin{center}
\begin{tikzpicture}[scale=1.0]
% Spire vue de profil : deux sections du fil
% Fil gauche (x=-1.8) : courant sortant (point)
\fill[poporange] (-1.8,0) circle (0.22);
\draw[popdark, thick] (-1.8,0) circle (0.22);
\fill[popdark] (-1.8,0) circle (0.03);
\node[popdark, below] at (-1.8,-0.35) {$I$ sortant ($\odot$)};
% Fil droit (x=1.8) : courant entrant (croix)
\fill[poporange] (1.8,0) circle (0.22);
\draw[popdark, thick] (1.8,0) circle (0.22);
\draw[popdark, thick] (1.73,-0.07) -- (1.87,0.07);
\draw[popdark, thick] (1.73,0.07) -- (1.87,-0.07);
\node[popdark, below] at (1.8,-0.35) {$I$ entrant ($\otimes$)};
% Axe et vecteur B au centre
\draw[popmuted, dashed] (-3.2,0) -- (3.2,0);
\draw[-{Stealth}, poppink, very thick] (0,0) -- (0,1.6) node[right, poppink]{$\vec{B}(O)$};
\fill[popdark] (0,0) circle (0.05) node[below left, popdark]{$O$};
% Lignes de champ intérieures
\draw[popblue, thick, -{Stealth}] (-1.5,0.55) .. controls (0,1.1) .. (1.5,0.55);
\draw[popblue, thick, -{Stealth}] (1.5,-0.55) .. controls (0,-1.1) .. (-1.5,-0.55);
% Faces Nord/Sud (Haut/Bas)
\node[popgreen, font=\bfseries] at (0,1.9) {Face Nord};
\node[popgreen, font=\bfseries] at (0,-1.9) {Face Sud};
\draw[popgreen, thick] (0,1.6) -- (0,1.2);
\draw[popgreen, thick] (0,-1.6) -- (0,-1.2);
% Rayon
\draw[{Stealth}-{Stealth}, popdark] (0,-1.5) -- (1.8,-1.5) node[midway, below]{$R$};
\draw[popmuted, dashed] (0,0) -- (0,-1.5);
\draw[popmuted, dashed] (1.8,0) -- (1.8,-1.5);
\end{tikzpicture}
\end{center}
\legende{Spire circulaire vue en coupe (plan horizontal). Le courant sort à gauche ($\odot$) et entre à droite ($\otimes$) : le champ $\vec{B}(O)$ au centre est vertical, dirigé vers le haut. Les lignes de champ sortent par la face Nord (dessus) et entrent par la face Sud (dessous).}
\end{popfigure}

\courssub{Bobine plate de $N$ spires}

Une \cle{bobine plate} est un enroulement de $N$ spires circulaires de même rayon $R$, jointives, parcourues par le même courant $I$. Les champs créés par chaque spire au centre s'ajoutent (principe de superposition) : ils ont même direction et même sens, donc les normes s'additionnent.

\begin{propriete}[Champ au centre d'une bobine plate (admise)]
Le champ magnétique au centre d'une bobine plate de $N$ spires de rayon $R$ parcourue par un courant $I$ a pour norme :
\[ B = \frac{\mu_0 N I}{2R}. \]
\end{propriete}

\begin{exemplebox}[Bobine plate de TP]
Une bobine plate de $N = 20$ spires de rayon $R = 10\,\text{cm} = 0{,}10\,\text{m}$ est parcourue par un courant $I = 3\,\text{A}$. Champ au centre :
\[ B = \frac{4\pi \times 10^{-7} \times 20 \times 3}{2 \times 0{,}10} = \frac{4\pi \times 10^{-7} \times 60}{0{,}20} \approx 3{,}8 \times 10^{-4}\,\text{T}, \]
soit environ 19 fois le champ terrestre : de quoi faire tourner nettement une aiguille de boussole placée au centre !
\end{exemplebox}

\courssub{Le solénoïde : un champ uniforme à l'intérieur}

\begin{definition}[Solénoïde]
Un \cle{solénoïde} est une bobine longue constituée de $N$ spires jointives enroulées sur un cylindre de longueur $L$ et de rayon $R$, avec $L \gg R$ (longueur grande devant le rayon). On note $n = \dfrac{N}{L}$ le nombre de spires par mètre (en $\text{m}^{-1}$).
\end{definition}

L'exploration du champ avec une petite boussole montre un résultat remarquable : \textbf{à l'intérieur du solénoïde (loin des extrémités), les lignes de champ sont parallèles à l'axe et équidistantes} : le champ y est \cle{uniforme} (même direction, même sens, même norme en tout point). À l'extérieur, le spectre ressemble à celui d'un aimant droit : les lignes sortent par la face Nord et rentrent par la face Sud.

\begin{propriete}[Champ à l'intérieur d'un solénoïde long (admise)]
À l'intérieur d'un solénoïde long parcouru par un courant $I$, le champ magnétique est \textbf{uniforme} :
\begin{itemize}
\item \textbf{direction :} celle de l'axe du solénoïde ;
\item \textbf{sens :} donné par la règle de la main droite (pouce = sens de $\vec{B}$ à l'intérieur, doigts = sens du courant) ;
\item \textbf{norme :}
\[ B = \mu_0\, n\, I = \mu_0\, \frac{N}{L}\, I, \]
avec $n$ en spires par mètre ($\text{m}^{-1}$), $I$ en ampères, $B$ en teslas.
\end{itemize}
\end{propriete}

\textbf{Vérification dimensionnelle.} $[\mu_0 n I] = \text{T}\cdot\text{m}\cdot\text{A}^{-1} \times \text{m}^{-1} \times \text{A} = \text{T}$. Correct.

\begin{popfigure}
\begin{center}
\begin{tikzpicture}[scale=0.9]
% Solénoïde stylisé : spires vues en coupe (axe horizontal)
% Courant : en haut (y>0) -> entrant (croix) ; en bas (y<0) -> sortant (point)
% Pour que B pointe vers la droite (Nord à droite), pouce à droite -> doigts : haut=rentrant, bas=sortant.
\foreach \x in {-2.4,-1.8,-1.2,-0.6,0,0.6,1.2,1.8,2.4}{
  % Brin supérieur : courant entrant (croix)
  \draw[poporange, thick] (\x,0.55) circle (0.14);
  \draw[popdark, thick] (\x-0.05,0.62) -- (\x+0.05,0.48);
  \draw[popdark, thick] (\x-0.05,0.48) -- (\x+0.05,0.62);
  % Brin inférieur : courant sortant (point)
  \fill[poporange] (\x,-0.55) circle (0.14);
  \draw[popdark, thick] (\x,-0.55) circle (0.14);
  \fill[popdark] (\x,-0.55) circle (0.028);
}
% Axe
\draw[popmuted, dashed] (-3.6,0) -- (3.6,0);
% Lignes de champ intérieures parallèles (vers la droite)
\foreach \y in {-0.3,0,0.3}{
  \draw[popblue, thick, -{Stealth}] (-2.9,\y) -- (2.9,\y);
}
% Lignes extérieures (retour)
\draw[popblue, thick, -{Stealth}] (3.3,0.9) .. controls (0,1.6) .. (-3.3,0.9);
\draw[popblue, thick, -{Stealth}] (-3.3,-0.9) .. controls (0,-1.6) .. (3.3,-0.9);
% Vecteur B intérieur
\draw[-{Stealth}, poppink, very thick] (-0.8,0) -- (0.8,0) node[midway, above, poppink]{$\vec{B}$};
% Faces
\node[popgreen, font=\bfseries] at (-3.1,1.5) {Face Sud};
\node[popgreen, font=\bfseries] at (3.1,1.5) {Face Nord};
% Dimensions
\draw[{Stealth}-{Stealth}, popdark] (-2.4,-1.15) -- (2.4,-1.15) node[midway, below]{$L$};
\node[popdark] at (0,-1.7) {$n = N/L$ spires par mètre};
\end{tikzpicture}
\end{center}
\legende{Solénoïde en coupe : à l'intérieur, les lignes de champ sont parallèles à l'axe (champ uniforme $B = \mu_0 n I$) ; à l'extérieur, le spectre est analogue à celui d'un aimant droit, les lignes allant de la face Nord (droite) vers la face Sud (gauche). Le courant entre par le haut ($\otimes$) et sort par le bas ($\odot$).}
\end{popfigure}

\begin{exoresolu}[Champ dans un solénoïde de laboratoire]
Un solénoïde comporte $N = 500$ spires régulièrement réparties sur une longueur $L = 40\,\text{cm}$. Il est parcouru par un courant d'intensité $I = 2\,\text{A}$.
\begin{enumerate}
\item Calculer le nombre de spires par mètre $n$.
\item En déduire la norme du champ magnétique à l'intérieur du solénoïde.
\item Comparer au champ magnétique terrestre ($B_T \approx 5 \times 10^{-5}\,\text{T}$).
\end{enumerate}
\tcblower
\textbf{Solution.}
\begin{enumerate}
\item Conversion : $L = 40\,\text{cm} = 0{,}40\,\text{m}$.
\[ n = \frac{N}{L} = \frac{500}{0{,}40} = 1\,250\,\text{m}^{-1}. \]
\item Champ intérieur :
\begin{align*}
B &= \mu_0\, n\, I = 4\pi \times 10^{-7} \times 1\,250 \times 2 \\
&= 4\pi \times 10^{-7} \times 2\,500 \\
&= 10\,000\pi \times 10^{-7} = \pi \times 10^{-3} \\
&\approx 3{,}14 \times 10^{-3}\,\text{T}.
\end{align*}
\item Rapport : $\dfrac{B}{B_T} = \dfrac{3{,}14 \times 10^{-3}}{5 \times 10^{-5}} \approx 63$. Le champ du solénoïde est environ \textbf{63 fois plus intense} que le champ terrestre : c'est lui qui imposera l'orientation d'une aiguille aimantée placée à l'intérieur.
\end{enumerate}
\end{exoresolu}

\begin{attention}[Erreurs fréquentes sur le solénoïde]
\begin{itemize}
\item \textbf{Oublier les conversions :} $L$ doit être exprimé en \emph{mètres} avant de calculer $n = N/L$. Avec $L = 40\,\text{cm}$ non converti, on trouve $n = 12{,}5$ au lieu de $1\,250\,\text{m}^{-1}$ : un facteur 100 d'erreur !
\item \textbf{Confondre $N$ et $n$ :} $N$ est le nombre \emph{total} de spires (sans unité), $n = N/L$ est le nombre de spires \emph{par mètre} (en $\text{m}^{-1}$). La formule $B = \mu_0 n I$ utilise $n$, jamais $N$ seul.
\item \textbf{Appliquer la formule hors de son domaine :} $B = \mu_0 n I$ n'est valable qu'\emph{à l'intérieur} d'un solénoïde \emph{long} ($L \gg R$), loin des extrémités. Au centre d'une spire unique, c'est $B = \dfrac{\mu_0 I}{2R}$ qu'il faut utiliser.
\end{itemize}
\end{attention}

\courssub{Superposition des champs magnétiques}

Lorsqu'un point $M$ de l'espace est soumis à plusieurs sources de champ magnétique (deux fils, un fil et le champ terrestre, etc.), chaque source crée en $M$ son propre vecteur champ, \emph{comme si elle était seule}.

\begin{propriete}[Principe de superposition]
Le champ magnétique total en un point $M$ est la \textbf{somme vectorielle} des champs créés par chaque source :
\[ \vec{B}(M) = \vec{B}_1(M) + \vec{B}_2(M) + \cdots \]
La norme du champ total s'obtient par la géométrie du triangle des vecteurs (théorème de Pythagore si les champs sont perpendiculaires, relation d'Al-Kashi dans le cas général).
\end{propriete}

\begin{exoresolu}[Aiguille aimantée au centre d'une bobine]
Une petite aiguille aimantée est placée au centre d'une bobine plate dont l'axe est horizontal et perpendiculaire au méridien magnétique. Sans courant, l'aiguille s'oriente selon la composante horizontale du champ terrestre $B_T = 2{,}0 \times 10^{-5}\,\text{T}$. On fait passer un courant tel que la bobine crée au centre un champ $B_b = 4{,}0 \times 10^{-5}\,\text{T}$, perpendiculaire à $\vec{B}_T$.
\begin{enumerate}
\item Représenter les vecteurs $\vec{B}_T$, $\vec{B}_b$ et le champ total $\vec{B}$.
\item Calculer la norme de $\vec{B}$ et l'angle $\alpha$ dont dévie l'aiguille.
\end{enumerate}
\tcblower
\textbf{Solution.}
\begin{enumerate}
\item $\vec{B}_T$ pointe vers le Nord magnétique ; $\vec{B}_b$ lui est perpendiculaire (axe de la bobine). Le champ total est la diagonale du rectangle construit sur ces deux vecteurs : l'aiguille s'oriente selon $\vec{B} = \vec{B}_T + \vec{B}_b$.
\item Les deux champs étant perpendiculaires, le théorème de Pythagore donne :
\[ B = \sqrt{B_T^2 + B_b^2} = \sqrt{(2{,}0 \times 10^{-5})^2 + (4{,}0 \times 10^{-5})^2} = \sqrt{20 \times 10^{-10}} \approx 4{,}5 \times 10^{-5}\,\text{T}. \]
L'angle de déviation vérifie :
\[ \tan\alpha = \frac{B_b}{B_T} = \frac{4{,}0 \times 10^{-5}}{2{,}0 \times 10^{-5}} = 2{,}0 \quad\Longrightarrow\quad \alpha \approx 63\degres. \]
L'aiguille tourne donc d'environ $63\degres$ par rapport au Nord magnétique. C'est le principe de la \emph{boussole des tangentes}, instrument historique de mesure des courants.
\end{enumerate}
\end{exoresolu}

\begin{methode}[Résoudre un problème de superposition de champs]
\begin{enumerate}
\item Identifier chaque source et \textbf{représenter} en $M$ chaque vecteur champ (direction et sens par la règle du tire-bouchon ou de la main droite, norme par la formule adaptée : fil, spire, solénoïde).
\item Choisir un repère et \textbf{décomposer} les vecteurs si nécessaire.
\item Additionner les composantes, puis calculer la norme et la direction du champ total.
\item Vérifier la cohérence des ordres de grandeur (le champ terrestre vaut $\approx 10^{-5}\,\text{T}$ : une source qui prétend l'emporter doit créer un champ au moins comparable).
\end{enumerate}
\end{methode}

\begin{aretenir}[L'essentiel de la section]
\begin{itemize}
\item Tout courant crée un champ magnétique (expérience d'Œrsted) ; le sens de $\vec{B}$ s'obtient par la règle du tire-bouchon ou de la main droite.
\item Fil rectiligne infini : lignes circulaires, $B = \dfrac{\mu_0 I}{2\pi d}$.
\item Spire de rayon $R$ : $B = \dfrac{\mu_0 I}{2R}$ au centre, perpendiculaire au plan ; faces Nord (courant trigonométrique) et Sud (courant horaire).
\item Bobine plate de $N$ spires : $B = \dfrac{\mu_0 N I}{2R}$.
\item Solénoïde long : champ \textbf{uniforme} à l'intérieur, $B = \mu_0 n I$ avec $n = N/L$ ; spectre extérieur d'aimant droit.
\item Les champs magnétiques s'ajoutent \textbf{vectoriellement} (superposition).
\item $\mu_0 = 4\pi \times 10^{-7}\,\text{T}\cdot\text{m}\cdot\text{A}^{-1}$ ; toutes les distances en mètres !
\end{itemize}
\end{aretenir}

Maintenant que nous savons \emph{créer} et \emph{calculer} un champ magnétique, une question s'impose : comment ce champ agit-il sur une charge électrique en mouvement ? C'est l'objet de la section suivante, consacrée à la force de Lorentz. Nous verrons ensuite (section 4) la force de Laplace sur un conducteur, puis (section 6) l'interaction entre deux conducteurs parallèles qui permettra de définir l'ampère.

\coursec{Force de Lorentz sur une charge en mouvement}

Dans la section précédente, nous avons vu qu'un courant électrique crée autour de lui un champ magnétique $\vec{B}$. Mais un courant, qu'est-ce que c'est, au fond ? Un déplacement ordonné de charges électriques. Si un champ magnétique est capable d'agir sur un conducteur parcouru par un courant (nous le verrons avec la force de Laplace), il est naturel de penser qu'il agit directement sur \emph{chaque charge en mouvement}. C'est exactement ce qu'a établi le physicien néerlandais Hendrik Lorentz à la fin du XIX\textsuperscript{e} siècle. Cette section est le cœur de la leçon : elle explique pourquoi les anciens téléviseurs à tube fonctionnaient, comment on trie les ions en médecine nucléaire, et pourquoi les particules du vent solaire sont canalisées vers les pôles par le champ magnétique terrestre.

\courssub{Mise en évidence expérimentale : déviation d'un faisceau d'électrons}

\begin{experience}[Déviation d'un faisceau d'électrons par un aimant]
\textbf{Matériel :} un tube à rayons cathodiques (tube de Crookes ou tube de déviation), une alimentation haute tension, un aimant en U.

\textbf{Protocole :}
\begin{enumerate}
\item On produit dans le tube un faisceau d'électrons (rayon cathodique) se déplaçant horizontalement à grande vitesse. Le faisceau est rendu visible par un écran fluorescent ou un gaz raréfié : on observe une trace rectiligne lumineuse.
\item On approche un aimant en U de manière que le champ magnétique $\vec{B}$ soit perpendiculaire à la trajectoire des électrons.
\item On inverse les pôles de l'aimant, puis on écarte l'aimant.
\end{enumerate}

\textbf{Observations :}
\begin{itemize}
\item En l'absence d'aimant, le faisceau est \cle{rectiligne}.
\item Avec l'aimant, le faisceau est \cle{dévié} et sa trajectoire s'incurve en arc de cercle.
\item Le sens de la déviation dépend du sens de $\vec{B}$ : en inversant les pôles, la déviation s'inverse.
\item La déviation est d'autant plus marquée que l'aimant est puissant (champ intense) et que les électrons sont lents.
\end{itemize}

\textbf{Interprétation :} une charge électrique en mouvement dans un champ magnétique subit une \cle{force} qui modifie la direction de sa trajectoire. Cette force est la \emph{force de Lorentz magnétique}. Remarque cruciale : si le faisceau est immobile (pas d'émission), aucune déviation n'est possible ; et si $\vec{B}$ est parallèle au faisceau, aucune déviation n'est observée. La force n'existe que si la charge \emph{se déplace} et si sa vitesse n'est pas colinéaire au champ.
\end{experience}

\begin{savaistu}[Le tube cathodique, ancêtre de nos écrans]
Les téléviseurs et oscilloscopes dits « à tube cathodique », encore présents dans certains laboratoires de lycées camerounais, utilisaient précisément ce principe : un faisceau d'électrons, dévié par des champs magnétiques (ou électriques), balayait l'écran pour dessiner l'image. Aujourd'hui, la force de Lorentz reste au cœur des spectromètres de masse (analyse médicale, datation au carbone 14) et des accélérateurs de particules.
\end{savaistu}

\courssub{Expression de la force de Lorentz magnétique}

\begin{definition}[Force de Lorentz magnétique]
Une particule de charge $q$, animée d'une vitesse $\vec{v}$ dans une région où règne un champ magnétique $\vec{B}$, subit la \cle{force de Lorentz magnétique} :
\[ \vec{F} = q\, \vec{v} \wedge \vec{B} \]
Sa norme vaut :
\[ F = |q|\, v\, B\, \sin\alpha \]
où $\alpha$ est l'angle entre les vecteurs $\vec{v}$ et $\vec{B}$. Unités : $F$ en newtons (N), $q$ en coulombs (C), $v$ en mètres par seconde ($\mathrm{m\,s^{-1}}$), $B$ en teslas (T).
\end{definition}

\textbf{Vérification dimensionnelle.} Le produit vectoriel conserve les dimensions, donc $[F] = [q][v][B]$. Or $1\ \mathrm{T} = 1\ \mathrm{N\,A^{-1}\,m^{-1}}$ et $1\ \mathrm{C} = 1\ \mathrm{A\,s}$, d'où :
\[ [q][v][B] = (\mathrm{A\,s})\times(\mathrm{m\,s^{-1}})\times(\mathrm{N\,A^{-1}\,m^{-1}}) = \mathrm{N}. \]
La formule est homogène : un bon réflexe à garder au Baccalauréat !

\textbf{Caractéristiques du vecteur force :}
\begin{itemize}
\item \textbf{Point d'application :} la particule chargée elle-même.
\item \textbf{Direction :} perpendiculaire au plan formé par $\vec{v}$ et $\vec{B}$ (donc perpendiculaire à la fois à $\vec{v}$ et à $\vec{B}$).
\item \textbf{Sens :} celui du produit vectoriel $\vec{v} \wedge \vec{B}$ si $q > 0$ ; \emph{sens opposé} si $q < 0$.
\item \textbf{Norme :} $F = |q|\,v\,B\,\sin\alpha$.
\end{itemize}

\begin{methode}[Déterminer le sens de la force de Lorentz]
\begin{enumerate}
\item Représenter les vecteurs $\vec{v}$ et $\vec{B}$ et identifier l'angle $\alpha$.
\item Appliquer la règle de la main droite au produit $\vec{v} \wedge \vec{B}$ : pouce selon $\vec{v}$, index selon $\vec{B}$, le majeur (perpendiculaire à la paume) indique le sens de $\vec{v} \wedge \vec{B}$. (Variante : la règle du tire-bouchon — on tourne de $\vec{v}$ vers $\vec{B}$, la vis avance dans le sens de $\vec{v} \wedge \vec{B}$.)
\item Si la charge est \textbf{positive} ($q>0$) : $\vec{F}$ a le sens trouvé à l'étape 2.
\item Si la charge est \textbf{négative} ($q<0$, cas de l'électron !) : on \textbf{inverse} le sens trouvé à l'étape 2.
\item Calculer la norme : $F = |q|\,v\,B\,\sin\alpha$.
\end{enumerate}
\end{methode}

\begin{attention}[Le piège classique de l'électron]
L'erreur la plus fréquente au Bac : oublier que l'électron porte une charge \textbf{négative} $q = -e = \num{-1,6e-19}\ \mathrm{C}$. Pour un électron, la force de Lorentz est de sens \textbf{opposé} au produit vectoriel $\vec{v} \wedge \vec{B}$. Rédige toujours $\vec{F} = q\,\vec{v}\wedge\vec{B}$ avec $q$ \emph{algébrique}, et conclus sur le sens en tenant compte du signe. Autre piège : la condition d'existence — si la charge est au repos ($v = 0$) ou si $\vec{v}$ est parallèle à $\vec{B}$ ($\sin\alpha = 0$), alors $\vec{F} = \vec{0}$ : le champ magnétique n'agit pas.
\end{attention}

\textbf{Cas particuliers à connaître par cœur :}
\begin{center}
\begin{tabularx}{\linewidth}{|l|X|X|}
\hline
\rowcolor{popblueL} \textbf{Configuration} & \textbf{Conséquence} & \textbf{Norme de $\vec{F}$} \\
\hline
$\vec{v} \perp \vec{B}$ ($\alpha = 90°$) & Force maximale, trajectoire incurvée & $F = |q|\,v\,B$ \\
\hline
$\vec{v} \parallel \vec{B}$ ($\alpha = 0$ ou $180°$) & Aucune force, mouvement rectiligne uniforme & $F = 0$ \\
\hline
Charge au repos ($v = 0$) & Aucune force & $F = 0$ \\
\hline
\end{tabularx}
\end{center}

\begin{popfigure}
\begin{center}
\begin{tikzpicture}[scale=1.2, >=Stealth]
% ---- Charge positive ----
\begin{scope}
% Axes guides
\draw[->, thin, popmuted] (-0.5,0) -- (2.5,0) node[below] {$x$};
\draw[->, thin, popmuted] (0,-0.5) -- (0,2.5) node[left] {$y$};
\draw[->, thin, popmuted] (0,0) -- (-0.5,-0.5) node[below left] {$z$};
% Vecteurs v et B dans le plan (x,y)
\draw[->, very thick, popblue] (0,0) -- (2,0) node[below right] {$\vec{v}$};
\draw[->, very thick, popgreen!80!black] (0,0) -- (0,2) node[above left] {$\vec{B}$};
% Force : sortie de page (point)
\draw[poporange, very thick] (0,0) circle (0.15);
\fill[poporange] (0,0) circle (0.05);
\node[poporange, above right] at (0.2,0.2) {$\vec{F}\ (\odot)$};
% Particule
\fill[popink] (0,0) circle (0.08);
\node[popink, below left] at (-0.2,-0.2) {$q>0$};
% Angle
\draw[popmuted] (0.5,0) arc (0:90:0.5) node[midway, right, font=\small] {$90^\circ$};
\node[popdark, font=\small\bfseries] at (1, -1.2) {Charge positive};
\end{scope}
% ---- Charge négative ----
\begin{scope}[xshift=8.5cm]
\draw[->, thin, popmuted] (-0.5,0) -- (2.5,0) node[below] {$x$};
\draw[->, thin, popmuted] (0,-0.5) -- (0,2.5) node[left] {$y$};
\draw[->, thin, popmuted] (0,0) -- (-0.5,-0.5) node[below left] {$z$};
\draw[->, very thick, popblue] (0,0) -- (2,0) node[below right] {$\vec{v}$};
\draw[->, very thick, popgreen!80!black] (0,0) -- (0,2) node[above left] {$\vec{B}$};
% Force : entrée de page (croix)
\draw[poppink, very thick] (-0.15,-0.15) -- (0.15,0.15);
\draw[poppink, very thick] (-0.15,0.15) -- (0.15,-0.15);
\node[poppink, above right] at (0.2,0.2) {$\vec{F}\ (\otimes)$};
\fill[popink] (0,0) circle (0.08);
\node[popink, below left] at (-0.2,-0.2) {$q<0$};
\draw[popmuted] (0.5,0) arc (0:90:0.5) node[midway, right, font=\small] {$90^\circ$};
\node[popdark, font=\small\bfseries] at (1, -1.2) {Charge n\'egative (\'electron)};
\end{scope}
\end{tikzpicture}
\end{center}
\legende{Force de Lorentz : $\vec{v}$ selon $x$, $\vec{B}$ selon $y$. Pour $q>0$, $\vec{F}$ sort du plan ($\odot$) ; pour $q<0$ (électron), $\vec{F}$ entre dans le plan ($\otimes$).}
\end{popfigure}

\courssub{Une force qui ne travaille pas : conséquence énergétique}

Voici l'une des propriétés les plus remarquables — et les plus interrogées au Bac — de la force de Lorentz magnétique.

\begin{propriete}[La force magnétique ne travaille pas]
La puissance de la force de Lorentz magnétique est nulle à tout instant :
\[ \mathcal{P} = \vec{F}\cdot\vec{v} = q\,(\vec{v}\wedge\vec{B})\cdot\vec{v} = 0. \]
Par conséquent, la force magnétique \cle{ne modifie pas l'énergie cinétique} de la particule : elle ne change ni la norme de la vitesse ni la température du système. Elle ne fait que \emph{courber la trajectoire}.

\textbf{Démonstration guidée (exigible).} Par définition du produit vectoriel, $\vec{v}\wedge\vec{B}$ est perpendiculaire à $\vec{v}$. Or le produit scalaire de deux vecteurs perpendiculaires est nul, donc $(\vec{v}\wedge\vec{B})\cdot\vec{v} = 0$, d'où $\mathcal{P} = 0$. Le travail sur un déplacement élémentaire vaut $\delta W = \vec{F}\cdot d\vec{r} = \vec{F}\cdot\vec{v}\,dt = 0$. D'après le théorème de l'énergie cinétique, $\Delta E_c = W = 0$ : la norme de $\vec{v}$ reste constante. \emph{CQFD.}
\end{propriete}

\begin{attention}[Perpendiculaire ne veut pas dire « sans effet »]
Beaucoup d'élèves concluent trop vite : « la force ne travaille pas, donc elle ne fait rien ». Faux ! Une force perpendiculaire à la vitesse ne modifie pas la \emph{norme} de $\vec{v}$, mais elle modifie en permanence sa \emph{direction} : c'est exactement le rôle de la force centripète dans le mouvement circulaire uniforme. Sans la force de Lorentz, pas d'arc de cercle dans le tube cathodique.
\end{attention}

\courssub{Complément au programme : mouvement circulaire uniforme dans un champ uniforme}

Ce résultat n'est pas exigible en tant que tel, mais il est régulièrement demandé comme application guidée au Baccalauréat (série C). Il couronne magnifiquement ce qui précède.

\begin{propriete}[Rayon de la trajectoire circulaire — complément Bac]
Une particule de charge $q$, de masse $m$, pénétrant avec une vitesse $\vec{v}$ \textbf{perpendiculaire} à un champ magnétique uniforme $\vec{B}$, décrit un \cle{cercle} de rayon :
\[ R = \frac{m\,v}{|q|\,B} \]
parcouru à vitesse constante (mouvement circulaire uniforme).

\textbf{Démonstration guidée (à savoir refaire).}
\begin{enumerate}
\item \textbf{Bilan des forces :} dans le référentiel du laboratoire galiléen, la particule n'est soumise qu'à la force de Lorentz $\vec{F} = q\,\vec{v}\wedge\vec{B}$ (le poids est négligeable devant $\vec{F}$ pour une particule élémentaire).
\item \textbf{Deuxième loi de Newton :} $\vec{F} = m\,\vec{a}$. Comme $\vec{F}\perp\vec{v}$, l'accélération est purement normale : $a_t = \dfrac{dv}{dt} = 0$ (vitesse constante) et $a_n = \dfrac{v^2}{R}$.
\item \textbf{Projection sur la normale} (base de Frenet) : comme $\vec{v}\perp\vec{B}$, $F = |q|\,v\,B$, d'où :
\[ |q|\,v\,B = m\,\frac{v^2}{R} \quad\Longrightarrow\quad R = \frac{m\,v}{|q|\,B}. \]
\end{enumerate}
La période de révolution $T = \dfrac{2\pi R}{v} = \dfrac{2\pi m}{|q|B}$ est remarquable : elle \emph{ne dépend pas} de la vitesse ! C'est le principe du cyclotron.
\end{propriete}

\begin{popfigure}
\begin{center}
\begin{tikzpicture}[scale=1.1, >=Stealth]
% Champ B entrant : croix
\foreach \x in {-2.5,-1.5,-0.5,0.5,1.5,2.5}{
  \foreach \y in {-2.5,-1.5,-0.5,0.5,1.5,2.5}{
    \draw[popmuted, thin] (\x-0.1,\y-0.1) -- (\x+0.1,\y+0.1);
    \draw[popmuted, thin] (\x-0.1,\y+0.1) -- (\x+0.1,\y-0.1);
  }
}
\node[popmuted, font=\small] at (2.5,3.0) {$\vec{B}$ entrant ($\otimes$)};
% Trajectoire circulaire (sens horaire pour électron)
\draw[very thick, popblue] (0,0) circle (2.0);
% Electron position 1 : (2, 0) - Point le plus à droite
\fill[popink] (2.0,0) circle (0.1) node[right=2pt, popink] {$e^-$};
\draw[->, thick, popgreen!80!black] (2.0,0) -- (2.0,-1.5) node[right] {$\vec{v}$};
\draw[->, very thick, poporange] (2.0,0) -- (0.5,0) node[midway, above] {$\vec{F}$};
% Electron position 2 : (-2, 0) - Point le plus à gauche
\fill[popink] (-2.0,0) circle (0.1) node[left=2pt, popink] {$e^-$};
\draw[->, thick, popgreen!80!black] (-2.0,0) -- (-2.0,1.5) node[left] {$\vec{v}$};
\draw[->, very thick, poporange] (-2.0,0) -- (-0.5,0) node[midway, below] {$\vec{F}$};
% Centre et Rayon
\draw[dashed, popdark] (0,0) -- (2.0,0) node[midway, below, font=\small] {$R=\dfrac{mv}{eB}$};
\fill[popdark] (0,0) circle (0.05);
% Sens de rotation
\draw[->, thick, popblue] (1.5, 1.2) arc (90:45:1.9);
\node[popblue, font=\small] at (1.8, 1.8) {Sens horaire};
\end{tikzpicture}
\end{center}
\legende{Trajectoire circulaire d'un électron (rotation horaire) dans un champ $\vec{B}$ uniforme entrant ($\otimes$). La force de Lorentz, centripète, courbe la trajectoire sans modifier la norme de la vitesse.}
\end{popfigure}

\courssub{Exemple chiffré et exercice résolu}

\begin{exoresolu}[Électron dévié dans le champ magnétique terrestre]
Un électron ($q = -e = \num{-1,6e-19}\ \mathrm{C}$ ; $m = \num{9,11e-31}\ \mathrm{kg}$) se déplace horizontalement à la vitesse $v = \num{2,0e7}\ \mathrm{m\,s^{-1}}$ dans une région où le champ magnétique terrestre vaut $B = \num{5,0e-5}\ \mathrm{T}$, dirigé perpendiculairement à $\vec{v}$.

\textbf{1.} Calculer la norme de la force de Lorentz subie par l'électron.

\textbf{2.} Comparer cette force au poids de l'électron. Conclure.

\textbf{3.} Déterminer le rayon de la trajectoire.

\tcblower

\textbf{1.} Comme $\vec{v}\perp\vec{B}$, $\sin\alpha = 1$ :
\[ F = |q|\,v\,B = \num{1,6e-19}\times\num{2,0e7}\times\num{5,0e-5} = \SI{1,6e-16}{\newton}. \]

\textbf{2.} Le poids vaut $P = m\,g = \num{9,11e-31}\times 9,8 \approx \SI{8,9e-30}{\newton}$. Le rapport est :
\[ \frac{F}{P} = \frac{\num{1,6e-16}}{\num{8,9e-30}} \approx \num{1,8e13}. \]
La force magnétique est environ $10^{13}$ fois plus grande que le poids : \textbf{le poids d'une particule élémentaire est toujours négligeable} devant la force de Lorentz. C'est une justification à rédiger clairement au Bac.

\textbf{3.} La trajectoire est circulaire de rayon :
\[ R = \frac{m\,v}{|q|\,B} = \frac{\num{9,11e-31}\times\num{2,0e7}}{\num{1,6e-19}\times\num{5,0e-5}} = \frac{\num{1,822e-23}}{\num{8,0e-24}} \approx \num{2,3}\ \mathrm{m}. \]
L'électron décrit un cercle de rayon $R \approx \SI{2,3}{m}$, à vitesse constante puisque la force ne travaille pas.
\end{exoresolu}

\begin{exercice}[Pour t'entraîner]
Un proton ($q = +e = \num{1,6e-19}\ \mathrm{C}$ ; $m_p = \num{1,67e-27}\ \mathrm{kg}$) pénètre avec une vitesse $v = \num{1,0e6}\ \mathrm{m\,s^{-1}}$ dans un champ uniforme $B = \num{0,20}\ \mathrm{T}$, perpendiculaire à $\vec{v}$.

\textbf{1.} Calculer la norme de la force de Lorentz.

\textbf{2.} Montrer que la trajectoire est circulaire et calculer son rayon $R$.

\textbf{3.} Calculer la période $T$ de révolution. Dépend-elle de $v$ ?
\end{exercice}

\begin{corrige}[Exercice : proton dans un champ uniforme]
\textbf{1.} $\vec{v}\perp\vec{B}$ donc $F = e\,v\,B = \num{1,6e-19}\times\num{1,0e6}\times\num{0,20} = \SI{3,2e-14}{\newton}$.

\textbf{2.} Le proton n'est soumis qu'à $\vec{F} = q\,\vec{v}\wedge\vec{B}$ (poids négligeable : $P = \num{1,6e-26}\ \mathrm{N} \ll F$). La deuxième loi de Newton projetée sur la normale donne $e\,v\,B = m_p\dfrac{v^2}{R}$, d'où :
\[ R = \frac{m_p\,v}{e\,B} = \frac{\num{1,67e-27}\times\num{1,0e6}}{\num{1,6e-19}\times\num{0,20}} \approx \num{5,2e-2}\ \mathrm{m} = \SI{5,2}{cm}. \]

\textbf{3.} $T = \dfrac{2\pi R}{v} = \dfrac{2\pi m_p}{eB} = \dfrac{2\pi\times\num{1,67e-27}}{\num{1,6e-19}\times\num{0,20}} \approx \SI{3,3e-7}{\second}$. La vitesse $v$ n'apparaît plus : la période est \textbf{indépendante de $v$}.
\end{corrige}

\begin{aretenir}[L'essentiel de la section]
\begin{itemize}
\item Force de Lorentz magnétique : $\vec{F} = q\,\vec{v}\wedge\vec{B}$, de norme $F = |q|\,v\,B\,\sin\alpha$.
\item Elle n'existe que si la charge \textbf{se déplace} et si $\vec{v}$ n'est \textbf{pas colinéaire} à $\vec{B}$.
\item Direction : perpendiculaire au plan $(\vec{v},\vec{B})$ ; sens donné par la main droite, \textbf{inversé si $q<0$} (électron !).
\item Son travail est \textbf{nul} : elle ne modifie pas l'énergie cinétique, seulement la direction du mouvement.
\item Complément Bac : si $\vec{v}\perp\vec{B}$ uniforme, la trajectoire est un cercle de rayon $R = \dfrac{mv}{|q|B}$, parcouru à vitesse constante.
\end{itemize}
\end{aretenir}

Dans la section suivante, nous verrons comment la somme des forces de Lorentz subies par les porteurs de charge d'un conducteur donne naissance à une force macroscopique sur le conducteur tout entier : la \emph{force de Laplace}, principe de fonctionnement des moteurs électriques.

\coursec{Force de Laplace sur un conducteur}

Dans la section précédente, nous avons vu qu'une particule chargée en mouvement dans un champ magnétique $\vec{B}$ subit la force de Lorentz $\vec{F} = q\,\vec{v} \wedge \vec{B}$. Or, qu'est-ce qu'un courant électrique, sinon un déplacement ordonné de porteurs de charge ? Un conducteur parcouru par un courant et placé dans un champ magnétique contient donc des milliards de milliards de charges en mouvement, chacune soumise à une force de Lorentz. La somme de toutes ces forces microscopiques se transmet au conducteur tout entier : c'est la \cle{force de Laplace}, nommée en l'honneur du mathématicien et physicien français Pierre-Simon Laplace. C'est elle qui fait tourner les moteurs, vibrer les haut-parleurs et fonctionner les pompes des forages. Cette section est le cœur de l'électromécanique.

\courssub{Mise en évidence expérimentale : le rail de Laplace}

\begin{experience}[L'expérience des rails de Laplace]
\textbf{Matériel :} deux rails de cuivre parallèles et horizontaux, une barre de cuivre mobile (tige cylindrique) posée perpendiculairement sur les rails, un aimant en U (champ $\vec{B}$ sensiblement uniforme et vertical entre ses branches), un générateur de courant continu réglable, un interrupteur, des fils de connexion.

\textbf{Protocole :}
\begin{enumerate}
\item On place la barre mobile entre les branches de l'aimant en U, au repos, circuit ouvert : rien ne se passe.
\item On ferme l'interrupteur : un courant $I$ circule dans le circuit, et donc dans la barre.
\item On observe le mouvement de la barre.
\item On inverse le sens du courant, puis on retourne l'aimant (on inverse $\vec{B}$), et on note les effets.
\end{enumerate}

\textbf{Observations :}
\begin{itemize}
\item À la fermeture du circuit, la barre \textbf{se met en mouvement} et roule sur les rails : elle subit donc une force.
\item Si l'on inverse le sens du courant $I$, la barre part dans le \textbf{sens opposé}.
\item Si l'on retourne l'aimant (inversion de $\vec{B}$), la barre part également dans le \textbf{sens opposé}.
\item Si l'on augmente $I$ ou si l'on utilise un aimant plus puissant, le mouvement est plus rapide : la force est plus grande.
\item Si la barre est placée parallèlement à $\vec{B}$, elle ne bouge pas.
\end{itemize}

\textbf{Interprétation :} un conducteur parcouru par un courant et placé dans un champ magnétique subit une force mécanique, appelée \cle{force de Laplace}. Son sens dépend du sens de $I$ et du sens de $\vec{B}$ ; sa norme croît avec $I$, avec $B$ et avec la longueur du conducteur baignant dans le champ.
\end{experience}

\begin{popfigure}
\begin{tikzpicture}[scale=1.05, >=Stealth]
% Aimant en U (vue stylisée)
\fill[popblue!20] (-0.9,-1.5) rectangle (0.9,1.5);
\node[popblue] at (0,-1.8) {\small Aimant en U (champ $\vec{B}$ vertical)};
% Rails
\draw[line width=2.2pt, popdark] (-3.2,-0.55) -- (3.2,-0.55);
\draw[line width=2.2pt, popdark] (-3.2,0.55) -- (3.2,0.55);
\node[popdark] at (-3.55,0.55) {\small rail};
\node[popdark] at (-3.55,-0.55) {\small rail};
% Barre mobile
\draw[line width=3.5pt, poporange] (1.2,-0.75) -- (1.2,0.75);
\node[poporange, right] at (1.28,0.95) {\small barre mobile};
% Courant dans la barre
\draw[->, line width=1.4pt, popgreen] (1.2,-0.6) -- (1.2,0.5);
\node[popgreen, left] at (1.05,0) {$I$};
% Champ B (sortant du plan de la figure) : représenté par des points
\foreach \x in {-0.5,0.1,2.0,2.6} {
  \draw[popblue] (\x,0) circle (0.09);
  \fill[popblue] (\x,0) circle (0.025);
}
\node[popblue] at (2.3,-0.35) {$\vec{B}$};
% Force de Laplace
\draw[->, line width=1.8pt, poppink] (1.2,0) -- (2.9,0);
\node[poppink, above] at (2.6,0.05) {$\vec{F}$};
% Générateur
\draw[popdark] (-3.2,-0.55) -- (-3.9,-0.55) -- (-3.9,-1.3);
\draw[popdark] (-3.2,0.55) -- (-3.9,0.55) -- (-3.9,-0.2);
\draw[popdark, line width=1pt] (-4.25,-0.2) -- (-3.55,-0.2);
\draw[popdark, line width=2.4pt] (-4.1,-1.3) -- (-3.7,-1.3);
\draw[popdark] (-3.9,-0.2) -- (-3.9,-0.75); \draw[popdark, dashed] (-3.9,-0.75) -- (-3.9,-1.3);
\node[popdark, left] at (-4.25,-0.75) {\small G};
\end{tikzpicture}
\legende{Dispositif des rails de Laplace : la barre parcourue par le courant $I$, plongée dans le champ $\vec{B}$ (ici sortant du plan de la figure, symbolisé par des points), subit la force $\vec{F}$ qui la fait rouler sur les rails.}
\end{popfigure}

\courssub{Expression vectorielle de la force de Laplace}

\begin{definition}[Force de Laplace]
Un conducteur rectiligne de longueur $\ell$, parcouru par un courant d'intensité $I$ et placé dans un champ magnétique uniforme $\vec{B}$, subit une force appelée \cle{force de Laplace} :
\[ \vec{F} = I\,\vec{\ell} \wedge \vec{B} \]
où $\vec{\ell}$ est le vecteur dont la norme est égale à la longueur du conducteur \textbf{située dans le champ magnétique} et dont le sens est celui du courant $I$ (sens conventionnel, c'est-à-dire le sens opposé au mouvement des électrons).
\end{definition}

\begin{propriete}[Caractéristiques de la force de Laplace]
La force de Laplace $\vec{F} = I\,\vec{\ell} \wedge \vec{B}$ possède les caractéristiques suivantes :
\begin{itemize}
\item \textbf{Point d'application :} le milieu de la portion de conducteur placée dans le champ magnétique ;
\item \textbf{Direction :} perpendiculaire au plan formé par $\vec{\ell}$ et $\vec{B}$ (donc perpendiculaire à la fois au conducteur et au champ) ;
\item \textbf{Sens :} donné par la règle de la main droite (voir méthode ci-dessous) ; il est tel que le trièdre $(\vec{\ell}, \vec{B}, \vec{F})$ soit direct ;
\item \textbf{Norme :} \[ F = I\,\ell\,B\,\sin\alpha \] où $\alpha$ est l'angle entre la direction du courant (vecteur $\vec{\ell}$) et le vecteur $\vec{B}$.
\end{itemize}
\textbf{Unités SI :} $I$ en ampères (\SI{}{A}), $\ell$ en mètres (\SI{}{m}), $B$ en teslas (\SI{}{T}), $F$ en newtons (\SI{}{N}).
\end{propriete}

Vérifions l'homogénéité de la relation $F = I\ell B$ : le tesla se définit justement par $\SI{1}{T} = \SI{1}{N.A^{-1}.m^{-1}}$, donc $[I\ell B] = \SI{}{A} \times \SI{}{m} \times \SI{}{N.A^{-1}.m^{-1}} = \SI{}{N}$. La formule est bien homogène à une force.

\textbf{Cas particuliers à connaître par cœur :}
\begin{itemize}
\item Conducteur \textbf{perpendiculaire} au champ ($\alpha = 90^\circ$, $\sin\alpha = 1$) : la force est maximale, $F = I\ell B$.
\item Conducteur \textbf{parallèle} au champ ($\alpha = 0$ ou $180^\circ$, $\sin\alpha = 0$) : la force est nulle, $F = 0$. C'est pourquoi la barre mal orientée dans l'expérience ne bouge pas.
\end{itemize}

\begin{methode}[La règle de la main droite pour trouver le sens de $\vec{F}$]
Pour déterminer le sens de la force de Laplace $\vec{F} = I\,\vec{\ell} \wedge \vec{B}$ :
\begin{enumerate}
\item Ouvre la main droite, pouce écarté, index et majeur perpendiculaires entre eux (comme trois arêtes d'un coin de cube).
\item Pointe l'\textbf{index} dans le sens du courant $I$ (c'est-à-dire le sens de $\vec{\ell}$).
\item Pointe le \textbf{majeur} dans le sens du champ magnétique $\vec{B}$.
\item Le \textbf{pouce} indique alors le sens de la force $\vec{F}$.
\end{enumerate}
Variante équivalente (règle du tire-bouchon) : fais tourner le vecteur $\vec{\ell}$ vers $\vec{B}$ par le plus court chemin ; le sens d'avancement du tire-bouchon donne le sens de $\vec{F}$.

\textbf{Vérification systématique :} $\vec{F}$ doit être perpendiculaire à $\vec{\ell}$ \emph{et} à $\vec{B}$. Si ce n'est pas le cas sur ton schéma, il y a une erreur.
\end{methode}

\begin{popfigure}
\begin{tikzpicture}[scale=1.1, >=Stealth]
% Axes représentant la main droite
\draw[->, line width=2pt, popgreen] (0,0) -- (3,0) node[right] {$I$ (index)};
\draw[->, line width=2pt, popblue] (0,0) -- (0,2.6) node[above] {$\vec{B}$ (majeur)};
\draw[->, line width=2.2pt, poppink] (0,0) -- (-1.6,-1.6) node[below left] {$\vec{F}$ (pouce)};
% Angle droit
\draw[popmuted] (0.35,0) -- (0.35,0.35) -- (0,0.35);
\node[popmuted] at (0.75,0.28) {\small $90^\circ$};
% Main stylisée
\begin{scope}[shift={(4.6,1.1)}, scale=0.85]
\fill[popgoldL, draw=popgold, line width=1pt] (0,0) ellipse (0.55 and 0.45);
% pouce
\draw[line width=5pt, popgold, rounded corners=3pt] (-0.35,-0.3) -- (-1.05,-0.95);
% index
\draw[line width=5pt, popgold, rounded corners=3pt] (0.5,0.05) -- (1.45,0.05);
% majeur
\draw[line width=5pt, popgold, rounded corners=3pt] (0.05,0.4) -- (0.05,1.35);
\draw[->, line width=1.3pt, popgreen] (0.5,0.05) -- (1.6,0.05) node[right, popgreen] {\small $I$};
\draw[->, line width=1.3pt, popblue] (0.05,0.4) -- (0.05,1.5) node[above, popblue] {\small $\vec{B}$};
\draw[->, line width=1.3pt, poppink] (-0.35,-0.3) -- (-1.2,-1.05) node[below, poppink] {\small $\vec{F}$};
\end{scope}
\end{tikzpicture}
\legende{Règle de la main droite : index dans le sens du courant $I$, majeur dans le sens de $\vec{B}$, le pouce donne le sens de la force de Laplace $\vec{F}$. Les trois vecteurs forment un trièdre direct.}
\end{popfigure}

\courssub{Du microscopique au macroscopique : la force de Laplace comme somme de forces de Lorentz}

D'où vient cette force ? Nous avons dit qu'un courant est un déplacement de porteurs de charge. Établissons rigoureusement le lien entre la force de Lorentz (microscopique) et la force de Laplace (macroscopique) : c'est une démonstration classique, très formatrice, qui montre l'unité de l'électromagnétisme.

\begin{propriete}[La force de Laplace est la résultante des forces de Lorentz]
Considérons un conducteur cylindrique de section $S$, de longueur $\ell$ placée dans un champ magnétique uniforme $\vec{B}$, contenant $n$ porteurs de charge par unité de volume (pour un métal, ce sont les électrons de conduction, de charge $q = -e$, animés de la vitesse d'ensemble $\vec{v}$).

\textbf{Démonstration guidée :}
\begin{enumerate}
\item \textbf{Nombre de porteurs.} Le volume du conducteur baignant dans le champ est $V = S\ell$. Le nombre total de porteurs de charge qu'il contient est donc :
\[ N = n\,S\,\ell \]
\item \textbf{Force sur un porteur.} Chaque électron, de charge $-e$ et de vitesse $\vec{v}$, subit la force de Lorentz :
\[ \vec{f} = -e\,\vec{v} \wedge \vec{B} \]
\item \textbf{Force totale.} Tous les électrons ont la même vitesse d'ensemble ; la force totale subie par l'ensemble des porteurs, transmise au réseau cristallin du conducteur (c'est-à-dire à la matière elle-même), est :
\[ \vec{F} = N\,\vec{f} = nS\ell\,(-e)\,\vec{v} \wedge \vec{B} = -nSe\,\ell\,\vec{v} \wedge \vec{B} \]
\item \textbf{Introduction de l'intensité.} L'intensité du courant s'écrit $I = nSev$ en valeur absolue. Vectoriellement, en introduisant le vecteur $\vec{\ell}$ orienté dans le \emph{sens conventionnel du courant} (opposé à la vitesse des électrons), on a la relation :
\[ -nSe\,v\,\vec{\ell}_{\text{électrons}} = I\,\vec{\ell} \]
car le signe moins de la charge des électrons compense exactement l'inversion de sens entre $\vec{v}$ et le sens conventionnel du courant. Autrement dit : $-nSe\,\ell\,\vec{v} = I\,\vec{\ell}$.
\item \textbf{Conclusion.} En reportant dans l'expression de $\vec{F}$ :
\[ \vec{F} = I\,\vec{\ell} \wedge \vec{B} \]
\end{enumerate}
La force de Laplace n'est donc pas une nouvelle force de la nature : c'est la \textbf{somme macroscopique des forces de Lorentz} subies par les porteurs de charge du conducteur.
\end{propriete}

\begin{attention}[Bien choisir la longueur $\ell$]
L'erreur la plus fréquente au Baccalauréat est d'utiliser la \textbf{longueur totale du circuit} dans la formule $F = I\ell B\sin\alpha$. C'est faux : seule la portion de conducteur \textbf{effectivement plongée dans le champ magnétique} subit la force. Dans le dispositif des rails de Laplace, si le champ de l'aimant en U n'existe que sur une largeur de \SI{5}{cm}, alors $\ell = \SI{5}{cm}$, même si la barre mesure \SI{20}{cm} et le circuit complet plusieurs mètres. Autre piège : l'angle $\alpha$ est l'angle entre le \emph{courant} et $\vec{B}$, pas entre le conducteur et un rail ou un bord de l'aimant. Enfin, n'oublie pas que $\vec{\ell}$ est orienté dans le \textbf{sens conventionnel} du courant (de $+$ vers $-$ à l'extérieur du générateur), et non dans le sens de déplacement des électrons.
\end{attention}

\courssub{Exemple chiffré et bilan énergétique}

\begin{exemplebox}[Force sur une barre de cuivre]
Une barre de cuivre dont la longueur plongée dans le champ est $\ell = \SI{10}{cm}$ est parcourue par un courant $I = \SI{5}{A}$. Elle est placée perpendiculairement aux lignes d'un champ magnétique uniforme $B = \SI{0,2}{T}$ (valeur typique entre les branches d'un aimant en U de laboratoire).

La barre étant perpendiculaire au champ, $\alpha = 90^\circ$ et $\sin\alpha = 1$ :
\[ F = I\,\ell\,B = 5 \times \num{0,10} \times \num{0,2} = \SI{0,10}{N} \]

Cette force peut sembler modeste, mais elle est du même ordre de grandeur que le poids d'une masse de \SI{10}{g} : largement suffisante pour mettre en mouvement une barre légère bien guidée, comme le confirme l'expérience des rails.
\end{exemplebox}

\textbf{Un point physique capital : la force de Laplace peut travailler.} Nous avons vu dans la section précédente que la force de Lorentz sur une particule chargée \emph{isolée} ne travaille jamais, car elle est à chaque instant perpendiculaire à la vitesse de la particule. La situation est différente pour le conducteur : la force de Laplace s'exerce sur la \emph{matière} du conducteur (le réseau cristallin), dont la vitesse macroscopique n'a aucune raison d'être perpendiculaire à $\vec{F}$. Si la barre se déplace dans le sens de la force sur une distance $d$, le travail est :
\[ W = F\,d = I\ell B\,d > 0 \]
Ce travail moteur provient de l'énergie fournie par le générateur : le rail de Laplace réalise une \textbf{conversion d'énergie électrique en énergie mécanique}. C'est le principe fondamental de tous les moteurs électriques, que nous détaillerons dans la section suivante. Notons au passage que cette apparente contradiction avec la force de Lorentz se lève en regardant de près : le générateur compense en permanence l'effet de la composante de la force magnétique qui s'oppose au mouvement des porteurs le long du fil (c'est le phénomène d'induction, étudié plus tard dans l'année).

\begin{exoresolu}[Rails de Laplace inclinés]
Une barre de cuivre de masse $m = \SI{20}{g}$ peut rouler sans frottement sur deux rails parallèles distants de $\ell = \SI{8,0}{cm}$, inclinés d'un angle $\beta = 15^\circ$ par rapport à l'horizontale. L'ensemble est plongé dans un champ magnétique $\vec{B}$ vertical, de norme $B = \SI{0,50}{T}$, la portion de barre dans le champ étant exactement $\ell$. On fait passer un courant $I$ dans la barre, perpendiculaire à $\vec{B}$.

\textbf{1.} Faire le bilan des forces s'exerçant sur la barre.

\textbf{2.} Calculer l'intensité $I$ nécessaire pour que la barre reste en équilibre sur les rails. On prend $g = \SI{9,8}{N.kg^{-1}}$.
\tcblower
\textbf{1.} La barre est soumise à trois forces :
\begin{itemize}
\item son poids $\vec{P} = m\vec{g}$, vertical, vers le bas, appliqué au centre de la barre ;
\item la réaction $\vec{R}$ des rails, perpendiculaire aux rails (frottements négligés) ;
\item la force de Laplace $\vec{F} = I\vec{\ell}\wedge\vec{B}$ : le courant est horizontal et perpendiculaire à $\vec{B}$ (vertical), donc $\vec{F}$ est horizontale, de norme $F = I\ell B$, dirigée vers le haut de la pente (sens du courant choisi en conséquence, vérifié par la règle de la main droite).
\end{itemize}

\textbf{2.} À l'équilibre, $\vec{P} + \vec{R} + \vec{F} = \vec{0}$. Projetons cette relation sur l'axe parallèle aux rails, orienté vers le haut de la pente. La réaction $\vec{R}$, perpendiculaire aux rails, n'a pas de projection sur cet axe. La composante du poids le long de la pente vaut $-mg\sin\beta$ ; la composante de la force de Laplace (horizontale) le long de la pente vaut $+F\cos\beta$. D'où :
\[ F\cos\beta - mg\sin\beta = 0 \quad\Longrightarrow\quad I\ell B\cos\beta = mg\sin\beta \]
\[ I = \dfrac{mg\tan\beta}{\ell B} = \dfrac{\num{0,020} \times \num{9,8} \times \tan 15^\circ}{\num{0,080} \times \num{0,50}} \]
\[ I = \dfrac{\num{0,020} \times \num{9,8} \times \num{0,268}}{\num{0,040}} \approx \SI{1,3}{A} \]

Il faut donc un courant d'environ $\SI{1,3}{A}$ pour maintenir la barre en équilibre. Vérification d'homogénéité : $\dfrac{\SI{}{kg}\times\SI{}{N.kg^{-1}}}{\SI{}{m}\times\SI{}{T}} = \dfrac{\SI{}{N}}{\SI{}{m.T}} = \SI{}{A}$, correct.
\end{exoresolu}

\begin{aretenir}[L'essentiel sur la force de Laplace]
\begin{itemize}
\item Un conducteur parcouru par un courant $I$ et placé dans un champ magnétique $\vec{B}$ subit la force de Laplace : $\vec{F} = I\,\vec{\ell}\wedge\vec{B}$, de norme $F = I\ell B\sin\alpha$.
\item $\vec{\ell}$ est orienté dans le sens conventionnel du courant et sa norme est la longueur du conducteur \textbf{dans le champ} uniquement.
\item Le sens de $\vec{F}$ s'obtient par la règle de la main droite ; $\vec{F}$ est perpendiculaire au plan $(\vec{\ell}, \vec{B})$.
\item Cas limites : $F = I\ell B$ si le conducteur est perpendiculaire à $\vec{B}$ ; $F = 0$ s'il est parallèle.
\item La force de Laplace est la résultante des forces de Lorentz subies par les porteurs de charge : $\vec{F} = N(-e)\,\vec{v}\wedge\vec{B}$ conduit à $\vec{F} = I\vec{\ell}\wedge\vec{B}$ grâce à $I = nSev$.
\item Contrairement à la force de Lorentz sur une charge isolée, la force de Laplace peut \textbf{travailler} : elle convertit l'énergie électrique en énergie mécanique (principe du moteur électrique).
\end{itemize}
\end{aretenir}

\begin{exercice}[À toi de jouer]
Dans un haut-parleur, une spire de fil de longueur totale $\ell = \SI{3,0}{m}$ est placée dans un champ magnétique radial de norme $B = \SI{0,80}{T}$, le fil étant en tout point perpendiculaire au champ. La bobine est parcourue par un courant $I = \SI{2,0}{A}$.
\begin{enumerate}
\item Calculer la norme de la force de Laplace totale subie par la bobine.
\item Que devient cette force si l'on inverse le sens du courant ? Quel mouvement en résulte-t-il pour la membrane du haut-parleur ?
\item La bobine a une masse $m = \SI{12}{g}$. Calculer l'accélération qu'elle subirait si la force précédente était la seule à agir.
\end{enumerate}
\end{exercice}

\begin{corrige}[Exercice : haut-parleur]
\textbf{1.} Le fil est partout perpendiculaire au champ ($\sin\alpha = 1$) et toute la longueur $\ell$ baigne dans le champ :
\[ F = I\ell B = \num{2,0} \times \num{3,0} \times \num{0,80} = \SI{4,8}{N} \]

\textbf{2.} Si l'on inverse le sens du courant, $\vec{\ell}$ change de sens, donc $\vec{F} = I\vec{\ell}\wedge\vec{B}$ change de sens : la force garde la même norme $\SI{4,8}{N}$ mais s'exerce en sens opposé. Avec un courant alternatif (le signal musical), la force alterne donc de sens à la fréquence du signal : la bobine et la membrane solidaire vibrent, ce qui produit le son.

\textbf{3.} D'après la deuxième loi de Newton appliquée à la bobine (si $\vec{F}$ est la seule force) :
\[ a = \dfrac{F}{m} = \dfrac{\num{4,8}}{\num{0,012}} = \SI{400}{m.s^{-2}} \]
soit environ $40$ fois l'accélération de la pesanteur : on comprend que la membrane puisse suivre des vibrations rapides et reproduire des sons aigus.
\end{corrige}

\begin{savaistu}[Des rails de Laplace aux trains à lévitation magnétique]
L'expérience des rails de Laplace, réalisée pour la première fois au début du XIX\textsuperscript{e} siècle, est l'ancêtre de tous les actionneurs électromagnétiques. Les trains Maglev (comme celui de Shanghai, qui atteint \SI{430}{km.h^{-1}}) utilisent des forces magnétiques géantes pour léviter et se propulser sans contact avec les rails. Au Cameroun, les forces de Laplace sont à l'œuvre chaque jour dans les moteurs des pompes à eau des forages ruraux, dans les ventilateurs, les compresseurs des chambres froides et les démarreurs des véhicules : chaque fois qu'un courant électrique produit un mouvement, c'est Laplace (et Lorentz derrière lui) qui travaille.
\end{savaistu}

\coursec{Applications : haut-parleur et moteur électrique}

Dans la section précédente, nous avons établi la loi de Laplace : une portion de conducteur de longueur $\ell$, parcourue par un courant $I$ et placée dans un champ magnétique $\vec{B}$, subit la force $\vec{F} = I\,\vec{\ell} \wedge \vec{B}$. Cette force, qui peut atteindre plusieurs newtons, n'est pas une simple curiosité de laboratoire : elle est le \cle{moteur invisible} de presque tous les appareils qui transforment l'électricité en mouvement. Du haut-parleur qui diffuse la musique makossa dans un bar de Douala au moteur de la broyeuse de manioc du village, en passant par le ventilateur du salon, le même principe physique est à l'œuvre. Cette section te montre comment, à partir de la seule loi de Laplace, on comprend le fonctionnement du haut-parleur électrodynamique et du moteur à courant continu.

\courssub{Le haut-parleur électrodynamique}

\textbf{Constitution.}\par
Un haut-parleur électrodynamique comporte trois éléments essentiels :
\begin{itemize}
\item un \cle{aimant annulaire} (en forme d'anneau) qui crée, dans l'entrefer circulaire où il est placé, un champ magnétique $\vec{B}$ \emph{radial} : en tout point de l'entrefer, $\vec{B}$ est dirigé perpendiculairement à l'axe du haut-parleur, du centre vers l'extérieur (ou l'inverse selon la polarité de l'aimant) ;
\item une \cle{bobine mobile} de $N$ spires, solidaire d'une membrane, placée dans l'entrefer ; elle peut coulisser le long de l'axe ;
\item une \cle{membrane} (cône en carton ou en matière composite), reliée au châssis par des suspensions élastiques qui jouent le rôle de ressorts et ramènent la bobine à sa position d'équilibre.
\end{itemize}

\begin{popfigure}
\begin{center}
\begin{tikzpicture}[scale=1.0,>=Stealth]
% Aimant annulaire (coupe)
\fill[popblueL] (-3.2,-1.6) rectangle (-1.6,1.6);
\fill[popblueL] (1.6,-1.6) rectangle (3.2,1.6);
\fill[white] (-1.6,-0.55) rectangle (1.6,0.55);
\draw[popdark,thick] (-3.2,-1.6) rectangle (-1.6,1.6);
\draw[popdark,thick] (1.6,-1.6) rectangle (3.2,1.6);
\node[popdark] at (-2.4,1.9) {\small Aimant (N)};
\node[popdark] at (2.4,1.9) {\small Aimant (S)};
% Champ B radial dans l'entrefer
\draw[->,popblue,very thick] (-1.55,0.9) -- (-0.75,0.9);
\node[popblue,above] at (-1.15,0.9) {$\vec{B}$};
\draw[->,popblue,very thick] (1.55,0.9) -- (0.75,0.9);
\node[popblue,above] at (1.15,0.9) {$\vec{B}$};
% Bobine
\draw[poporange,very thick] (-1.1,-0.55) rectangle (1.1,0.55);
\foreach \x in {-0.9,-0.45,0,0.45,0.9}{
  \node[poporange] at (\x,0.28) {\small $\otimes$};
  \node[poporange] at (\x,-0.28) {\small $\odot$};}
\node[poporange,below] at (0,-0.7) {\small Bobine mobile ($N$ spires, courant $I$)};
% Membrane
\draw[popgreen,very thick] (1.1,0.55) -- (2.6,1.5);
\draw[popgreen,very thick] (1.1,-0.55) -- (2.6,-1.5);
\draw[popgreen,very thick] (2.6,-1.5) -- (2.6,1.5);
\node[popgreen,right] at (2.6,0) {\small Membrane};
% Force de Laplace
\draw[->,poppink,ultra thick] (0,0.75) -- (1.6,0.75);
\node[poppink,above] at (1.35,0.75) {$\vec{F}$};
% Axe
\draw[dashed,popmuted] (-3.4,0) -- (3.6,0);
\end{tikzpicture}
\end{center}
\legende{Coupe schématique d'un haut-parleur électrodynamique. Le champ $\vec{B}$ est radial dans l'entrefer ; la bobine parcourue par le courant $I$ subit la force de Laplace $\vec{F}$ dirigée selon l'axe, qui entraîne la membrane.}
\end{popfigure}

\textbf{Analyse physique : pourquoi la membrane vibre.}\par
Appliquons la démarche systématique de la loi de Laplace. Chaque spire de la bobine est un cercle de rayon $r$ placé dans le champ radial $\vec{B}$. En tout point d'une spire, orientons $\vec{\ell}$ dans le sens du courant : $\vec{\ell}$ est tangent à la spire, donc perpendiculaire à $\vec{B}$. Le produit vectoriel $\vec{\ell} \wedge \vec{B}$ est alors dirigé selon l'axe du haut-parleur, et ceci \emph{en tout point de la spire} : toutes les forces élémentaires s'ajoutent. Pour une spire de circonférence $2\pi r$, la force a pour norme $F_1 = I \times 2\pi r \times B$ (car $\sin\theta = 1$). Pour $N$ spires :
\[ F = N\,I\,(2\pi r)\,B. \]
Le courant délivré par l'amplificateur est \emph{variable} : il reproduit le signal musical, par exemple $i(t) = I_{\max}\sin(2\pi f t)$ pour une note pure de fréquence $f$. La force $F(t) = N(2\pi r)B\,i(t)$ varie donc en phase avec le courant.

\textbf{Le va-et-vient de la bobine.}\par
C'est ici que le sens du courant devient décisif. Lorsque $i > 0$, la règle des trois doigts de la main droite (ou le produit vectoriel) donne $\vec{F}$ vers l'avant : la membrane est poussée vers l'extérieur. Lorsque $i < 0$, le courant change de sens, donc $\vec{\ell}$ change de sens, donc $\vec{F}$ s'inverse : la membrane est tirée vers l'arrière. La bobine effectue ainsi un \cle{mouvement de va-et-vient} à la fréquence $f$ du signal. La membrane, de grande surface, met en mouvement l'air environnant : elle crée des zones successives de compression et de dilatation de l'air, c'est-à-dire une \emph{onde sonore}. Le haut-parleur réalise la conversion : énergie électrique $\rightarrow$ énergie mécanique $\rightarrow$ énergie sonore.

\begin{methode}[Analyser un dispositif à force de Laplace]
Pour étudier un haut-parleur, un moteur ou tout dispositif électromagnétique, suis toujours la même démarche :
\begin{enumerate}
\item \textbf{Identifier} la portion de conducteur placée dans le champ magnétique (la bobine dans l'entrefer, les côtés du cadre, etc.) ;
\item \textbf{Orienter} le vecteur $\vec{\ell}$ dans le sens du courant $I$ ;
\item \textbf{Appliquer} $\vec{F} = I\,\vec{\ell} \wedge \vec{B}$ : direction par la règle de la main droite, norme $F = I\ell B \sin\theta$ ;
\item \textbf{Conclure} sur le mouvement : translation si la résultante des forces est non nulle, rotation si les forces forment un couple.
\end{enumerate}
\end{methode}

\begin{exemplebox}[Force sur la bobine d'un haut-parleur]
La bobine mobile d'un haut-parleur comporte $N = 50$ spires de rayon moyen $r = \SI{2,0}{cm}$, placées dans un champ radial de norme $B = \SI{0,40}{T}$. Elle est parcourue par un courant $I = \SI{0,50}{A}$. La longueur totale de fil baignant dans le champ est $\ell_{\text{tot}} = N \times 2\pi r = 50 \times 2\pi \times \num{0,020} \approx \SI{6,28}{m}$. Comme $\vec{\ell} \perp \vec{B}$ en tout point :
\[ F = N\,I\,(2\pi r)\,B = 50 \times \num{0,50} \times 2\pi \times \num{0,020} \times \num{0,40} \approx \SI{1,26}{N}. \]
Une force de plus d'un newton sur une bobine de quelques grammes : l'accélération est considérable, ce qui permet à la membrane de suivre des vibrations rapides (jusqu'à plusieurs milliers de hertz). Vérification dimensionnelle : $[I\ell B] = A \cdot m \cdot T = A \cdot m \cdot \dfrac{N}{A \cdot m} = N$. \checkmark
\end{exemplebox}

\courssub{Le cadre rectangulaire dans un champ uniforme : le couple magnétique}

\textbf{Expérience et observation.}\par
Suspendons un cadre rectangulaire mobile autour d'un axe vertical $\Delta$, parcourons-le d'un courant $I$ et plaçons-le dans le champ magnétique uniforme $\vec{B}$ d'un aimant en U. On observe que le cadre \emph{tourne} et s'oriente, puis s'immobilise dans une position d'équilibre. Pourtant, la résultante des forces de Laplace sur le cadre est nulle... Que se passe-t-il ? Analysons en détail.

\begin{popfigure}
\begin{center}
\begin{tikzpicture}[scale=1.05,>=Stealth]
% Champ B
\foreach \y in {-1.4,-0.7,0,0.7,1.4}{
  \draw[->,popblue,thick] (-3.6,\y) -- (3.6,\y);}
\node[popblue,above right] at (3.6,1.4) {$\vec{B}$ (uniforme)};
% Axe de rotation
\draw[dashed,popdark,thick] (0,-2.3) -- (0,2.3);
\node[popdark,above] at (0,2.3) {$\Delta$};
% Cadre vu en perspective
\draw[poporange,very thick] (-1.6,-1.0) -- (1.2,-0.6) -- (1.2,1.0) -- (-1.6,0.6) -- cycle;
% Courant
\draw[->,poporange,thick] (-1.6,-0.2) -- (-1.6,0.3);
\draw[->,poporange,thick] (1.2,0.2) -- (1.2,-0.3);
\node[poporange,left] at (-1.6,0.05) {\small $I$};
\node[poporange,right] at (1.2,-0.05) {\small $I$};
% Forces sur les côtés actifs
\draw[->,poppink,ultra thick] (-1.6,0) -- (-2.6,0.55);
\node[poppink,above left] at (-2.5,0.5) {$\vec{F}_1$};
\draw[->,poppink,ultra thick] (1.2,0) -- (2.2,-0.55);
\node[poppink,below right] at (2.1,-0.5) {$\vec{F}_2$};
% Couple
\draw[->,popgold,very thick] (0.25,1.5) arc (70:110:0.9);
\node[popgold,above] at (0,1.75) {\small couple};
\end{tikzpicture}
\end{center}
\legende{Cadre rectangulaire parcouru par un courant $I$ dans un champ $\vec{B}$ uniforme. Les forces $\vec{F}_1$ et $\vec{F}_2$, opposées mais non colinéaires, forment un couple qui fait tourner le cadre autour de $\Delta$. Les forces sur les côtés horizontaux sont nulles car $\vec{\ell} \parallel \vec{B}$.}
\end{popfigure}

\textbf{Bilan des forces.}\par
Notons $a$ la longueur des côtés perpendiculaires à l'axe $\Delta$ (les côtés « actifs », perpendiculaires à $\vec{B}$) et $b$ celle des deux autres côtés. Appliquons la méthode :
\begin{itemize}
\item Sur les deux côtés actifs (perpendiculaires à $\vec{B}$, de longueur $a$) : les forces $\vec{F}_1$ et $\vec{F}_2$ ont même norme $F = I a B$ (car $\sin\theta = 1$), des sens opposés (car les courants y circulent en sens opposés), et leurs lignes d'action sont \emph{distinctes}, situées de part et d'autre de l'axe $\Delta$.
\item Sur les deux autres côtés (parallèles à $\vec{B}$, de longueur $b$) : $\vec{\ell} \parallel \vec{B}$, donc $\vec{F}_3 = \vec{F}_4 = \vec{0}$. Il n'y a aucune force sur ces côtés.
\end{itemize}
La résultante totale est $\vec{F}_1 + \vec{F}_2 = \vec{0}$ : le centre de masse du cadre n'est pas accéléré. Mais $\vec{F}_1$ et $\vec{F}_2$ forment un \cle{couple de forces} : deux forces opposées de lignes d'action distinctes. Un couple ne translate pas, il \emph{fait tourner}. C'est le \cle{couple magnétique}.

\textbf{Moment du couple.}\par
Justifions qualitativement l'expression du moment. Le moment d'un couple vaut le produit de la norme commune des forces par la distance $d$ entre leurs lignes d'action (le « bras de levier ») : $\mathcal{M} = F \times d$. Ici $F = I a B$ et $d = b\sin\theta$, où $\theta$ est l'angle entre la normale au plan du cadre et $\vec{B}$ (le bras de levier est maximal quand le plan du cadre est parallèle à $\vec{B}$, et nul quand le cadre est perpendiculaire à $\vec{B}$). Donc :
\[ \mathcal{M} = I a B \times b \sin\theta = B I (a b) \sin\theta = B I S \sin\theta, \]
où $S = ab$ est la surface du cadre. Pour une bobine de $N$ spires, les moments s'ajoutent.

\begin{propriete}[Moment du couple magnétique sur un cadre — admis]
Un cadre (ou bobine plate) de $N$ spires, de surface totale $S$ par spire, parcouru par un courant $I$ et placé dans un champ magnétique uniforme $\vec{B}$, subit un couple de forces de Laplace dont le moment par rapport à l'axe de rotation est :
\[ \mathcal{M} = N\,B\,I\,S\,\sin\theta \]
où $\theta$ est l'angle entre la normale au plan du cadre et $\vec{B}$. La résultante des forces est nulle : seul le couple subsiste. Le moment est \emph{maximal} ($\mathcal{M}_{\max} = NBIS$) quand le plan du cadre est parallèle à $\vec{B}$ ($\theta = 90\degres$), et \emph{nul} quand le plan du cadre est perpendiculaire à $\vec{B}$ ($\theta = 0$) : c'est la position d'équilibre. Cette expression est admise ; la justification qualitative à partir des forces sur les côtés actifs est exigible.
\end{propriete}

\begin{attention}[La résultante est nulle, pas le couple !]
Erreur très fréquente au Bac : affirmer qu'un cadre fermé parcouru par un courant dans un champ uniforme est « attiré » ou « repoussé » par l'aimant. Dans un champ \emph{uniforme}, la somme vectorielle des forces de Laplace sur un circuit fermé complet est \textbf{toujours nulle} : $\sum \vec{F} = \vec{0}$. Le cadre ne se déplace pas en translation ; il subit uniquement un \textbf{couple} qui le fait pivoter. Ne confonds jamais « résultante nulle » (pas de translation) avec « aucun effet » (le couple fait tourner !). Remarque au passage : c'est parce que le champ de l'aimant du haut-parleur est radial et non uniforme sur un circuit plan que la bobine, elle, subit une force nette axiale.
\end{attention}

\courssub{Le moteur à courant continu}

\textbf{Le problème à résoudre.}\par
Le couple magnétique fait tourner le cadre... mais pas longtemps ! Parti de la position où le couple est maximal, le cadre tourne jusqu'à la position d'équilibre ($\theta = 0$, plan du cadre perpendiculaire à $\vec{B}$) où le couple s'annule ; s'il la dépasse par inertie, le couple change de signe et le ramène en arrière. Le cadre oscille puis s'immobilise. Pour obtenir une \emph{rotation continue}, il faut un dispositif ingénieux : le \cle{collecteur}.

\textbf{Constitution et principe.}\par
Le moteur à courant continu comprend :
\begin{itemize}
\item un \cle{stator} (partie fixe) : aimant ou électro-aimant créant le champ $\vec{B}$ ;
\item un \cle{rotor} (partie tournante) : bobine de $N$ spires montée sur un axe ;
\item un \cle{collecteur} : anneau cylindrique fendu en deux demi-anneaux, solidaire du rotor et tournant avec lui ; chaque demi-anneau est relié à une extrémité de la bobine ;
\item deux \cle{balais} (charbons) fixes, frottant sur le collecteur, qui relient le rotor au générateur extérieur.
\end{itemize}

\begin{popfigure}
\begin{center}
\begin{tikzpicture}[scale=1.0,>=Stealth]
% Stator (aimants)
\fill[popblueL] (-3.4,-1.9) rectangle (-2.2,1.9);
\fill[popblueL] (2.2,-1.9) rectangle (3.4,1.9);
\draw[popdark,thick] (-3.4,-1.9) rectangle (-2.2,1.9);
\draw[popdark,thick] (2.2,-1.9) rectangle (3.4,1.9);
\node[popdark] at (-2.8,0) {\small N};
\node[popdark] at (2.8,0) {\small S};
% Champ B
\foreach \y in {-1.2,-0.4,0.4,1.2}{
  \draw[->,popblue,thick] (-2.1,\y) -- (2.1,\y);}
\node[popblue,above] at (0,1.35) {$\vec{B}$};
% Rotor : cadre
\draw[poporange,very thick,rotate=25] (-1.3,-0.55) rectangle (1.3,0.55);
\node[poporange] at (0.1,1.15) {\small rotor (bobine de $N$ spires)};
% Axe
\draw[popdark,thick] (0,-2.6) -- (0,-1.6);
\fill[popdark] (0,0) circle (0.06);
% Collecteur (deux demi-anneaux)
\draw[popgold,very thick] (-0.28,-1.75) arc (180:360:0.28 and 0.12);
\draw[popgold,very thick] (0.28,-1.95) arc (0:180:0.28 and 0.12);
\draw[popgold,thick] (0,-1.35) -- (0,-1.6);
\node[popgold,right] at (0.35,-1.85) {\small collecteur};
% Balais
\fill[popmuted] (-0.55,-2.25) rectangle (-0.3,-2.0);
\fill[popmuted] (0.3,-2.25) rectangle (0.55,-2.0);
\node[popmuted,below] at (0,-2.3) {\small balais};
% Circuit
\draw[popdark,thick] (-0.42,-2.25) -- (-0.42,-2.9) -- (0.42,-2.9) -- (0.42,-2.25);
\draw[popdark,thick] (-0.25,-2.9) -- (0.05,-2.9);
\draw[popdark,thick] (-0.17,-3.0) -- (-0.03,-3.0);
\node[popdark,below] at (-0.1,-3.05) {\small $+$};
\node[popdark] at (0.75,-2.9) {\small $-$};
% Rotation
\draw[->,poppink,very thick] (1.7,0.9) arc (40:-40:1.15);
\node[poppink,right] at (2.0,0.55) {\small rotation};
\end{tikzpicture}
\end{center}
\legende{Schéma de principe du moteur à courant continu : le couple magnétique fait tourner le rotor ; le collecteur, via les balais, inverse le sens du courant à chaque demi-tour.}
\end{popfigure}

\textbf{Rôle du collecteur : l'inversion à chaque demi-tour.}\par
Voici l'idée géniale. Quand le rotor franchit la position d'équilibre ($\theta = 0$), chaque balai passe d'un demi-anneau à l'autre : le sens du courant dans la bobine \emph{s'inverse automatiquement}. Or, inverser $I$ inverse le sens de toutes les forces de Laplace, donc le sens du couple. Résultat : au lieu de devenir un couple de rappel, le couple reste \emph{moteur}, toujours dans le même sens de rotation. À chaque demi-tour, l'opération se répète : le rotor est relancé en permanence et tourne de façon continue. Sans collecteur, pas de moteur !

\textbf{Bilan énergétique.}\par
Le moteur reçoit de l'énergie électrique du générateur (puissance $P_{\text{élec}} = UI$) et la convertit en énergie mécanique de rotation (puissance $\mathcal{M}\,\omega$, où $\omega$ est la vitesse angulaire), avec des pertes par effet Joule dans les bobinages et par frottements :
\[ \text{énergie électrique} \;\longrightarrow\; \text{énergie mécanique} + \text{pertes (chaleur)}. \]
C'est cette conversion qui fait tourner les pompes à eau des forages villageois, les ventilateurs pendant la saison sèche à Garoua, les broyeuses de manioc et les moteurs des mototaxis électriques qui commencent à circuler à Yaoundé.

\begin{savaistu}[De la broyeuse de manioc au barrage de Song Loulou]
La force de Laplace est à l'origine d'une symétrie magnifique de la physique. Dans le \textbf{moteur}, un courant placé dans un champ magnétique subit des forces qui produisent du mouvement : électricité $\rightarrow$ mouvement. Mais le phénomène inverse existe : en imposant un mouvement à un conducteur dans un champ magnétique, on fait naître un courant — c'est l'\emph{induction électromagnétique}, que tu étudieras dans une leçon ultérieure. C'est le principe des \textbf{alternateurs} des barrages de Song Loulou et d'Edéa sur la Sanaga : l'eau fait tourner d'énormes turbines solidaires de bobines dans des champs magnétiques intenses, et le mouvement devient électricité. Mouvement $\rightarrow$ électricité. Moteur et alternateur sont deux visages du même couple électromagnétique : l'un consomme le courant pour créer le couple, l'autre utilise le couple pour créer le courant.
\end{savaistu}

\begin{exoresolu}[Couple moteur d'une petite pompe villageoise]
Le rotor du moteur à courant continu d'une pompe à eau est une bobine de $N = 200$ spires rectangulaires de dimensions $a = \SI{5,0}{cm}$ et $b = \SI{4,0}{cm}$, placée dans un champ magnétique uniforme de norme $B = \SI{0,25}{T}$. Le courant d'alimentation vaut $I = \SI{2,0}{A}$.
\begin{enumerate}
\item Calculer le moment maximal du couple magnétique.
\item Le rotor tourne à la vitesse angulaire $\omega = \SI{100}{rad.s^{-1}}$. En supposant que le couple utile vaut en moyenne la moitié du couple maximal, calculer la puissance mécanique fournie.
\end{enumerate}
\tcblower
\textbf{1.} La surface d'une spire est $S = a \times b = \num{0,050} \times \num{0,040} = \SI{2,0e-3}{m^2}$. Le moment maximal correspond à $\sin\theta = 1$ :
\[ \mathcal{M}_{\max} = N B I S = 200 \times \num{0,25} \times \num{2,0} \times \num{2,0e-3} = \SI{0,20}{N.m}. \]
Vérification dimensionnelle : $T \cdot A \cdot m^2 = \dfrac{N}{A \cdot m} \cdot A \cdot m^2 = N \cdot m$. \checkmark

\textbf{2.} Le couple utile moyen vaut $\mathcal{M} = \dfrac{\mathcal{M}_{\max}}{2} = \SI{0,10}{N.m}$. La puissance mécanique d'un couple en rotation est $P = \mathcal{M}\,\omega$ :
\[ P = \num{0,10} \times 100 = \SI{10}{W}. \]
La pompe fournit une puissance mécanique de \SI{10}{W}, suffisante pour remonter l'eau d'un petit forage. La puissance électrique absorbée est supérieure à cause des pertes par effet Joule dans les bobinages.
\end{exoresolu}

\begin{aretenir}[L'essentiel de la section]
\begin{itemize}
\item \textbf{Haut-parleur} : la bobine mobile, placée dans le champ radial de l'aimant, subit une force de Laplace $F = N I (2\pi r) B$ proportionnelle au courant ; le courant variable du signal sonore fait vibrer la membrane, qui émet le son.
\item \textbf{Cadre dans un champ uniforme} : la résultante des forces de Laplace est nulle, mais les forces sur les côtés actifs forment un couple de moment $\mathcal{M} = N B I S \sin\theta$, maximal quand le plan du cadre est parallèle à $\vec{B}$.
\item \textbf{Moteur à courant continu} : le couple magnétique fait tourner le rotor ; le collecteur inverse le courant à chaque demi-tour pour que le couple reste toujours moteur. Conversion : énergie électrique $\rightarrow$ énergie mécanique.
\item \textbf{Démarche universelle} : identifier le conducteur dans le champ, orienter $\vec{\ell}$ dans le sens de $I$, appliquer $\vec{F} = I\,\vec{\ell} \wedge \vec{B}$, conclure sur le mouvement.
\end{itemize}
\end{aretenir}

\coursec{Interaction entre deux conducteurs parallèles et définition de l'ampère}

Dans la section précédente, nous avons vu comment la force de Laplace met en mouvement : le haut-parleur transforme un courant variable en vibration sonore, le moteur électrique transforme l'énergie électrique en rotation. Dans les deux cas, un conducteur parcouru par un courant est placé dans un champ magnétique \emph{extérieur}, créé par un aimant ou une bobine. Mais une question naturelle se pose alors : un courant crée un champ magnétique (section 2), et un courant placé dans un champ magnétique subit une force (section 4). Que se passe-t-il donc lorsque \emph{deux} courants sont proches l'un de l'autre ? Chacun baigne dans le champ magnétique créé par l'autre : les deux conducteurs doivent s'exercer mutuellement des forces. C'est exactement ce que l'expérience confirme, et c'est cette interaction qui a servi, pendant plus d'un demi-siècle, à \emph{définir} l'unité d'intensité du courant électrique : l'ampère.

\courssub{Mise en évidence expérimentale}

\begin{experience}[Deux fils parallèles parcourus par des courants]
\textbf{Matériel :} deux fils de cuivre souples et légers, suspendus verticalement et parallèlement, distants de quelques centimètres ; un générateur capable de débiter un courant intense (quelques dizaines d'ampères) ; des interrupteurs permettant de choisir le sens du courant dans chaque fil.

\textbf{Protocole :}
\begin{enumerate}
\item On fait passer dans les deux fils des courants \cle{de même sens} (montants tous les deux). On observe les fils.
\item On inverse le sens du courant dans l'un des fils : les courants sont maintenant \cle{de sens opposés}. On observe à nouveau.
\end{enumerate}

\textbf{Observations :}
\begin{itemize}
\item Courants de même sens : les deux fils \textbf{se rapprochent} et finissent par se toucher : ils \textbf{s'attirent}.
\item Courants de sens opposés : les deux fils \textbf{s'écartent} l'un de l'autre : ils \textbf{se repoussent}.
\end{itemize}

\textbf{Interprétation :} les fils sont en cuivre, matériau non magnétique : il ne s'agit pas d'une interaction entre aimants. Les fils ne portent pas de charge électrique nette : il ne s'agit pas non plus d'une interaction électrostatique. L'interaction observée est purement \cle{magnétique} : le courant $I_1$ crée autour du fil 1 un champ magnétique $\vec{B}_1$ ; le fil 2, parcouru par $I_2$, est plongé dans ce champ et subit donc une force de Laplace. Symétriquement, le fil 1 subit la force de Laplace due au champ $\vec{B}_2$ créé par le fil 2.
\end{experience}

\begin{attention}[Ce que cette expérience n'est pas]
Ne confonds pas cette interaction avec l'attraction électrostatique entre charges de signes opposés. Un fil parcouru par un courant reste \emph{électriquement neutre} (autant d'électrons que de charges positives fixes dans le réseau cristallin). La force mise en jeu ici est d'origine \textbf{magnétique} : elle n'existe que parce que les charges sont \emph{en mouvement}. Coupe le courant : la force disparaît instantanément.
\end{attention}

\courssub{Interprétation théorique : combinaison du champ du fil et de la force de Laplace}

Reprenons les deux résultats fondamentaux établis dans les sections précédentes :

\begin{itemize}
\item \textbf{Champ créé par un fil rectiligne infini} (section 2) : un fil rectiligne parcouru par un courant $I_1$ crée, à la distance $d$, un champ magnétique orthoradial (dont les lignes de champ sont des cercles centrés sur le fil) de norme :
\[ B_1 = \frac{\mu_0\, I_1}{2\pi d} \]
\item \textbf{Force de Laplace} (section 4) : un conducteur de longueur $\ell$, parcouru par un courant $I_2$, placé dans un champ $\vec{B}_1$, subit la force :
\[ \vec{F}_{1\to 2} = I_2\, \vec{\ell} \wedge \vec{B}_1 \]
\end{itemize}

Le fil 2 est situé à la distance $d$ du fil 1 : il baigne donc dans le champ $\vec{B}_1$. Comme les fils sont parallèles, le champ $\vec{B}_1$ est, en tout point du fil 2, \emph{perpendiculaire} au fil 2. Le produit vectoriel se simplifie alors : $\|\vec{\ell} \wedge \vec{B}_1\| = \ell\, B_1 \sin 90^\circ = \ell B_1$. Il suffit de combiner les deux formules.

\begin{propriete}[Force d'interaction entre deux fils parallèles]
Deux conducteurs rectilignes, parallèles, de longueur $\ell$ (grande devant leur écartement), distants de $d$, parcourus par des courants $I_1$ et $I_2$, exercent l'un sur l'autre des forces de norme :
\[ F = I_2\, \ell\, B_1 = \frac{\mu_0\, I_1\, I_2\, \ell}{2\pi d} \]
soit, par unité de longueur :
\[ \boxed{\ \frac{F}{\ell} = \frac{\mu_0\, I_1\, I_2}{2\pi d}\ } \]
Ces deux forces forment un couple d'\cle{action-réaction} : $\vec{F}_{1\to 2} = -\vec{F}_{2\to 1}$, conformément à la troisième loi de Newton.

\textbf{Démonstration guidée (exigible) :}
\begin{enumerate}
\item Le fil 1 crée à la distance $d$ (là où se trouve le fil 2) le champ $B_1 = \dfrac{\mu_0 I_1}{2\pi d}$, perpendiculaire au plan contenant les deux fils.
\item Le fil 2, parcouru par $I_2$ et de longueur $\ell$, plongé dans ce champ perpendiculaire, subit la force de Laplace de norme $F_{1\to 2} = I_2\, \ell\, B_1\, \sin 90^\circ = I_2 \ell B_1$.
\item En remplaçant $B_1$ par son expression : $F_{1\to 2} = I_2 \ell \times \dfrac{\mu_0 I_1}{2\pi d} = \dfrac{\mu_0 I_1 I_2 \ell}{2\pi d}$.
\item En divisant par $\ell$, on obtient la force par unité de longueur, indépendante de $\ell$ : $\dfrac{F}{\ell} = \dfrac{\mu_0 I_1 I_2}{2\pi d}$.
\item Le même raisonnement appliqué au fil 1 plongé dans $B_2 = \dfrac{\mu_0 I_2}{2\pi d}$ donne $F_{2\to 1} = \dfrac{\mu_0 I_1 I_2 \ell}{2\pi d} = F_{1\to 2}$ : les normes sont égales, les sens opposés.
\end{enumerate}

\textbf{Vérification dimensionnelle :} $\mu_0$ s'exprime en $\mathrm{N\,A^{-2}}$, donc $\dfrac{[\mu_0]\,[I_1]\,[I_2]}{[d]} = \dfrac{\mathrm{N\,A^{-2}} \times \mathrm{A} \times \mathrm{A}}{\mathrm{m}} = \mathrm{N\,m^{-1}}$, qui est bien une force par unité de longueur.
\end{propriete}

\begin{attention}[Pièges fréquents]
\begin{itemize}
\item La formule $B = \dfrac{\mu_0 I}{2\pi d}$ suppose des fils \textbf{longs} devant $d$ (modèle du fil infini). Elle ne s'applique pas à deux petits tronçons de fils.
\item N'oublie pas le facteur $2\pi$ au dénominateur : il vient de la circonférence des lignes de champ circulaires autour du fil.
\item $F/\ell$ ne dépend \textbf{pas} de la longueur des fils : c'est bien une force \emph{par unité de longueur}. Pour obtenir la force totale sur une portion de longueur $\ell$, multiplie par $\ell$.
\item Attention à l'homogénéité : $d$ doit être exprimé en \emph{mètres} dans les applications numériques.
\end{itemize}
\end{attention}

\courssub{Analyse vectorielle : attraction ou répulsion ?}

La norme de la force ne dit pas si les fils s'attirent ou se repoussent : c'est l'\emph{orientation} des vecteurs, obtenue par la règle de la main droite appliquée au produit vectoriel $\vec{F} = I\,\vec{\ell} \wedge \vec{B}$, qui tranche.

\begin{propriete}[Règle d'interaction entre courants parallèles]
\begin{itemize}
\item Deux courants parallèles \cle{de même sens} \textbf{s'attirent}.
\item Deux courants parallèles \cle{de sens opposés} \textbf{se repoussent}.
\end{itemize}

\textbf{Vérification par la règle de la main droite (cas de même sens) :} plaçons le fil 1 à gauche, le fil 2 à droite, les deux courants sortant du plan de la feuille (symbole $\odot$).
\begin{enumerate}
\item \textbf{Direction de $\vec{B}_1$ au niveau du fil 2 :} les lignes de champ du fil 1 sont des cercles centrés sur le fil 1, orientés par la règle du tire-bouchon (ou de la main droite : pouce dans le sens de $I_1$, les doigts s'enroulent dans le sens de $\vec{B}_1$). À droite du fil 1, le champ $\vec{B}_1$ est dirigé \emph{vers le haut} de la feuille.
\item \textbf{Force sur le fil 2 :} $\vec{F}_{1\to 2} = I_2\,\vec{\ell} \wedge \vec{B}_1$. Avec $\vec{\ell}$ sortant de la feuille et $\vec{B}_1$ vers le haut, la règle de la main droite donne $\vec{F}_{1\to 2}$ dirigée \emph{vers la gauche}, c'est-à-dire \textbf{vers le fil 1} : attraction.
\item Par symétrie (action-réaction), $\vec{F}_{2\to 1}$ est dirigée vers la droite, vers le fil 2.
\end{enumerate}
Si l'on inverse le sens de $I_2$ (courants de sens opposés), le produit vectoriel change de signe : la force sur le fil 2 s'oriente \emph{vers la droite}, donc \textbf{à l'opposé du fil 1} : répulsion. Retiens bien le contraste avec l'électrostatique : en électricité, les charges de \emph{même signe} se repoussent ; en magnétisme, les courants de \emph{même sens} s'attirent.
\end{propriete}

\begin{popfigure}
\begin{center}
\begin{tikzpicture}[scale=1.0, >=Stealth]
% ---- Cas 1 : courants de même sens (attraction) ----
\begin{scope}[xshift=-4.2cm]
  \node[font=\bfseries] at (0,2.6) {Même sens : attraction};
  % Fil 1
  \draw[thick, popblue, fill=popblueL] (-1.6,0) circle (0.28);
  \node[popblue] at (-1.6,0) {\large $\odot$};
  \node[below, popblue] at (-1.6,-0.35) {fil 1 : $I_1$};
  % Fil 2
  \draw[thick, poporange, fill=poporangeL] (1.6,0) circle (0.28);
  \node[poporange] at (1.6,0) {\large $\odot$};
  \node[below, poporange] at (1.6,-0.35) {fil 2 : $I_2$};
  % distance d
  \draw[<->, popmuted] (-1.6,-0.85) -- (1.6,-0.85) node[midway, below] {$d$};
  % Champ B1 au niveau du fil 2 (vers le haut)
  \draw[->, very thick, popgreen] (1.6,0.45) -- (1.6,1.55) node[right] {$\vec{B}_1$};
  % Champ B2 au niveau du fil 1 (vers le bas)
  \draw[->, very thick, popgreen] (-1.6,-0.45) -- (-1.6,-1.55) node[right] {$\vec{B}_2$};
  % Forces
  \draw[->, very thick, poppink] (1.25,0) -- (0.35,0) node[midway, above] {$\vec{F}_{1\to 2}$};
  \draw[->, very thick, poppink] (-1.25,0) -- (-0.35,0) node[midway, above] {$\vec{F}_{2\to 1}$};
\end{scope}
% ---- Cas 2 : courants de sens opposés (répulsion) ----
\begin{scope}[xshift=4.2cm]
  \node[font=\bfseries] at (0,2.6) {Sens opposés : répulsion};
  % Fil 1
  \draw[thick, popblue, fill=popblueL] (-1.6,0) circle (0.28);
  \node[popblue] at (-1.6,0) {\large $\odot$};
  \node[below, popblue] at (-1.6,-0.35) {fil 1 : $I_1$};
  % Fil 2
  \draw[thick, poporange, fill=poporangeL] (1.6,0) circle (0.28);
  \node[poporange] at (1.6,0) {\large $\otimes$};
  \node[below, poporange] at (1.6,-0.35) {fil 2 : $I_2$};
  % distance d
  \draw[<->, popmuted] (-1.6,-0.85) -- (1.6,-0.85) node[midway, below] {$d$};
  % Champ B1 au niveau du fil 2 (vers le haut)
  \draw[->, very thick, popgreen] (1.6,0.45) -- (1.6,1.55) node[right] {$\vec{B}_1$};
  % Champ B2 au niveau du fil 1 (vers le bas)
  \draw[->, very thick, popgreen] (-1.6,-0.45) -- (-1.6,-1.55) node[right] {$\vec{B}_2$};
  % Forces (repoussées vers l'extérieur)
  \draw[->, very thick, poppink] (1.95,0) -- (2.85,0) node[midway, above] {$\vec{F}_{1\to 2}$};
  \draw[->, very thick, poppink] (-1.95,0) -- (-2.85,0) node[midway, above] {$\vec{F}_{2\to 1}$};
\end{scope}
\end{tikzpicture}
\end{center}
\legende{Deux fils parallèles vus en coupe ($\odot$ : courant sortant, $\otimes$ : courant entrant). À gauche, les courants de même sens s'attirent ; à droite, les courants de sens opposés se repoussent. Les champs $\vec{B}_1$ et $\vec{B}_2$ sont tangents aux lignes de champ circulaires de chaque fil.}
\end{popfigure}

\courssub{Application numérique : des forces faibles mais non négligeables}

\begin{exemplebox}[Deux câbles d'une installation industrielle]
Deux câbles parallèles d'une installation électrique sont distants de $d = \SI{5}{cm} = \SI{0,05}{m}$ et parcourus chacun par un courant $I_1 = I_2 = \SI{100}{A}$. Calculons la force par unité de longueur qu'ils exercent l'un sur l'autre :
\begin{align*}
\frac{F}{\ell} &= \frac{\mu_0\, I_1\, I_2}{2\pi d} = \frac{4\pi \times 10^{-7} \times 100 \times 100}{2\pi \times 0{,}05} \\
&= \frac{4\pi \times 10^{-7} \times 10^{4}}{0{,}1\pi} = \frac{4\pi \times 10^{-3}}{0{,}1\pi} = 4 \times 10^{-2}\ \mathrm{N\,m^{-1}}
\end{align*}
\[ \boxed{\ \frac{F}{\ell} = \SI{0,04}{N/m}\ } \]
Sur un tronçon de $\ell = \SI{10}{m}$, cela représente $F = 0{,}4\ \mathrm{N}$, soit le poids d'une masse d'environ $40\ \mathrm{g}$. C'est modeste en régime normal... mais lors d'un \emph{court-circuit}, le courant peut atteindre $10\,000\ \mathrm{A}$ : la force est alors multipliée par $\left(\frac{10\,000}{100}\right)^2 = 10^{4}$, car elle est \textbf{proportionnelle au produit} $I_1 I_2$ ! On obtient $F/\ell = \SI{400}{N/m}$ : de quoi déformer violemment les câbles et arracher leurs supports.
\end{exemplebox}

\begin{savaistu}[Les efforts électrodynamiques dans les réseaux d'Eneo]
Dans les postes de transformation et les lignes à haute tension exploités par Eneo et la Sonatrel (comme ceux qui évacuent l'énergie des barrages de Nachtigal ou de Lom Pangar), les barres omnibus et les câbles parallèles véhiculent des centaines d'ampères en régime normal. Lors d'un court-circuit — par exemple une branche tombée sur une ligne — le courant peut dépasser $\SI{20}{kA}$ pendant une fraction de seconde. Les forces de Laplace entre conducteurs voisins croissent alors comme le \emph{carré} du courant et peuvent atteindre plusieurs milliers de newtons par mètre : les câbles claquent violemment les uns contre les autres ou s'écartent brutalement. C'est pourquoi les ingénieurs \emph{dimensionnent mécaniquement} les supports, entretoises et isolateurs pour résister à ces efforts électrodynamiques, et pourquoi les disjoncteurs doivent couper le courant en quelques dizaines de millisecondes.
\end{savaistu}

\courssub{La définition historique de l'ampère}

La formule $\dfrac{F}{\ell} = \dfrac{\mu_0 I_1 I_2}{2\pi d}$ possède une vertu remarquable : elle relie une grandeur électrique (l'intensité $I$) à des grandeurs purement mécaniques (une force, une longueur), que l'on sait mesurer avec une balance et une règle. Les physiciens ont donc renversé la perspective : au lieu de calculer la force à partir des courants, on peut \emph{définir} l'unité de courant à partir de la force. C'est ce qu'a fait la Conférence générale des poids et mesures en 1948.

\begin{definition}[L'ampère (définition historique, 1948)]
L'\cle{ampère} (symbole A) est l'intensité d'un courant constant qui, maintenu dans deux conducteurs parallèles, rectilignes, de longueur infinie, de section circulaire négligeable et placés à une distance de $\SI{1}{m}$ l'un de l'autre dans le vide, produit entre ces conducteurs une force égale à $2 \times 10^{-7}\ \mathrm{N}$ par mètre de longueur.
\end{definition}

Vérifions la cohérence de cette définition avec la valeur de $\mu_0$. Posons $I_1 = I_2 = \SI{1}{A}$ et $d = \SI{1}{m}$ dans la formule :
\[ \frac{F}{\ell} = \frac{\mu_0 \times 1 \times 1}{2\pi \times 1} = \frac{\mu_0}{2\pi} \]
Pour que cette force vaille exactement $2\times 10^{-7}\ \mathrm{N\,m^{-1}}$, il faut :
\[ \frac{\mu_0}{2\pi} = 2\times 10^{-7}\ \mathrm{N\,A^{-2}} \quad \Longrightarrow \quad \boxed{\ \mu_0 = 4\pi \times 10^{-7}\ \mathrm{N\,A^{-2}}\ } \]
La valeur de la perméabilité magnétique du vide, $\mu_0 = 4\pi \times 10^{-7}$ S.I. (soit environ $1{,}257 \times 10^{-6}\ \mathrm{H\,m^{-1}}$), n'est donc pas un hasard : elle était \emph{fixée exactement} par la définition de l'ampère. L'ampère était l'unité de base, et $\mu_0$ en découlait.

\begin{savaistu}[La nouvelle définition de l'ampère (complément culturel)]
Depuis le 20 mai 2019, le Système international d'unités a été profondément réformé : l'ampère n'est plus défini par l'expérience des deux fils, mais en fixant la valeur de la \textbf{charge élémentaire} : $e = 1{,}602\,176\,634 \times 10^{-19}\ \mathrm{C}$ exactement. Un ampère correspond alors au passage d'environ $6{,}24 \times 10^{18}$ charges élémentaires par seconde ($I = \dfrac{\Delta Q}{\Delta t}$, avec $\SI{1}{A} = \SI{1}{C/s}$). Conséquence : $\mu_0$ n'est plus une constante fixée exactement, mais une grandeur \emph{mesurée} (sa valeur reste $4\pi \times 10^{-7}$ S.I. à une précision relative meilleure que $10^{-9}$). Au niveau du Baccalauréat, tu continueras d'utiliser $\mu_0 = 4\pi \times 10^{-7}$ S.I. sans aucun souci, et l'expérience des deux fils parallèles demeure la meilleure illustration concrète de ce qu'est un ampère.
\end{savaistu}

\courssub{Exercice résolu de synthèse}

\begin{exoresolu}[Balance de Cotton simplifiée : mesurer un courant par une force]
Deux rails conducteurs parallèles, très longs, sont distants de $d = \SI{10}{cm}$ dans le vide. Ils sont parcourus par le même courant $I$, dans le même sens. On mesure, sur une portion de longueur $\ell = \SI{50}{cm}$ de l'un des rails, une force d'attraction $F = 2{,}0 \times 10^{-4}\ \mathrm{N}$.
\begin{enumerate}
\item Calculer la force par unité de longueur $F/\ell$.
\item En déduire l'intensité $I$ du courant.
\item Que devient cette force si l'on double l'intensité du courant, la distance restant inchangée ?
\end{enumerate}
\tcblower
\textbf{Solution détaillée.}
\begin{enumerate}
\item La force par unité de longueur est :
\[ \frac{F}{\ell} = \frac{2{,}0 \times 10^{-4}}{0{,}50} = 4{,}0 \times 10^{-4}\ \mathrm{N\,m^{-1}} \]
\item On applique la loi d'interaction avec $I_1 = I_2 = I$ :
\[ \frac{F}{\ell} = \frac{\mu_0 I^2}{2\pi d} \quad \Longrightarrow \quad I^2 = \frac{2\pi d}{\mu_0}\times \frac{F}{\ell} = \frac{2\pi \times 0{,}10}{4\pi \times 10^{-7}} \times 4{,}0 \times 10^{-4} \]
\[ I^2 = \frac{0{,}10}{2 \times 10^{-7}} \times 4{,}0 \times 10^{-4} = 5{,}0 \times 10^{5} \times 4{,}0 \times 10^{-4} = 200\ \mathrm{A^2} \]
d'où $I = \sqrt{200} \approx \SI{14}{A}$. (On vérifie l'homogénéité : $\mathrm{m} \times \mathrm{N\,m^{-1}} / (\mathrm{N\,A^{-2}}) = \mathrm{A^2}$.)
\item La force est proportionnelle à $I^2$ : doubler $I$ quadruple la force. La nouvelle force par unité de longueur vaut $4 \times 4{,}0 \times 10^{-4} = 1{,}6 \times 10^{-3}\ \mathrm{N\,m^{-1}}$, soit $F' = 8{,}0 \times 10^{-4}\ \mathrm{N}$ sur la portion de $\SI{50}{cm}$.
\end{enumerate}
\textbf{Remarque :} c'est précisément le principe des anciennes \emph{balances de courant} des laboratoires de métrologie : on « pesait » littéralement l'ampère.
\end{exoresolu}

\begin{aretenir}[L'essentiel de la section]
\begin{itemize}
\item Deux fils parallèles parcourus par des courants exercent l'un sur l'autre des forces de Laplace, car chacun baigne dans le champ magnétique créé par l'autre.
\item Force par unité de longueur : $\dfrac{F}{\ell} = \dfrac{\mu_0 I_1 I_2}{2\pi d}$ — obtenue en combinant $B_1 = \dfrac{\mu_0 I_1}{2\pi d}$ et $F = I_2 \ell B_1$.
\item Courants de \textbf{même sens} : attraction ; courants de \textbf{sens opposés} : répulsion (à retrouver par la règle de la main droite).
\item Définition historique de l'ampère : le courant qui, dans deux fils infinis distants de $\SI{1}{m}$ dans le vide, produit $2 \times 10^{-7}\ \mathrm{N}$ par mètre ; elle fixait $\mu_0 = 4\pi \times 10^{-7}$ S.I. Depuis 2019, l'ampère est défini via la charge élémentaire $e$ fixée.
\item La force croît comme le \emph{produit} des courants : en cas de court-circuit, les efforts électrodynamiques deviennent destructeurs (dimensionnement des installations d'Eneo).
\end{itemize}
\end{aretenir}

\newpage

\coursec{Activité d'intégration}

\courssub{Objectifs de l'activité}

Cette activité d'intégration te place dans la peau d'un jeune technicien en électronique. À travers la conception et l'analyse d'un mini-haut-parleur, tu devras mobiliser l'ensemble des savoir-faire de la leçon :

\begin{itemize}
\item représenter le vecteur champ magnétique $\vec{B}$ créé par un aimant dans son entrefer et retrouver sa direction et son sens ;
\item appliquer la règle de la main droite pour déterminer le sens de la force de Laplace ;
\item calculer la norme de la force de Laplace s'exerçant sur un conducteur bobiné ;
\item exploiter le caractère sinusoïdal d'un courant pour relier force, mouvement et son émis ;
\item calculer le champ magnétique créé par une bobine plate et le comparer à celui d'un aimant.
\end{itemize}

La restitution se fera sous la forme d'un \cle{poster technique annoté} : schéma en coupe, calculs justifiés, conclusions rédigées en phrases complètes.

\courssub{Situation}

\textbf{Le club scientifique du lycée de Ngaoundéré relève un défi.}

Dans le cadre de la semaine nationale de la science, le club de physique de ton lycée a décidé de fabriquer un mini-haut-parleur capable de diffuser la note « la » ($\SI{440}{Hz}$) pour accompagner la chorale de l'établissement. Le proviseur a mis à disposition un vieux haut-parleur démonté, et le club dispose de sa fiche technique simplifiée :

\begin{itemize}
\item un \cle{aimant annulaire} (en forme d'anneau) créant dans son entrefer un champ magnétique \emph{radial}, uniforme, de norme $B = \SI{0,50}{T}$ : en tout point de l'entrefer, $\vec{B}$ est dirigé du centre vers l'extérieur, perpendiculairement à l'axe du haut-parleur ;
\item une \cle{bobine mobile} solidaire d'une membrane (cône en carton), constituée de $N = 80$ spires circulaires de rayon moyen $r = \SI{1,5}{cm}$, plongée dans l'entrefer de l'aimant ;
\item un amplificateur délivrant à la bobine un courant sinusoïdal :
\[ i(t) = 0{,}8 \cos(2\pi \times 440\, t) \quad \text{(en ampères, $t$ en secondes)}. \]
\end{itemize}

On rappelle la perméabilité magnétique du vide : $\mu_0 = 4\pi \times 10^{-7}\ \mathrm{H\,m^{-1}}$.

Le principe est le suivant : le courant $i(t)$ circule dans les spires de la bobine ; chaque portion de fil plongée dans le champ radial $\vec{B}$ subit une force de Laplace dirigée selon l'axe du haut-parleur ; la bobine et la membrane se déplacent alors d'avant en arrière, compriment l'air et produisent une onde sonore.

\textbf{Mission du club :} avant de brancher l'amplificateur, il faut dimensionner le dispositif, prévoir le sens de déplacement de la membrane et vérifier que le champ créé par la bobine elle-même ne perturbe pas le fonctionnement. C'est ton équipe qui est chargée de l'étude !

\courssub{Consignes}

\begin{enumerate}
\item \textbf{Schéma en coupe.} Représente le haut-parleur en coupe par un plan contenant son axe : l'aimant annulaire, l'entrefer, une spire de la bobine vue en coupe (deux sections de fil, courant entrant $\otimes$ d'un côté et sortant $\odot$ de l'autre). À un instant où $i > 0$, place sur ce schéma les vecteurs $\vec{B}$, le sens du courant $\vec{I}$ et la force de Laplace $\vec{F}$ sur chaque section de fil. Justifie le sens de $\vec{F}$ par la règle de la main droite.
\item \textbf{Longueur de fil dans le champ.} Calcule la longueur totale $\ell$ de fil de la bobine plongée dans le champ magnétique.
\item \textbf{Force sur la bobine.} Donne l'expression de la force de Laplace totale $F(t)$ s'exerçant sur la bobine en fonction de $i(t)$, puis calcule sa valeur maximale $F_{\max}$. Commente l'ordre de grandeur obtenu.
\item \textbf{Du courant au son.} Détermine la fréquence $f$ de vibration de la membrane. Justifie que le haut-parleur émet bien la note « la ».
\item \textbf{Inversion des bornes.} Explique qualitativement ce qui se passe si l'on inverse les bornes de l'alimentation de la bobine. Le son produit est-il modifié ?
\item \textbf{Prolongement — auto-influence de la bobine.} En assimilant la bobine mobile à une bobine plate de $N$ spires de rayon $R = r$, calcule le champ magnétique $B_{\text{bobine}}$ qu'elle crée en son centre lorsque le courant est maximal ($I_{\max} = \SI{0,8}{A}$). Compare au champ $B$ de l'aimant et conclus sur la validité de l'hypothèse « le champ dans l'entrefer est celui de l'aimant seul ».
\end{enumerate}

\courssub{Ressources à mobiliser}

\begin{aretenir}[Boîte à outils de la leçon]
\begin{itemize}
\item \textbf{Force de Laplace} sur un conducteur de longueur $\ell$ parcouru par un courant $I$ et plongé dans un champ $\vec{B}$ uniforme :
\[ \vec{F} = I\,\vec{\ell} \wedge \vec{B}, \qquad F = I\,\ell\, B \sin\alpha, \]
où $\vec{\ell}$ est orienté dans le sens du courant et $\alpha$ est l'angle entre le conducteur et $\vec{B}$. Si le conducteur est perpendiculaire à $\vec{B}$ : $F = I\ell B$.
\item \textbf{Règle de la main droite} (ou des trois doigts) : les trois vecteurs $I\vec{\ell}$, $\vec{B}$, $\vec{F}$ forment un trièdre direct.
\item \textbf{Champ au centre d'une bobine plate} de $N$ spires de rayon $R$ parcourues par un courant $I$ :
\[ B = \dfrac{\mu_0 N I}{2R}. \]
\item \textbf{Courant sinusoïdal} $i(t) = I_{\max}\cos(2\pi f t)$ : $f$ est la fréquence en hertz ; la force de Laplace, proportionnelle à $i$, oscille à la même fréquence.
\item Analyse dimensionnelle : $[F] = \mathrm{A \cdot m \cdot T} = \mathrm{N}$ car $\SI{1}{T} = \SI{1}{N\,A^{-1}\,m^{-1}}$.
\end{itemize}
\end{aretenir}

\courssub{Démarche et corrigé commenté}

\begin{exoresolu}[Conception du mini-haut-parleur : corrigé complet]
\textbf{Question 1.} Schéma en coupe et sens de la force de Laplace.

\textbf{Question 2.} Longueur totale de fil dans le champ.

\textbf{Question 3.} Force maximale sur la bobine.

\textbf{Question 4.} Fréquence de vibration de la membrane.

\textbf{Question 5.} Inversion des bornes de l'alimentation.

\textbf{Question 6.} Champ créé par la bobine et comparaison.
\tcblower
\textbf{1. Schéma en coupe.} Dans l'entrefer, le champ $\vec{B}$ est radial : il pointe du centre de l'aimant vers l'extérieur. La spire, vue en coupe, apparaît comme deux sections de fil : le courant entre ($\otimes$) d'un côté et sort ($\odot$) de l'autre. Appliquons la règle de la main droite à la section de gauche : $I\vec{\ell}$ est dirigé vers l'arrière de la feuille (courant entrant), $\vec{B}$ pointe vers la gauche (vers l'extérieur de l'anneau) ; le produit vectoriel $\vec{F} = I\vec{\ell} \wedge \vec{B}$ donne une force dirigée \emph{vers la membrane}. Sur la section de droite, le courant est sortant ($\odot$) et $\vec{B}$ pointe vers la droite : le même produit vectoriel donne une force également dirigée vers la membrane. \textbf{Toutes les spires sont poussées dans le même sens} : les forces s'ajoutent, la bobine et la membrane avancent lorsque $i > 0$.

\begin{popfigure}
\begin{center}
\begin{tikzpicture}[scale=1.0,>=Stealth]
% Aimant (coupe)
\fill[poppurple!25] (-3.4,-1.4) rectangle (-2.2,1.4);
\fill[poppurple!25] (2.2,-1.4) rectangle (3.4,1.4);
\node at (-2.8,1.7) {\small Aimant (N)};
\node at (2.8,1.7) {\small Aimant (S)};
% Entrefer : champ radial
\draw[->,popblue,very thick] (-2.2,0.6) -- (-1.2,0.6);
\draw[->,popblue,very thick] (2.2,0.6) -- (1.2,0.6);
\node[popblue] at (0,0.95) {\small $\vec{B}$ radial};
% Axe
\draw[dashed,popmuted] (-3.6,-0.9) -- (3.6,-0.9) node[right]{\small axe};
% Sections de la spire
\draw[fill=popcream,draw=popdark,thick] (-1.0,0.6) circle (0.16);
\node at (-1.0,0.6) {\scriptsize $\otimes$};
\draw[fill=popcream,draw=popdark,thick] (1.0,0.6) circle (0.16);
\node at (1.0,0.6) {\scriptsize $\odot$};
\node[popdark] at (0,0.25) {\small spire de la bobine (coupe)};
% Forces de Laplace
\draw[->,poporange,very thick] (-1.0,0.6) -- (-1.0,-1.5) node[below]{\small $\vec{F}$};
\draw[->,poporange,very thick] (1.0,0.6) -- (1.0,-1.5) node[below]{\small $\vec{F}$};
% Membrane
\draw[popgreen,very thick] (-2.6,-2.3) -- (2.6,-2.3);
\node[popgreen] at (0,-2.65) {\small membrane (cône)};
\end{tikzpicture}
\end{center}
\legende{Coupe du haut-parleur à l'instant où $i>0$ : champ radial $\vec{B}$, courant dans la spire et forces de Laplace $\vec{F}$ poussant la membrane.}
\end{popfigure}

\textbf{2. Longueur de fil dans le champ.} Chacune des $N = 80$ spires est une circonférence de rayon $r = \SI{1,5}{cm} = \SI{1,5e-2}{m}$, entièrement plongée dans le champ radial. La longueur totale de fil soumise à la force de Laplace est donc :
\[ \ell = N \times 2\pi r = 80 \times 2\pi \times \num{1,5e-2} \approx \SI{7,54}{m}. \]

\textbf{3. Force maximale.} Le fil est enroulé en cercle autour de l'axe et le champ est radial : en tout point du fil, le courant est \emph{perpendiculaire} à $\vec{B}$ ($\alpha = 90°$, donc $\sin\alpha = 1$). C'est précisément l'intérêt du champ radial : la géométrie garantit l'angle droit partout, quelle que soit la position de la spire. La force totale sur la bobine est donc :
\[ F(t) = i(t)\,\ell\, B = 0{,}8\cos(2\pi \times 440\, t) \times 7{,}54 \times 0{,}50. \]
Vérification dimensionnelle : $\mathrm{A \times m \times T} = \mathrm{A \times m \times N\,A^{-1}\,m^{-1}} = \mathrm{N}$. La valeur maximale est atteinte quand $\cos(2\pi \times 440\, t) = \pm 1$ :
\[ F_{\max} = 0{,}8 \times 7{,}54 \times 0{,}50 \approx \SI{3,0}{N}. \]
C'est l'ordre de grandeur du poids d'une masse de $\SI{300}{g}$ : une force tout à fait suffisante pour faire vibrer une membrane légère en carton.

\textbf{4. Fréquence de vibration.} La force $F(t) = F_{\max}\cos(2\pi \times 440\, t)$ est sinusoïdale de fréquence $f = \SI{440}{Hz}$. La membrane, entraînée par cette force périodique, vibre à la \emph{même fréquence} et impose à l'air une onde sonore de fréquence $\SI{440}{Hz}$ : c'est précisément la note « la » du diapason. Le haut-parleur remplit sa mission.

\textbf{5. Inversion des bornes.} Inverser les bornes revient à changer le sens du courant : $i(t)$ devient $-i(t)$, c'est-à-dire un déphasage de $\pi$. La force de Laplace change de sens à chaque instant : la membrane recule quand elle avançait auparavant. L'onde sonore émise est simplement \emph{en opposition de phase} avec la précédente : pour un auditeur seul, le son perçu (hauteur $\SI{440}{Hz}$, timbre) est \textbf{inchangé}. En revanche, si deux haut-parleurs fonctionnent ensemble avec des branchements opposés, leurs ondes peuvent se compenser (annulations par interférences) : c'est pourquoi les techniciens respectent la « polarité » des haut-parleurs dans les sonorisations, par exemple lors des concerts au palais des sports de Yaoundé.

\textbf{6. Champ propre de la bobine.} En assimilant la bobine à une bobine plate de $N = 80$ spires de rayon $R = \SI{1,5e-2}{m}$, parcourue par $I_{\max} = \SI{0,8}{A}$ :
\[ B_{\text{bobine}} = \dfrac{\mu_0 N I_{\max}}{2R} = \dfrac{4\pi \times 10^{-7} \times 80 \times 0{,}8}{2 \times \num{1,5e-2}} \approx \SI{2,7e-3}{T}. \]
Comparaison :
\[ \dfrac{B_{\text{bobine}}}{B} = \dfrac{\num{2,7e-3}}{0{,}50} \approx \num{5,4e-3} \approx 0{,}5\,\%. \]
Le champ créé par la bobine est environ \textbf{200 fois plus faible} que celui de l'aimant. L'hypothèse « le champ dans l'entrefer est celui de l'aimant seul » est donc excellente : le calcul de la force de Laplace mené à la question 3 est valide.
\end{exoresolu}

\begin{attention}[Pièges fréquents dans ce type d'exercice]
\begin{itemize}
\item \textbf{Oublier le nombre de spires} : la force de Laplace s'exerce sur \emph{chaque} spire ; il faut multiplier par $N$ (ou utiliser $\ell = 2\pi r N$). Écrire $F = B I \times 2\pi r$ sans le facteur $N$ est l'erreur la plus classique.
\item \textbf{Confondre les règles d'orientation} : la règle de la main droite pour la force de Laplace porte sur le trièdre $(I\vec{\ell}, \vec{B}, \vec{F})$ ; ne pas la confondre avec la règle du tire-bouchon (ou du bonhomme d'Ampère) qui donne le sens du champ créé par un courant.
\item \textbf{Mélanger force de Lorentz et force de Laplace} : la force de Lorentz $\vec{F} = q\,\vec{v}\wedge\vec{B}$ s'applique à une \emph{charge} en mouvement ; la force de Laplace $\vec{F} = I\vec{\ell}\wedge\vec{B}$ s'applique à un \emph{conducteur} parcouru par un courant. Ici, c'est bien Laplace — d'ailleurs, la force de Laplace n'est rien d'autre que la somme des forces de Lorentz subies par les électrons de conduction du fil.
\item \textbf{Négliger la géométrie} : ici $\sin\alpha = 1$ grâce au champ radial ; si le champ était axial (parallèle au déplacement), la force serait nulle ! Toujours vérifier l'angle entre le courant et $\vec{B}$ avant d'écrire $F = I\ell B$.
\end{itemize}
\end{attention}

\courssub{Bilan de l'activité}

Cette activité t'a permis de mettre en évidence les points essentiels de la leçon dans un dispositif technologique réel :

\begin{itemize}
\item un \cle{aimant} crée dans son entrefer un champ magnétique dont on sait représenter la direction et le sens (ici, un champ radial) ;
\item la \cle{force de Laplace} $\vec{F} = I\vec{\ell}\wedge\vec{B}$ est le moteur du haut-parleur : proportionnelle au courant, elle hérite de sa forme temporelle — un courant sinusoïdal à $\SI{440}{Hz}$ produit une vibration sonore à $\SI{440}{Hz}$ ;
\item la \cle{règle de la main droite} permet de prévoir le sens de déplacement de la membrane et de comprendre l'effet d'une inversion de polarité ;
\item le calcul du champ d'une \cle{bobine plate} a montré que le champ propre de la bobine ($\approx \SI{2,7}{mT}$) est négligeable devant celui de l'aimant ($\SI{0,50}{T}$), ce qui valide le modèle ;
\item enfin, le dimensionnement ($F_{\max} \approx \SI{3,0}{N}$) illustre la démarche de l'ingénieur : partir des lois physiques pour prévoir quantitativement le comportement d'un objet technique, comme le font chaque jour les techniciens qui installent et réparent les systèmes audio dans tout le Cameroun.
\end{itemize}

Tu es maintenant prêt à aborder l'étude du moteur électrique à courant continu, qui repose exactement sur les mêmes principes : un bobinage parcouru par un courant, placé dans le champ d'un aimant, et des forces de Laplace dont le moment fait tourner le rotor.

\newpage

\coursec{Méthodes et exercices résolus}

\courssub{Méthode 1 : Représenter le vecteur champ magnétique en un point}

\begin{methode}[Déterminer et représenter $\vec{B}$ en un point]
\begin{enumerate}
\item \textbf{Identifier les sources} de champ magnétique : aimant(s), courant(s), champ terrestre $\vec{B}_T$.
\item \textbf{Rappeler la règle d'orientation} : le vecteur $\vec{B}$ en un point $M$ est tangent à la ligne de champ passant par $M$, orienté du \cle{pôle Nord} vers le \cle{pôle Sud} (à l'extérieur de l'aimant).
\item \textbf{Utiliser l'aiguille aimantée} : une petite boussole placée en $M$ s'oriente selon $\vec{B}$ ; son pôle Nord indique le sens de $\vec{B}$.
\item \textbf{Si plusieurs sources} agissent : appliquer le \cle{principe de superposition} $\vec{B} = \vec{B}_1 + \vec{B}_2$ (somme vectorielle, construction du parallélogramme).
\item \textbf{Calculer la norme} si nécessaire : $B = \sqrt{B_1^2 + B_2^2}$ si les champs sont perpendiculaires ; vérifier l'unité : le tesla (T).
\item \textbf{Représenter} le vecteur à l'échelle (par exemple $1\ \mathrm{cm} \leftrightarrow 10^{-4}\ \mathrm{T}$), en précisant l'échelle sur la copie.
\end{enumerate}
\textbf{Réflexe Bac} : toujours indiquer sur le schéma la direction, le sens et l'échelle ; un vecteur sans échelle ni légende coûte des points.
\end{methode}

\begin{exoresolu}[Superposition du champ d'un aimant et du champ terrestre]
Une petite aiguille aimantée est placée en un point $M$. En l'absence d'aimant, elle s'oriente selon le champ magnétique terrestre $\vec{B}_T$ de valeur $B_T = \num{4,0e-5}\ \mathrm{T}$, dirigé vers le Nord géographique. On approche un aimant droit dont le pôle Sud fait face à $M$, situé à l'Est de $M$ ; l'aimant crée en $M$ un champ $\vec{B}_A$ de valeur $B_A = \num{3,0e-5}\ \mathrm{T}$, dirigé vers l'Est. Déterminer les caractéristiques du champ résultant $\vec{B}$ en $M$ et l'angle $\alpha$ dont tourne l'aiguille par rapport au Nord.
\tcblower
\textbf{1. Schéma et données.} $\vec{B}_T$ pointe vers le Nord, $\vec{B}_A$ pointe vers l'Est : les deux vecteurs sont \textbf{perpendiculaires}. Le champ résultant est, par superposition :
\[ \vec{B} = \vec{B}_T + \vec{B}_A. \]
Remarque de cohérence : le pôle Sud de l'aimant fait face à $M$, donc les lignes de champ de l'aimant \emph{rentrent} par ce pôle Sud ; en $M$, situé à l'Ouest de l'aimant, $\vec{B}_A$ pointe bien vers l'Est, c'est-à-dire vers l'aimant.

\textbf{2. Norme du champ résultant.} Le triangle des vecteurs est rectangle en $M$, donc d'après le théorème de Pythagore :
\[ B = \sqrt{B_T^2 + B_A^2} = \sqrt{(4,0\times 10^{-5})^2 + (3,0\times 10^{-5})^2} = 10^{-5}\sqrt{16 + 9} = \num{5,0e-5}\ \mathrm{T}. \]

\textbf{3. Direction.} L'angle $\alpha$ entre $\vec{B}$ et $\vec{B}_T$ (le Nord) vérifie :
\[ \tan\alpha = \dfrac{B_A}{B_T} = \dfrac{3,0\times 10^{-5}}{4,0\times 10^{-5}} = 0,75 \quad\Longrightarrow\quad \alpha \approx 37^\circ. \]

\textbf{4. Conclusion.} Le champ résultant en $M$ a pour norme $B = \num{5,0e-5}\ \mathrm{T}$ ; l'aiguille aimantée tourne d'environ $37^\circ$ vers l'Est par rapport à sa direction initiale Nord. On représente $\vec{B}$ à l'échelle $1\ \mathrm{cm} \leftrightarrow 10^{-5}\ \mathrm{T}$ : $\vec{B}_T$ mesure $4\ \mathrm{cm}$, $\vec{B}_A$ $3\ \mathrm{cm}$, $\vec{B}$ $5\ \mathrm{cm}$.
\end{exoresolu}

\courssub{Méthode 2 : Représenter les lignes d'un champ magnétique créé par un courant}

\begin{methode}[Tracer les lignes de champ d'un conducteur parcouru par un courant]
\begin{enumerate}
\item \textbf{Identifier la géométrie} du conducteur : fil rectiligne, spire circulaire, bobine plate ou solénoïde.
\item \textbf{Appliquer la règle du bonhomme d'Ampère} (ou du tire-bouchon) : pour un fil rectiligne, l'observateur est traversé par le courant qui entre par ses pieds et sort par sa tête ; regardant vers le point $M$, son bras gauche tendu indique le sens de $\vec{B}$. Les lignes de champ sont des \cle{cercles concentriques} centrés sur le fil, dans des plans perpendiculaires au fil.
\item \textbf{Pour une spire ou un solénoïde} : la face où le courant tourne dans le sens trigonométrique (vue par l'observateur) est une \cle{face Nord} (le champ en sort) ; l'autre est une face Sud. À l'intérieur d'un solénoïde long, le champ est \cle{quasi uniforme} : lignes parallèles à l'axe et équidistantes.
\item \textbf{Orienter les lignes} : à l'extérieur, elles sortent par le Nord et rentrent par le Sud ; elles sont \textbf{fermées} (contrairement aux lignes de champ électrostatique).
\item \textbf{Écrire la norme} si demandée :
\begin{itemize}
\item fil rectiligne long : $B = \dfrac{\mu_0 I}{2\pi d}$ ($d$ : distance du fil au point $M$) ;
\item solénoïde long : $B = \mu_0 n I$ avec $n = \dfrac{N}{L}$ le nombre de spires par mètre.
\end{itemize}
\item \textbf{Vérifier l'homogénéité} : $\mu_0$ en $\mathrm{T\,m\,A^{-1}}$, $I$ en A, $d$ et $L$ en m, $n$ en $\mathrm{m^{-1}}$ : $B$ est bien en tesla.
\end{enumerate}
\end{methode}

\begin{exoresolu}[Champ créé par un solénoïde — alimentation d'une bobine de laboratoire]
Au laboratoire du lycée, on réalise un solénoïde en enroulant $N = 800$ spires de fil de cuivre sur un cylindre de longueur $L = 40\ \mathrm{cm}$. On y fait circuler un courant d'intensité $I = 2,5\ \mathrm{A}$. On donne $\mu_0 = 4\pi \times 10^{-7}\ \mathrm{T\,m\,A^{-1}}$.
\begin{enumerate}
\item Représenter les lignes de champ à l'intérieur du solénoïde et préciser la nature du champ.
\item Calculer la valeur $B$ du champ magnétique à l'intérieur du solénoïde.
\item Comment faut-il modifier $I$ pour doubler $B$ ?
\end{enumerate}
\tcblower
\textbf{1. Lignes de champ.} Le solénoïde est long ($L$ grand devant son diamètre) : à l'intérieur, les lignes de champ sont des \textbf{droites parallèles à l'axe}, équidistantes, orientées de la face Sud vers la face Nord (à l'intérieur). Le champ est donc \textbf{uniforme} à l'intérieur du solénoïde. La face où le courant circule dans le sens trigonométrique (vue par un observateur) est la face Nord.

\textbf{2. Calcul de $B$.} Le nombre de spires par mètre est :
\[ n = \dfrac{N}{L} = \dfrac{800}{0,40} = 2{,}0\times 10^{3}\ \text{spires/m}. \]
D'où :
\[ B = \mu_0 n I = 4\pi \times 10^{-7} \times 2{,}0\times 10^{3} \times 2,5 = 4\pi \times 10^{-7} \times 5{,}0\times 10^{3} = 2\pi \times 10^{-3}\ \mathrm{T}. \]
\[ \boxed{B \approx \num{6,3e-3}\ \mathrm{T} = 6,3\ \mathrm{mT}} \]
\textit{Vérification dimensionnelle} : $[\mu_0 n I] = \mathrm{T\,m\,A^{-1}} \times \mathrm{m^{-1}} \times \mathrm{A} = \mathrm{T}$. Correct.

\textbf{3. Doubler $B$.} Comme $B = \mu_0 n I$ est proportionnelle à $I$ ($n$ fixé), il faut doubler l'intensité : $I' = 2I = 5,0\ \mathrm{A}$.

\textbf{Conclusion} : le champ à l'intérieur du solénoïde est uniforme, de valeur $B \approx 6,3\ \mathrm{mT}$, soit environ $150$ fois le champ terrestre ($B_T \approx 4\times 10^{-5}\ \mathrm{T}$) ; il double si l'intensité double.
\end{exoresolu}

\courssub{Méthode 3 : Exprimer la force de Lorentz sur une charge en mouvement}

\begin{methode}[Caractéristiques de la force de Lorentz $\vec{F} = q\,\vec{v} \wedge \vec{B}$]
\begin{enumerate}
\item \textbf{Vérifier les conditions d'application} : particule de charge $q$ (algébrique, en C), de vitesse $\vec{v}$, placée dans un champ $\vec{B}$ ; la force est nulle si $\vec{v} = \vec{0}$ ou si $\vec{v} \parallel \vec{B}$.
\item \textbf{Direction} : $\vec{F}$ est \textbf{perpendiculaire au plan} $(\vec{v}, \vec{B})$ — toujours le dire dans la rédaction.
\item \textbf{Sens} : règle de la main droite (ou du tire-bouchon) appliquée à $\vec{v} \wedge \vec{B}$, puis \textbf{inverser le sens si $q < 0$} (électron !).
\item \textbf{Norme} : $F = |q|\,v\,B\,|\sin\alpha|$ où $\alpha = (\vec{v}, \vec{B})$. Cas fréquent au Bac : $\vec{v} \perp \vec{B}$, alors $F = |q|\,v\,B$.
\item \textbf{Conséquence énergétique} : $\vec{F} \perp \vec{v}$ donc la puissance $\vec{F}\cdot\vec{v} = 0$ : la force de Lorentz \textbf{ne travaille pas} et ne modifie pas la norme de la vitesse, seulement sa direction (mouvement circulaire uniforme si $\vec{v} \perp \vec{B}$, de rayon $R = \dfrac{mv}{|q|B}$).
\item \textbf{Unités} : $q$ en C, $v$ en $\mathrm{m\,s^{-1}}$, $B$ en T, $F$ en N.
\end{enumerate}
\end{methode}

\begin{exoresolu}[Électron dans un champ magnétique uniforme — tube cathodique]
Un électron (charge $q = -e = \num{-1,6e-19}\ \mathrm{C}$, masse $m = \num{9,1e-31}\ \mathrm{kg}$) pénètre avec une vitesse $v = \num{2,0e7}\ \mathrm{m\,s^{-1}}$ dans une zone où règne un champ magnétique uniforme $B = \num{1,0e-3}\ \mathrm{T}$, perpendiculaire à $\vec{v}$. On néglige le poids de l'électron.
\begin{enumerate}
\item Donner les caractéristiques de la force de Lorentz subie par l'électron à l'entrée de la zone.
\item Justifier que le mouvement est circulaire uniforme et calculer le rayon $R$ de la trajectoire.
\item Vérifier que le poids est bien négligeable devant la force de Lorentz.
\end{enumerate}
\tcblower
\textbf{1. Caractéristiques de la force de Lorentz.}
\begin{itemize}
\item \textbf{Point d'application} : l'électron (particule ponctuelle).
\item \textbf{Direction} : perpendiculaire au plan $(\vec{v}, \vec{B})$.
\item \textbf{Sens} : celui de $\vec{v} \wedge \vec{B}$ \textbf{renversé} car $q = -e < 0$.
\item \textbf{Norme} : comme $\vec{v} \perp \vec{B}$ ($\sin\alpha = 1$) :
\[ F = |q|\,v\,B = 1,6\times 10^{-19} \times 2,0\times 10^{7} \times 1,0\times 10^{-3} = \num{3,2e-15}\ \mathrm{N}. \]
\end{itemize}

\textbf{2. Nature du mouvement.} La force de Lorentz est à chaque instant perpendiculaire à la vitesse : sa puissance est nulle, donc d'après le théorème de l'énergie cinétique, $E_c$ et donc $v$ sont \textbf{constantes} : le mouvement est \textbf{uniforme}. De plus $\vec{F}$ est centripète de norme constante : la trajectoire est un \textbf{cercle}. En appliquant la deuxième loi de Newton selon l'axe normal : $evB = m\dfrac{v^2}{R}$, d'où :
\[ R = \dfrac{mv}{eB} = \dfrac{9,1\times 10^{-31} \times 2,0\times 10^{7}}{1,6\times 10^{-19} \times 1,0\times 10^{-3}} = \dfrac{1,82\times 10^{-23}}{1,6\times 10^{-22}} \approx \num{0,11}\ \mathrm{m}. \]
\[ \boxed{R \approx 11\ \mathrm{cm}} \]

\textbf{3. Comparaison avec le poids.}
\[ P = mg = 9,1\times 10^{-31} \times 9,8 \approx \num{8,9e-30}\ \mathrm{N} \quad\Longrightarrow\quad \dfrac{F}{P} \approx \dfrac{3,2\times 10^{-15}}{8,9\times 10^{-30}} \approx 3,6\times 10^{14}. \]
La force de Lorentz est environ $10^{14}$ fois plus grande que le poids : celui-ci est tout à fait négligeable, l'hypothèse est justifiée.

\textbf{Conclusion} : l'électron décrit un cercle de rayon $R \approx 11\ \mathrm{cm}$ à vitesse constante $v = \num{2,0e7}\ \mathrm{m\,s^{-1}}$ ; c'est le principe de la déviation magnétique utilisée dans les anciens téléviseurs à tube cathodique.
\end{exoresolu}

\courssub{Méthode 4 : Exprimer la force de Laplace sur un conducteur}

\begin{methode}[Caractéristiques de la force de Laplace $\vec{F} = I\,\vec{\ell} \wedge \vec{B}$]
\begin{enumerate}
\item \textbf{Conditions} : portion de conducteur \textbf{rectiligne} de longueur $\ell$, parcourue par un courant $I$, placée dans un champ $\vec{B}$ uniforme ; $\vec{\ell}$ est orienté \textbf{dans le sens du courant}.
\item \textbf{Direction} : perpendiculaire au plan formé par le conducteur et $\vec{B}$.
\item \textbf{Sens} : \cle{règle de la main droite} (ou des trois doigts) : le courant $I$ remplace la vitesse des charges positives ; orienter $\vec{\ell} \wedge \vec{B}$.
\item \textbf{Norme} : $F = I\,\ell\,B\,|\sin\alpha|$ ; si le conducteur est perpendiculaire à $\vec{B}$ : $F = I\ell B$.
\item \textbf{Point d'application} : le \textbf{milieu} de la portion de conducteur (force répartie représentée par une force équivalente).
\item \textbf{Réflexe} : la force de Laplace est la somme des forces de Lorentz sur les porteurs de charge du conducteur — savoir le dire pour justifier le lien entre les deux forces.
\end{enumerate}
\end{methode}

\begin{exoresolu}[Balance de Cotton simplifiée — mesure d'un champ magnétique]
Une tige conductrice $CD$ de longueur $\ell = 10\ \mathrm{cm}$, horizontale et mobile autour d'un axe horizontal, est parcourue par un courant $I = 5,0\ \mathrm{A}$. Sa partie centrale baigne dans un champ magnétique uniforme $\vec{B}$, horizontal et perpendiculaire à la tige. Pour rétablir l'équilibre, on doit placer une masse $m = 2,0\ \mathrm{g}$ du côté opposé (bras de levier identiques). On donne $g = 9,8\ \mathrm{N\,kg^{-1}}$.
\begin{enumerate}
\item Représenter les forces en jeu et préciser le sens de la force de Laplace.
\item Calculer la valeur $B$ du champ magnétique.
\end{enumerate}
\tcblower
\textbf{1. Forces en jeu.} La tige est soumise à :
\begin{itemize}
\item son poids et les réactions de l'axe (moments nuls ou compensés par la géométrie du dispositif) ;
\item la \textbf{force de Laplace} $\vec{F} = I\,\vec{\ell} \wedge \vec{B}$ : le courant et le champ étant horizontaux et perpendiculaires entre eux, $\vec{F}$ est \textbf{verticale} ; d'après la règle de la main droite, avec le sens du courant choisi, elle est dirigée \textbf{vers le bas} (c'est elle que la masse $m$ doit compenser) ;
\item le poids $P = mg$ de la masse ajoutée, vertical vers le bas, de l'autre côté de l'axe.
\end{itemize}

\textbf{2. Calcul de $B$.} À l'équilibre, les moments des deux forces par rapport à l'axe se compensent ; les bras de levier étant égaux, les normes sont égales :
\[ F = P \quad\Longleftrightarrow\quad I\ell B = mg \quad\Longrightarrow\quad B = \dfrac{mg}{I\ell}. \]
Application numérique (avec $m = 2,0\times 10^{-3}\ \mathrm{kg}$ et $\ell = 0,10\ \mathrm{m}$) :
\[ B = \dfrac{2,0\times 10^{-3} \times 9,8}{5,0 \times 0,10} = \dfrac{1,96\times 10^{-2}}{0,50} = \num{3,9e-2}\ \mathrm{T}. \]
\textit{Vérification dimensionnelle} : $\left[\dfrac{mg}{I\ell}\right] = \dfrac{\mathrm{kg\,m\,s^{-2}}}{\mathrm{A\,m}} = \dfrac{\mathrm{N}}{\mathrm{A\,m}} = \mathrm{T}$, car $1\ \mathrm{T} = 1\ \mathrm{N\,A^{-1}\,m^{-1}}$. Correct.

\textbf{Conclusion} : le champ magnétique vaut $B \approx 39\ \mathrm{mT}$ ; la force de Laplace mise en jeu est $F = I\ell B \approx \num{2,0e-2}\ \mathrm{N}$, égale au poids de la masse de $2,0\ \mathrm{g}$. C'est le principe de la balance de Cotton, méthode historique de mesure des champs magnétiques.
\end{exoresolu}

\courssub{Méthode 5 : Analyser le fonctionnement d'un haut-parleur ou d'un moteur}

\begin{methode}[Expliquer un dispositif électromécanique]
\begin{enumerate}
\item \textbf{Repérer les trois ingrédients} : un champ magnétique $\vec{B}$ (aimant permanent), un conducteur mobile (bobine, spire) parcouru par un courant $I$, et une force de Laplace motrice.
\item \textbf{Haut-parleur} : le courant \textbf{variable} $i(t)$ (signal audio) traverse la bobine placée dans le champ radial de l'aimant ; la force de Laplace variable $F = i\,\ell\,B$ fait vibrer la bobine solidaire de la membrane, qui émet le son. La fréquence de vibration est celle du courant.
\item \textbf{Moteur à courant continu} : une spire (ou un ensemble de spires) placée dans $\vec{B}$ et parcourue par $I$ subit un \textbf{couple de forces de Laplace} dont le moment met la spire en rotation ; le \cle{collecteur} inverse le courant à chaque demi-tour pour maintenir le sens de rotation.
\item \textbf{Rédaction type} : nommer la force, donner son expression, expliquer le lien courant $\rightarrow$ force $\rightarrow$ mouvement, citer le rôle de chaque pièce.
\end{enumerate}
\end{methode}

\begin{exoresolu}[Haut-parleur électrodynamique]
La bobine d'un haut-parleur comporte $N = 50$ spires de rayon moyen $r = 2,0\ \mathrm{cm}$. Elle est placée dans un champ magnétique radial de valeur $B = 0,80\ \mathrm{T}$, perpendiculaire en tout point aux spires. Un courant $i(t) = I_{\max}\cos(2\pi f t)$ avec $I_{\max} = 0,50\ \mathrm{A}$ et $f = 440\ \mathrm{Hz}$ la parcourt.
\begin{enumerate}
\item Exprimer puis calculer la valeur maximale $F_{\max}$ de la force de Laplace s'exerçant sur la bobine.
\item Expliquer pourquoi la membrane émet un son et préciser sa hauteur.
\end{enumerate}
\tcblower
\textbf{1. Force maximale.} Chaque spire de circonférence $2\pi r$, parcourue par $i$ et placée dans un champ perpendiculaire au fil en tout point, subit une force axiale de norme $i \times 2\pi r \times B$. Pour $N$ spires, les forces s'ajoutent :
\[ F(t) = N \times 2\pi r \times B \times i(t). \]
La valeur maximale est :
\[ F_{\max} = 2\pi N r B I_{\max} = 2\pi \times 50 \times 2,0\times 10^{-2} \times 0,80 \times 0,50. \]
\[ F_{\max} = 2\pi \times 50 \times 8,0\times 10^{-3} = 2\pi \times 0,40 \approx \num{2,5}\ \mathrm{N}. \]
\textit{Vérification dimensionnelle} : $[N r B I] = \mathrm{m \times T \times A} = \mathrm{m \times N\,A^{-1}\,m^{-1} \times A} = \mathrm{N}$. Correct.

\textbf{2. Émission du son.} La force $F(t) = F_{\max}\cos(2\pi f t)$ est \textbf{alternative} : elle pousse puis tire la bobine $440$ fois par seconde. La bobine, solidaire de la membrane, la fait vibrer à la même fréquence ; la membrane comprime successivement l'air, créant une onde sonore de fréquence $f = 440\ \mathrm{Hz}$ : c'est la note \emph{la} (la$_3$). L'amplitude du son dépend de $F_{\max}$, donc de $I_{\max}$ : plus le courant est fort, plus le son est intense.

\textbf{Conclusion} : la force de Laplace convertit le signal électrique en vibration mécanique : $F_{\max} \approx 2,5\ \mathrm{N}$ et le son émis a une fréquence de $440\ \mathrm{Hz}$.
\end{exoresolu}

\courssub{Méthode 6 : Interaction entre deux conducteurs parallèles et définition de l'ampère}

\begin{methode}[Force entre deux fils parallèles]
\begin{enumerate}
\item \textbf{Raisonner en deux étapes} : le fil 1 crée en tout point du fil 2 un champ $B_1 = \dfrac{\mu_0 I_1}{2\pi d}$ ; le fil 2, parcouru par $I_2$ dans ce champ, subit une force de Laplace.
\item \textbf{Norme de la force} sur une longueur $\ell$ :
\[ F = I_2\,\ell\,B_1 = \dfrac{\mu_0\,I_1 I_2\,\ell}{2\pi d}. \]
\item \textbf{Nature de l'interaction} (règle de la main droite) : courants de \textbf{même sens} $\Rightarrow$ \textbf{attraction} ; courants de \textbf{sens opposés} $\Rightarrow$ \textbf{répulsion}. À savoir démontrer sur schéma.
\item \textbf{Définition de l'ampère} (définition historique issue de cette expérience) : l'ampère est l'intensité d'un courant constant qui, maintenu dans deux conducteurs rectilignes, parallèles, de longueur infinie, de section négligeable et distants de $1\ \mathrm{m}$ dans le vide, produit entre eux une force de $2\times 10^{-7}\ \mathrm{N}$ par mètre de longueur.
\item \textbf{Vérifier} : $\dfrac{F}{\ell} = \dfrac{\mu_0 I_1 I_2}{2\pi d}$ ; avec $I_1 = I_2 = 1\ \mathrm{A}$, $d = 1\ \mathrm{m}$ : $\dfrac{F}{\ell} = \dfrac{4\pi\times 10^{-7}}{2\pi} = 2\times 10^{-7}\ \mathrm{N/m}$.
\end{enumerate}
\end{methode}

\begin{exoresolu}[Deux lignes électriques parallèles]
Deux câbles rectilignes parallèles d'une ligne de transport d'électricité, distants de $d = 50\ \mathrm{cm}$, sont parcourus par des courants continus de même sens $I_1 = I_2 = 200\ \mathrm{A}$ (modèle simplifié). On donne $\mu_0 = 4\pi\times 10^{-7}\ \mathrm{T\,m\,A^{-1}}$.
\begin{enumerate}
\item Calculer le champ magnétique $B_1$ créé par le câble 1 à l'emplacement du câble 2.
\item En déduire la force électromagnétique par mètre de câble ; préciser sa nature (attractive ou répulsive).
\end{enumerate}
\tcblower
\textbf{1. Champ créé par le câble 1.} Pour un fil rectiligne long, avec $d = 0,50\ \mathrm{m}$ :
\[ B_1 = \dfrac{\mu_0 I_1}{2\pi d} = \dfrac{4\pi\times 10^{-7} \times 200}{2\pi \times 0,50} = \dfrac{8\pi\times 10^{-5}}{\pi} = 8,0\times 10^{-5}\ \mathrm{T}. \]
\[ \boxed{B_1 = \num{8,0e-5}\ \mathrm{T}} \]
Ce champ est perpendiculaire au plan contenant les deux fils (règle du bonhomme d'Ampère), soit environ deux fois le champ terrestre.

\textbf{2. Force par mètre.} Le câble 2 baigne dans $B_1$ qui lui est perpendiculaire ; la force de Laplace sur une longueur $\ell = 1\ \mathrm{m}$ vaut :
\[ F = I_2\,\ell\,B_1 = 200 \times 1 \times 8,0\times 10^{-5} = \num{1,6e-2}\ \mathrm{N}. \]
On retrouve la formule générale : $\dfrac{F}{\ell} = \dfrac{\mu_0 I_1 I_2}{2\pi d} = \dfrac{4\pi\times 10^{-7}\times 200\times 200}{2\pi\times 0,50} = 1,6\times 10^{-2}\ \mathrm{N/m}$.

\textbf{Nature de la force} : les courants sont de \textbf{même sens} ; la règle de la main droite appliquée à $I_2\vec{\ell}\wedge\vec{B}_1$ montre que $\vec{F}$ est dirigée \textbf{vers le câble 1} : l'interaction est \textbf{attractive}.

\textbf{Conclusion} : chaque mètre de câble subit une force attractive de $16\ \mathrm{mN}$. Faible ici, cette force devient considérable lors des courts-circuits (courants de plusieurs kiloampères, et $F$ proportionnelle au \emph{produit} des courants) : les câbles doivent être solidement ancrés, et c'est précisément cette interaction qui a servi à définir l'ampère.
\end{exoresolu}

\begin{attention}[Les 5 erreurs qui coûtent des points au Bac]
\begin{enumerate}
\item \textbf{Oublier le signe de la charge dans la force de Lorentz.} Pour un électron ($q = -e < 0$), le sens de $\vec{F}$ est \textbf{opposé} à celui de $\vec{v}\wedge\vec{B}$. C'est l'erreur n°1 : annonce explicitement « $q<0$, donc on inverse le sens ».
\item \textbf{Confondre les deux forces.} Force de \textbf{Lorentz} : sur une \emph{charge en mouvement}, $F = |q|vB\sin\alpha$ ; force de \textbf{Laplace} : sur un \emph{conducteur parcouru par un courant}, $F = I\ell B\sin\alpha$. Écrire $F = qvB$ pour un fil conducteur (ou l'inverse) est une faute de cours éliminatoire.
\item \textbf{Oublier le facteur $\sin\alpha$ ou le conserver à tort.} Si $\vec{v}$ (ou le courant) est parallèle à $\vec{B}$, la force est \textbf{nulle} ; si l'énoncé précise « perpendiculaire », écris $\sin\alpha = 1$ et simplifie. Ne jamais écrire $F = qvB$ sans avoir vérifié l'angle.
\item \textbf{Erreurs de conversion et de puissances de 10.} Convertis systématiquement : longueurs en mètres ($10\ \mathrm{cm} = 0,10\ \mathrm{m}$), masses en kg ($2\ \mathrm{g} = 2\times 10^{-3}\ \mathrm{kg}$), et pour le solénoïde $n = N/L$ en spires \textbf{par mètre}. Présente chaque résultat avec son unité SI et 2 ou 3 chiffres significatifs.
\item \textbf{Affirmer que la force de Lorentz accélère la particule.} $\vec{F}\perp\vec{v}$ donc le travail est nul : la \textbf{norme} de la vitesse ne change pas, seule sa direction change (mouvement circulaire \emph{uniforme}). Écrire « la particule accélère » ou « le mouvement est uniformément accéléré » est faux ; de même, ne jamais oublier de justifier que le poids est négligeable par un calcul de rapport $F/P$.
\end{enumerate}
\end{attention}

\newpage

\coursec{Exercices}

\courssub{Je vérifie mes connaissances}

\begin{exercice}[QCM : notions fondamentales]
Pour chaque affirmation, choisir la (ou les) bonne(s) réponse(s) en justifiant brièvement.

\textbf{1.} Le vecteur champ magnétique $\vec{B}$ créé en un point $M$ par un fil rectiligne infini parcouru par un courant $I$ est :
\begin{itemize}
\item[a)] dirigé selon le fil ;
\item[b)] tangent au cercle passant par $M$, centré sur le fil et situé dans un plan perpendiculaire au fil ;
\item[c)] dirigé vers le fil.
\end{itemize}

\textbf{2.} L'intensité du champ magnétique à l'intérieur d'un solénoïde long comportant $n$ spires par mètre, parcouru par un courant $I$, vaut :
\begin{itemize}
\item[a)] $B = \mu_0\, n\, I$ ;
\item[b)] $B = \dfrac{\mu_0 I}{2\pi n}$ ;
\item[c)] $B = \mu_0\, n\, I^2$.
\end{itemize}

\textbf{3.} La force de Lorentz s'exerçant sur une charge $q$ animée d'une vitesse $\vec{v}$ dans un champ $\vec{B}$ :
\begin{itemize}
\item[a)] est toujours parallèle à $\vec{v}$ ;
\item[b)] est perpendiculaire à $\vec{v}$ et à $\vec{B}$ ;
\item[c)] ne travaille pas.
\end{itemize}

\textbf{4.} Deux conducteurs parallèles parcourus par des courants de sens opposés :
\begin{itemize}
\item[a)] s'attirent ;
\item[b)] se repoussent ;
\item[c)] n'interagissent pas.
\end{itemize}
\end{exercice}

\begin{exercice}[Vrai ou faux, à justifier]
Répondre par vrai ou faux en justifiant chaque réponse par une propriété du cours.

\textbf{1.} Une ligne de champ magnétique est orientée dans le sens sud $\rightarrow$ nord à l'extérieur d'un aimant.

\textbf{2.} Une charge électrique immobile placée dans un champ magnétique uniforme subit une force de Lorentz non nulle.

\textbf{3.} La force de Laplace s'exerçant sur un conducteur rectiligne de longueur $\ell$ parcouru par un courant $I$ et placé dans un champ magnétique uniforme $\vec{B}$ a pour expression vectorielle $\vec{F} = I\,\vec{\ell} \wedge \vec{B}$.

\textbf{4.} La face d'une bobine plate d'où le courant semble tourner dans le sens trigonométrique est une face nord.

\textbf{5.} Le champ magnétique terrestre est uniforme à la surface du globe.
\end{exercice}

\begin{exercice}[Définitions et représentations]
\textbf{1.} Définir : vecteur champ magnétique ; ligne de champ ; spectre magnétique.

\textbf{2.} Énoncer la règle qui permet de déterminer le sens du champ magnétique créé par un fil rectiligne parcouru par un courant.

\textbf{3.} Sur un schéma en coupe, représenter un fil rectiligne horizontal parcouru par un courant $I$ dirigé vers la droite, puis les lignes de champ et le vecteur $\vec{B}$ en un point $M$ situé au-dessus du fil.

\textbf{4.} Donner l'unité SI du champ magnétique et celle de la perméabilité du vide $\mu_0$. Vérifier par analyse dimensionnelle que $\dfrac{\mu_0 I}{2\pi d}$ est homogène à un champ magnétique.
\end{exercice}

\begin{exercice}[Calculs directs]
\textbf{1.} Un fil rectiligne est parcouru par un courant $I = \SI{5,0}{A}$. Calculer l'intensité du champ magnétique créé à la distance $d = \SI{2,0}{cm}$ du fil. On donne $\mu_0 = 4\pi \times 10^{-7}\ \mathrm{S.I.}$

\textbf{2.} Un proton ($q = e = \num{1,6e-19}$ C) se déplace à la vitesse $v = \SI{1,0e6}{m.s^{-1}}$ perpendiculairement à un champ magnétique uniforme $B = \SI{0,50}{T}$. Calculer l'intensité de la force de Lorentz qu'il subit.

\textbf{3.} Un conducteur rectiligne de longueur $\ell = \SI{10}{cm}$, parcouru par un courant $I = \SI{4,0}{A}$, est placé perpendiculairement à un champ magnétique uniforme $B = \SI{0,20}{T}$. Calculer l'intensité de la force de Laplace.

\textbf{4.} Reprendre la question 3 avec le conducteur faisant un angle $\alpha = 30^{\circ}$ avec $\vec{B}$.
\end{exercice}

\courssub{Je m'entraîne}

\begin{exercice}[Solénoïde d'un laboratoire de lycée]
Un solénoïde comporte $n = 800$ spires par mètre et est parcouru par un courant d'intensité $I$ réglable. On donne $\mu_0 = 4\pi \times 10^{-7}\ \mathrm{S.I.}$

\textbf{1.} Rappeler l'expression de l'intensité du champ magnétique à l'intérieur du solénoïde, en précisant les conditions de validité.

\textbf{2.} Déterminer la valeur de $I$ permettant d'obtenir un champ intérieur $B = \SI{2,5}{mT}$.

\textbf{3.} Représenter le solénoïde en coupe, le sens du courant, les lignes de champ intérieures et extérieures, et identifier les faces nord et sud.

\textbf{4.} Que devient la valeur de $B$ si l'on double à la fois $n$ et $I$ ?
\end{exercice}

\begin{exercice}[Rails de Laplace horizontaux]
Une barre métallique $MN$ de masse $m = \SI{20}{g}$ peut glisser sans frottement sur deux rails horizontaux parallèles distants de $\ell = \SI{12}{cm}$. L'ensemble est plongé dans un champ magnétique uniforme vertical, dirigé vers le haut, d'intensité $B = \SI{0,40}{T}$. Un générateur fait circuler un courant d'intensité $I$ dans le circuit. On donne $g = \SI{9,8}{N.kg^{-1}}$.

\textbf{1.} Faire un schéma du dispositif en vue de dessus et en coupe, en représentant les vecteurs $\vec{B}$, $I\vec{\ell}$ et la force de Laplace $\vec{F}$.

\textbf{2.} Calculer l'intensité $I$ du courant pour que la barre subisse une force de Laplace d'intensité $F = \SI{0,050}{N}$.

\textbf{3.} La barre étant initialement au repos, déterminer la nature de son mouvement et calculer son accélération.

\textbf{4.} Comment faut-il modifier le sens du courant ou celui de $\vec{B}$ pour inverser le sens du déplacement de la barre ? Justifier à l'aide de la règle de la main droite.
\end{exercice}

\begin{exercice}[Interaction entre deux câbles de transport d'énergie]
Deux câbles électriques parallèles, assimilés à des fils infinis, sont distants de $d = \SI{10}{cm}$. Ils sont parcourus par des courants de même sens d'intensités $I_1 = \SI{20}{A}$ et $I_2 = \SI{30}{A}$. On donne $\mu_0 = 4\pi \times 10^{-7}\ \mathrm{S.I.}$

\textbf{1.} Exprimer puis calculer l'intensité du champ magnétique $\vec{B}_1$ créé par le fil 1 à l'emplacement du fil 2.

\textbf{2.} En déduire l'expression littérale puis la valeur de la force par unité de longueur $\dfrac{F}{\ell}$ s'exerçant sur chaque fil.

\textbf{3.} Préciser, en le justifiant, si cette force est attractive ou répulsive.

\textbf{4.} Représenter sur un schéma en coupe les vecteurs $\vec{B}_1$, $\vec{B}_2$ et les forces $\vec{F}_{1/2}$ et $\vec{F}_{2/1}$.

\textbf{5.} Que devient la force par unité de longueur si les courants sont de sens opposés, les intensités étant inchangées ?
\end{exercice}

\begin{exercice}[Superposition du champ d'un fil et du champ terrestre]
Un fil rectiligne vertical est parcouru par un courant ascendant d'intensité $I = \SI{10}{A}$. On considère un point $M$ situé à $d = \SI{5,0}{cm}$ à l'est du fil, au niveau du sol. La composante horizontale du champ magnétique terrestre en ce lieu vaut $B_h = \num{2,0e-5}$ T et est dirigée vers le nord. On donne $\mu_0 = 4\pi \times 10^{-7}\ \mathrm{S.I.}$

\textbf{1.} Déterminer la direction, le sens et l'intensité du champ $\vec{B}_{\text{fil}}$ créé par le fil au point $M$. On s'aidera d'un schéma vu de dessus (nord en haut de la page).

\textbf{2.} Montrer que $\vec{B}_{\text{fil}}$ et $\vec{B}_h$ sont perpendiculaires. Calculer l'intensité du champ résultant $\vec{B}$ en $M$.

\textbf{3.} Calculer l'angle que fait le champ résultant avec la direction du nord.

\textbf{4.} Une aiguille aimantée mobile autour d'un axe vertical est placée en $M$. Que indique-t-elle ? Ce dispositif rappelle une expérience historique célèbre : laquelle ?
\end{exercice}

\courssub{Je me prépare au Bac}

\begin{exercice}[Type Bac : électron dans un champ magnétique uniforme]
Dans un tube cathodique de démonstration, un électron pénètre en un point $O$ dans une région où règne un champ magnétique uniforme $\vec{B}$, avec une vitesse $\vec{v}_0$ perpendiculaire à $\vec{B}$.

\textbf{Données :} $v_0 = \num{2,0e7}$ m/s ; $B = \SI{1,0}{mT}$ ; charge élémentaire $e = \num{1,6e-19}$ C ; masse de l'électron $m = \num{9,1e-31}$ kg. On néglige le poids de l'électron devant la force magnétique.

\textbf{1.} Représenter sur un schéma les vecteurs $\vec{v}_0$, $\vec{B}$ et la force de Lorentz $\vec{F}$ au point $O$ (on choisira $\vec{B}$ sortant du plan de la figure). Donner les caractéristiques (direction, sens, norme) de $\vec{F}$ et calculer son intensité.

\textbf{2.} Justifier que le poids de l'électron est effectivement négligeable devant la force de Lorentz (on prendra $g = \SI{9,8}{N.kg^{-1}}$).

\textbf{3.} Montrer que la force de Lorentz ne travaille pas. En déduire que la norme de la vitesse de l'électron reste constante.

\textbf{4.} En appliquant la deuxième loi de Newton dans le repère de Frenet, établir que le mouvement de l'électron est circulaire uniforme et que le rayon de la trajectoire s'écrit $R = \dfrac{m v_0}{e B}$.

\textbf{5.} Calculer la valeur numérique de $R$.

\textbf{6.} Calculer la période $T$ du mouvement. Montrer qu'elle ne dépend pas de la vitesse $v_0$ ; commenter ce résultat en citant une application technologique.
\end{exercice}

\begin{exercice}[Type Bac : cadre dans un champ magnétique — le moteur électrique]
Un cadre rectangulaire $ABCD$ de $N = 100$ spires, de dimensions $AB = \SI{10}{cm}$ et $BC = \SI{5,0}{cm}$, est parcouru par un courant d'intensité $I = \SI{2,0}{A}$. Il est placé dans un champ magnétique uniforme $B = \SI{0,30}{T}$ et peut tourner sans frottement autour d'un axe vertical $(\Delta)$ passant par les milieux des côtés $AD$ et $BC$, perpendiculaire à $\vec{B}$. On note $\theta$ l'angle entre le vecteur normal $\vec{n}$ au plan du cadre et le vecteur $\vec{B}$.

\textbf{1.} Faire un schéma soigné du dispositif en vue de dessus, en représentant $\vec{B}$, $\vec{n}$, l'axe $(\Delta)$ et le sens du courant.

\textbf{2.} Exprimer puis calculer l'intensité des forces de Laplace s'exerçant sur chacun des quatre côtés du cadre. Préciser la direction et le sens de chacune.

\textbf{3.} Montrer que la somme vectorielle des forces de Laplace est nulle. Le cadre peut-il être mis en translation ?

\textbf{4.} Montrer que ces forces forment un couple dont le moment par rapport à $(\Delta)$ s'écrit $\mathcal{M} = N\, I\, S\, B \sin\theta$, où $S$ est l'aire d'une spire.

\textbf{5.} Calculer la valeur maximale $\mathcal{M}_{\max}$ de ce moment. Pour quelle position du cadre est-elle atteinte ?

\textbf{6.} Expliquer qualitativement comment ce dispositif, muni d'un collecteur qui inverse le sens du courant à chaque demi-tour, permet de réaliser un moteur à courant continu. Citer une application concrète au Cameroun.
\end{exercice}

\begin{exercice}[Type Bac : balance de Cotton et mesure d'un champ magnétique]
Pour mesurer l'intensité $B$ du champ magnétique régnant dans l'entrefer d'un électroaimant, un technicien du laboratoire d'un lycée de Yaoundé utilise une balance de Cotton. Une tige conductrice rigide, mobile autour d'un axe horizontal $(\Delta)$, comporte un tronçon rectiligne $MN$ de longueur $\ell = \SI{2,0}{cm}$ placé dans l'entrefer, perpendiculairement à $\vec{B}$. Le bras de levier du tronçon $MN$ par rapport à $(\Delta)$ vaut $d_1 = \SI{10}{cm}$. De l'autre côté de l'axe, à la distance $d_2 = \SI{15}{cm}$, on place une masse marquée $m$ pour rétablir l'équilibre de la balance lorsqu'un courant $I$ traverse la tige.

\textbf{Données :} $g = \SI{9,8}{N.kg^{-1}}$.

\textbf{1.} Faire un schéma de la balance en représentant l'axe $(\Delta)$, le tronçon $MN$, le champ $\vec{B}$, le courant $I$, la force de Laplace $\vec{F}$ et le poids $\vec{P}$ de la masse marquée.

\textbf{2.} Donner l'expression de l'intensité de la force de Laplace s'exerçant sur le tronçon $MN$. Pourquoi les autres tronçons de la tige ne contribuent-ils pas au moment mesuré ?

\textbf{3.} En écrivant la condition d'équilibre de la balance (théorème des moments), établir la relation :
\[ B = \dfrac{m\, g\, d_2}{I\, \ell\, d_1}. \]

\textbf{4.} Application numérique : pour $I = \SI{5,0}{A}$, l'équilibre est obtenu avec $m = \SI{1,63}{g}$. Calculer $B$.

\textbf{5.} Le technicien double l'intensité du courant. Que doit-il faire de la masse marquée pour conserver l'équilibre ? Justifier.

\textbf{6.} Citer une autre méthode de mesure d'un champ magnétique vue au laboratoire, et préciser le nom de l'appareil utilisé.
\end{exercice}

\newpage

\coursec{Corrigés des exercices}

\coursec{Exercices d'application et type Bac — Champ magnétique, force de Lorentz et force de Laplace}

\begin{corrige}[Exercice 1 — QCM : notions fondamentales]
\textbf{Énoncé :} Questions à choix multiples sur le champ magnétique d'un fil, d'un solénoïde, la force de Lorentz et l'interaction entre conducteurs.
\tcblower
\textbf{1. Réponse b).} Le champ magnétique créé par un fil rectiligne infini est, en tout point $M$, \cle{tangent au cercle} passant par $M$, centré sur le fil et situé dans le plan perpendiculaire au fil (lignes de champ circulaires). Son sens est donné par la règle du tire-bouchon ou de la main droite (pouce selon $I$). Il n'est ni radial (vers le fil) ni parallèle au fil.

\textbf{2. Réponse a).} À l'intérieur d'un solénoïde long, le champ est uniforme et vaut $\boxed{B = \mu_0\, n\, I}$. Vérification dimensionnelle : $[\mu_0 n I] = \mathrm{N.A^{-2} \times m^{-1} \times A} = \mathrm{N.A^{-1}.m^{-1}} = \mathrm{T}$ (car $1\ \mathrm{T} = 1\ \mathrm{N.A^{-1}.m^{-1}}$, d'après $F = I\ell B$). La réponse b) est homogène à $\mathrm{T.m^2}$ et la réponse c) à $\mathrm{T.A}$ : toutes deux sont fausses dimensionnellement.

\textbf{3. Réponses b) et c).} $\vec{F} = q\,\vec{v} \wedge \vec{B}$ est un produit vectoriel : $\vec{F}$ est donc \cle{perpendiculaire} à la fois à $\vec{v}$ et à $\vec{B}$. Étant toujours perpendiculaire au déplacement, son travail est nul : $W = \vec{F}\cdot\overrightarrow{dM} = 0$. La force de Lorentz magnétique ne modifie ni l'énergie cinétique ni la norme de la vitesse ; elle ne fait que courber la trajectoire.

\textbf{4. Réponse b).} Deux conducteurs parallèles parcourus par des courants de \cle{sens opposés se repoussent} (de même sens : ils s'attirent). Vérification par la règle de la main droite : le champ créé par le fil 1 à l'emplacement du fil 2, combiné au courant $I_2$ opposé à $I_1$, donne une force de Laplace dirigée vers l'extérieur.
\end{corrige}

\begin{corrige}[Exercice 2 — Vrai ou faux, à justifier]
\textbf{Énoncé :} Cinq affirmations sur les lignes de champ, la force de Lorentz, la force de Laplace, la règle des faces et le champ terrestre.
\tcblower
\textbf{1. Faux.} À l'\emph{extérieur} d'un aimant, les lignes de champ sont orientées du pôle \cle{nord vers le pôle sud} (elles vont du sud vers le nord à l'\emph{intérieur} de l'aimant, ce qui referme les lignes).

\textbf{2. Faux.} La force de Lorentz s'écrit $\vec{F} = q\,\vec{v}\wedge\vec{B}$. Si la charge est immobile, $\vec{v} = \vec{0}$, donc $\vec{F} = \vec{0}$ quelle que soit la valeur de $\vec{B}$. Un champ magnétique n'agit que sur des charges \emph{en mouvement}.

\textbf{3. Vrai.} C'est l'expression vectorielle de la \cle{force de Laplace} : $\vec{F} = I\,\vec{\ell}\wedge\vec{B}$, où $\vec{\ell}$ est un vecteur de norme $\ell$ orienté dans le sens du courant. En norme : $F = I\ell B \sin\alpha$, $\alpha$ étant l'angle entre le conducteur et $\vec{B}$.

\textbf{4. Vrai.} Règle des faces d'une bobine plate : la face où le courant semble tourner dans le \cle{sens trigonométrique} (antihoraire) est une face \cle{nord} ; vue dans le sens horaire, c'est une face sud. C'est cohérent avec la règle du tire-bouchon : un courant trigonométrique crée un champ sortant par cette face, or les lignes sortent par le nord.

\textbf{5. Faux.} Le champ magnétique terrestre n'est \cle{pas uniforme} à la surface du globe : son intensité varie d'environ \SI{2,5e-5}{T} à l'équateur à environ \SI{6e-5}{T} aux pôles, et sa direction (inclinaison) varie avec la latitude. Il n'est assimilable à un champ uniforme que localement, à l'échelle d'un laboratoire.
\end{corrige}

\begin{corrige}[Exercice 3 — Définitions et représentations]
\textbf{Énoncé :} Définir le vecteur champ magnétique, la ligne de champ, le dipôle magnétique ; énoncer la règle du tire-bouchon ; décrire le schéma d'un fil rectiligne ; donner les unités et faire l'analyse dimensionnelle.
\tcblower
\textbf{1. Définitions.}
\begin{itemize}
\item \cle{Vecteur champ magnétique} $\vec{B}$ en un point $M$ : grandeur vectorielle caractérisant l'état magnétique de l'espace en $M$. Sa direction et son sens sont ceux qu'indique le pôle nord d'une petite aiguille aimantée placée en $M$ ; sa norme se mesure en tesla (T) avec un teslamètre.
\item \cle{Ligne de champ} : courbe de l'espace tangente en chacun de ses points au vecteur $\vec{B}$ et orientée dans le sens de $\vec{B}$. Les lignes de champ sont denses là où le champ est intense et ne se croisent jamais.
\item \cle{Dipôle magnétique / moment magnétique} : un petit circuit parcouru par un courant (spire, aiguille aimantée, atome) se comporte comme un dipôle magnétique caractérisé par son moment magnétique $\vec{\mathcal{M}} = N I S\,\vec{n}$ (en $\mathrm{A.m^2}$), où $\vec{n}$ est le vecteur normal à la spire orienté par la règle du tire-bouchon.
\end{itemize}

\textbf{2. Règle du tire-bouchon (ou de la main droite) :} on place le pouce de la main droite dans le sens du courant $I$ le long du fil ; les doigts qui s'enroulent autour du fil indiquent alors le sens du vecteur champ magnétique $\vec{B}$ le long des lignes de champ circulaires.

\textbf{3. Schéma.} Le fil est horizontal, courant vers la droite. Les lignes de champ sont des cercles centrés sur le fil, dans des plans perpendiculaires au fil. En un point $M$ situé au-dessus du fil, la règle de la main droite donne $\vec{B}_M$ \emph{sortant} du plan de la figure (vers l'observateur) si le courant va vers la droite. Le schéma de l'énoncé illustre cette configuration : lignes circulaires concentriques et $\vec{B}_M$ tangent au cercle en $M$.

\textbf{4. Unités et analyse dimensionnelle.}
Unité SI du champ magnétique : le \cle{tesla} (T). Unité de $\mu_0$ : $\mathrm{N.A^{-2}}$ (ou $\mathrm{T.m.A^{-1}}$, ou $\mathrm{H.m^{-1}}$).
D'après la force de Laplace $F = I\ell B$, on tire $[B] = \dfrac{[F]}{[I][\ell]} = \mathrm{N.A^{-1}.m^{-1}}$, donc $1\ \mathrm{T} = 1\ \mathrm{N.A^{-1}.m^{-1}}$.
Vérifions :
\[ \left[\frac{\mu_0 I}{2\pi d}\right] = \frac{\mathrm{N.A^{-2}}\times\mathrm{A}}{\mathrm{m}} = \mathrm{N.A^{-1}.m^{-1}} = \mathrm{T}. \]
L'expression $\dfrac{\mu_0 I}{2\pi d}$ est bien homogène à un champ magnétique.
\end{corrige}

\begin{corrige}[Exercice 4 — Calculs directs]
\textbf{Énoncé :} Calculer le champ d'un fil, la force de Lorentz sur un proton, la force de Laplace sur un conducteur perpendiculaire puis incliné.
\tcblower
\textbf{1.} Champ créé par un fil rectiligne infini : $B = \dfrac{\mu_0 I}{2\pi d}$.
\[ B = \frac{4\pi \times 10^{-7} \times 5{,}0}{2\pi \times 0{,}020} = \frac{2 \times 10^{-7} \times 5{,}0}{0{,}020} = \frac{1{,}0 \times 10^{-6}}{0{,}020} = \SI{5,0e-5}{T} = \SI{50}{\micro T}. \]
C'est du même ordre de grandeur que la composante horizontale du champ terrestre : d'où l'expérience d'Ørsted.

\textbf{2.} Le proton se déplace perpendiculairement à $\vec{B}$ ($\sin\alpha = 1$) :
\[ F = |q|\,v\,B = 1{,}60\times 10^{-19} \times 1{,}0\times 10^{6} \times 0{,}50 = \SI{8,0e-14}{N}. \]

\textbf{3.} Conducteur perpendiculaire à $\vec{B}$ ($\sin\alpha = 1$), $\ell = 0{,}10\ \mathrm{m}$ :
\[ F = I\ell B = 4{,}0 \times 0{,}10 \times 0{,}20 = \SI{8,0e-2}{N} = \SI{80}{mN}. \]

\textbf{4.} Avec $\alpha = 30^{\circ}$ entre le conducteur et $\vec{B}$ :
\[ F = I\ell B \sin\alpha = 4{,}0 \times 0{,}10 \times 0{,}20 \times \sin 30^{\circ} = 0{,}080 \times 0{,}50 = \SI{4,0e-2}{N} = \SI{40}{mN}. \]
\begin{attention}
N'oublie jamais le facteur $\sin\alpha$ : la force de Laplace est maximale quand le conducteur est perpendiculaire à $\vec{B}$ et \emph{nulle} quand il est parallèle à $\vec{B}$.
\end{attention}
\end{corrige}

\begin{corrige}[Exercice 5 — Solénoïde d'un laboratoire de lycée]
\textbf{Énoncé :} Expression du champ dans un solénoïde long, calcul du courant pour un champ donné, description des lignes de champ, effet de la doublement de $n$ et $I$.
\tcblower
\textbf{1.} À l'intérieur d'un solénoïde \cle{long} (longueur très supérieure au rayon, $L \gg R$) et loin de ses extrémités, le champ magnétique est uniforme, dirigé selon l'axe, et vaut :
\[ B = \mu_0\, n\, I \]
où $n = N/L$ est le nombre de spires par mètre. Conditions de validité : solénoïde long, point situé à l'intérieur, loin des bords.

\textbf{2.} On isole $I$ :
\[ I = \frac{B}{\mu_0 n} = \frac{2{,}5\times 10^{-3}}{4\pi\times 10^{-7} \times 800} = \frac{2{,}5\times 10^{-3}}{1{,}005\times 10^{-3}} \approx \SI{2,5}{A}. \]

\textbf{3.} Voir la figure de l'énoncé : en coupe, les spires du haut sont traversées par un courant sortant ($\odot$) et celles du bas par un courant rentrant ($\otimes$) — ou l'inverse selon le branchement. Les lignes de champ sont des droites parallèles à l'axe à l'intérieur (champ uniforme), orientées de la face sud vers la face nord, et se referment à l'extérieur comme pour un aimant droit. La face d'où le courant semble tourner dans le sens trigonométrique est la face \cle{nord}.

\textbf{4.} $B = \mu_0 n I$ est proportionnelle au produit $nI$. Si $n \to 2n$ et $I \to 2I$, alors $B \to 4B$ : le champ est \textbf{multiplié par 4}, soit $B' = \SI{10}{mT}$.
\end{corrige}

\begin{corrige}[Exercice 6 — Rails de Laplace horizontaux]
\textbf{Énoncé :} Schéma des forces sur une barre mobile sur rails, calcul du courant pour une force donnée, nature du mouvement, inversion de la force.
\tcblower
\textbf{1.} Voir les schémas de l'énoncé. Vue de dessus : $\vec{B}$ vertical vers le haut est représenté sortant du plan ($\odot$) ; le courant $I$ traverse la barre $MN$ ; la règle de la main droite (ou des trois doigts) appliquée à $\vec{F} = I\vec{\ell}\wedge\vec{B}$ donne $\vec{F}_L$ horizontale, perpendiculaire à la barre, dirigée le long des rails.

\textbf{2.} La barre est perpendiculaire à $\vec{B}$ ($\sin\alpha = 1$) : $F = I\ell B$, d'où
\[ I = \frac{F}{\ell B} = \frac{0{,}050}{0{,}12 \times 0{,}40} = \frac{0{,}050}{0{,}048} \approx \SI{1,0}{A}. \]

\textbf{3.} Bilan des forces sur la barre : poids $\vec{P} = m\vec{g}$ (vertical), réactions des rails $\vec{R}$ (verticales, sans frottement), force de Laplace $\vec{F}_L$ (horizontale). Les forces verticales se compensent ; la résultante est donc $\vec{F}_L$, constante en norme, direction et sens.
Deuxième loi de Newton : $m\vec{a} = \vec{F}_L$, donc $\vec{a}$ constante : le mouvement est \cle{rectiligne uniformément accéléré} (partant du repos), avec
\[ a = \frac{F}{m} = \frac{0{,}050}{0{,}020} = \SI{2,5}{m.s^{-2}}. \]

\textbf{4.} Le sens de $\vec{F}_L = I\vec{\ell}\wedge\vec{B}$ change si l'on inverse \emph{un seul} des deux vecteurs :
\begin{itemize}
\item soit on inverse le sens du courant $I$ (en permutant les bornes du générateur) ;
\item soit on inverse le sens de $\vec{B}$ (en retournant l'aimant ou l'électroaimant).
\end{itemize}
La règle de la main droite le confirme : inverser l'index ($I\vec{\ell}$) ou le majeur ($\vec{B}$) retourne le pouce ($\vec{F}$). Attention : inverser \emph{les deux} laisse la force inchangée.
\end{corrige}

\begin{corrige}[Exercice 7 — Interaction entre deux câbles de transport d'énergie]
\textbf{Énoncé :} Champ créé par un fil, force sur l'autre, nature de l'interaction (attraction/répulsion), schémas, définition de l'ampère.
\tcblower
\textbf{1.} Le fil 1 crée à la distance $d$ un champ d'intensité :
\[ B_1 = \frac{\mu_0 I_1}{2\pi d} = \frac{4\pi\times 10^{-7}\times 20}{2\pi \times 0{,}10} = \frac{2\times 10^{-7}\times 20}{0{,}10} = \SI{4,0e-5}{T}. \]
Direction : tangent au cercle centré sur le fil 1 passant par le fil 2 ; sens donné par la règle du tire-bouchon.

\textbf{2.} Le fil 2, parcouru par $I_2$, est plongé dans $\vec{B}_1$ qui lui est perpendiculaire. La force de Laplace sur une longueur $\ell$ vaut $F = I_2 \ell B_1$, soit par unité de longueur :
\[ \frac{F}{\ell} = I_2 B_1 = \frac{\mu_0 I_1 I_2}{2\pi d} = \frac{4\pi\times 10^{-7}\times 20\times 30}{2\pi\times 0{,}10} = \frac{2\times 10^{-7}\times 600}{0{,}10} = \SI{1,2e-3}{N.m^{-1}}. \]
D'après le principe des actions réciproques, le fil 1 subit une force de même intensité, de sens opposé.

\textbf{3.} Courants de \cle{même sens} : la force est \textbf{attractive}. Justification (figure de l'énoncé, courants sortants) : en $M_2$, $\vec{B}_1$ est tangent au cercle centré en $M_1$, orienté vers le haut ; $\vec{F}_{1/2} = I_2\vec{\ell}\wedge\vec{B}_1$ est alors dirigée vers le fil 1 (règle de la main droite). Symétriquement, $\vec{F}_{2/1}$ est dirigée vers le fil 2 : les deux fils s'attirent.

\textbf{4.} Voir la figure de l'énoncé : $\vec{B}_1$ (tangent au cercle centré sur le fil 1), $\vec{B}_2$ (tangent au cercle centré sur le fil 2), et les forces $\vec{F}_{1/2}$ et $\vec{F}_{2/1}$ dirigées l'une vers l'autre.

\textbf{5.} Si les courants sont de \cle{sens opposés}, le sens de l'un des produits vectoriels change : les forces deviennent \textbf{répulsives}, mais leur norme est inchangée car elle ne dépend que des intensités et de $d$ : $F/\ell = \SI{1,2e-3}{N.m^{-1}}$.
\begin{savaistu}
C'est exactement cette expérience (Ampère, 1820) qui fonde la \cle{définition historique de l'ampère} : deux fils parallèles distants de 1 m, parcourus par des courants de 1 A, s'attirent avec une force de $2\times 10^{-7}\ \mathrm{N}$ par mètre.
\end{savaistu}
\end{corrige}

\begin{corrige}[Exercice 8 — Superposition du champ d'un fil et du champ terrestre]
\textbf{Énoncé :} Champ d'un fil vertical au point M au nord, superposition avec le champ terrestre horizontal, déviation de l'aiguille aimantée, expérience d'Ørsted.
\tcblower
\textbf{1.} Vue de dessus, courant ascendant (sortant de la page, $\odot$). Règle du tire-bouchon : les lignes de champ sont des cercles orientés dans le sens trigonométrique. Au point $M$ situé au \textbf{nord} du fil, la tangente au cercle est dirigée vers l'\textbf{est}. Donc $\vec{B}_{\text{fil}}$ : direction est-ouest, sens \textbf{vers l'est}, norme :
\[ B_{\text{fil}} = \frac{\mu_0 I}{2\pi d} = \frac{4\pi\times 10^{-7}\times 10}{2\pi\times 0{,}050} = \frac{2\times 10^{-7}\times 10}{0{,}050} = \SI{4,0e-5}{T}. \]

\textbf{2.} $\vec{B}_h$ est dirigé vers le nord et $\vec{B}_{\text{fil}}$ vers l'est : les deux vecteurs sont bien \textbf{perpendiculaires}. Le champ résultant $\vec{B} = \vec{B}_h + \vec{B}_{\text{fil}}$ a pour norme (Pythagore) :
\[ B = \sqrt{B_h^2 + B_{\text{fil}}^2} = \sqrt{(2{,}0\times 10^{-5})^2 + (4{,}0\times 10^{-5})^2} = 10^{-5}\sqrt{4 + 16} = \sqrt{20}\times 10^{-5} \approx \SI{4,5e-5}{T}. \]

\textbf{3.} L'angle $\alpha$ entre $\vec{B}$ et la direction du nord vérifie :
\[ \tan\alpha = \frac{B_{\text{fil}}}{B_h} = \frac{4{,}0\times 10^{-5}}{2{,}0\times 10^{-5}} = 2{,}0 \quad\Rightarrow\quad \alpha = \arctan(2{,}0) \approx 63^{\circ}. \]
Le champ résultant pointe vers le nord-est, à environ $63^{\circ}$ à l'est du nord.

\textbf{4.} L'aiguille aimantée, mobile autour d'un axe vertical, s'aligne sur la composante horizontale du champ magnétique total en $M$ : elle indique donc la direction de $\vec{B}$, c'est-à-dire environ $63^{\circ}$ à l'est du nord géographique, et non plus le nord magnétique terrestre. Si l'on coupe le courant, elle retrouve sa direction initiale.
Ce dispositif reproduit l'\textbf{expérience d'Ørsted (1820)} : la déviation d'une aiguille aimantée par un courant électrique, première preuve expérimentale du lien entre électricité et magnétisme.
\end{corrige}

\begin{corrige}[Exercice 9 — Type Bac : électron dans un champ magnétique uniforme]
\textbf{Énoncé :} Électron injecté perpendiculairement à $\vec{B}$ : force, négligeabilité du poids, conservation de l'énergie, nature de la trajectoire, rayon, période, isochronisme et cyclotron.
\tcblower
\textbf{1. (2 pts)} Schéma : $\vec{B}$ sortant du plan ($\odot$), $\vec{v}_0$ horizontale vers la droite en $O$.
Force de Lorentz : $\vec{F} = q\,\vec{v}_0\wedge\vec{B} = -e\,\vec{v}_0\wedge\vec{B}$.
Le produit $\vec{v}_0\wedge\vec{B}$ est vertical vers le haut (règle de la main droite) ; comme $q = -e < 0$, $\vec{F}$ est \textbf{verticale, vers le bas}, appliquée en $O$, de norme :
\[ F = e v_0 B = 1{,}60\times 10^{-19}\times 2{,}0\times 10^{7}\times 1{,}0\times 10^{-3} = \SI{3,2e-15}{N}. \]

\textbf{2. (1 pt)} Poids de l'électron :
\[ P = mg = 9{,}11\times 10^{-31}\times 9{,}8 \approx \SI{8,9e-30}{N}. \]
Rapport : $\dfrac{F}{P} \approx \dfrac{3{,}2\times 10^{-15}}{8{,}9\times 10^{-30}} \approx 3{,}6\times 10^{14} \gg 1$. Le poids est environ $10^{14}$ fois plus faible que la force magnétique : il est bien négligeable.

\textbf{3. (2 pts)} Puissance de la force de Lorentz : $\mathcal{P} = \vec{F}\cdot\vec{v} = q(\vec{v}\wedge\vec{B})\cdot\vec{v} = 0$ car $\vec{v}\wedge\vec{B} \perp \vec{v}$. La force ne travaille donc pas : $W = 0$.
Théorème de l'énergie cinétique entre deux instants quelconques : $\Delta E_c = W = 0$, donc $E_c = \tfrac{1}{2}mv^2$ est constante, et comme $m$ est constante, la norme $v = v_0$ est \textbf{constante} : le mouvement est uniforme.

\textbf{4. (3 pts)} Système : électron. Référentiel galiléen du laboratoire. Seule force : $\vec{F}$ (poids négligé).
Dans le repère de Frenet $(\vec{\tau}, \vec{n})$ : $\vec{a} = \dfrac{dv}{dt}\vec{\tau} + \dfrac{v^2}{\rho}\vec{n}$.
Comme $v$ est constante, $\dfrac{dv}{dt} = 0$ : l'accélération est purement normale, $\vec{a} = \dfrac{v_0^2}{\rho}\vec{n}$.
Deuxième loi de Newton : $m\vec{a} = \vec{F}$. Or $\vec{F}$ est perpendiculaire à $\vec{v}$, donc colinéaire à $\vec{n}$, de norme $ev_0B$. En projection sur $\vec{n}$ :
\[ m\,\frac{v_0^2}{\rho} = e v_0 B \quad\Rightarrow\quad \rho = \frac{m v_0}{e B} = \text{constante}. \]
Le rayon de courbure est constant et le mouvement est uniforme : la trajectoire est un \textbf{cercle} de rayon $R = \dfrac{mv_0}{eB}$, parcouru à vitesse constante : mouvement \cle{circulaire uniforme}.

\textbf{5. (1 pt)}
\[ R = \frac{9{,}11\times 10^{-31}\times 2{,}0\times 10^{7}}{1{,}60\times 10^{-19}\times 1{,}0\times 10^{-3}} = \frac{1{,}822\times 10^{-23}}{1{,}60\times 10^{-22}} \approx \SI{0,114}{m} \approx \SI{11}{cm}. \]

\textbf{6. (2 pts)} La période est la durée d'un tour complet :
\[ T = \frac{2\pi R}{v_0} = \frac{2\pi}{v_0}\times\frac{mv_0}{eB} = \frac{2\pi m}{eB}. \]
$v_0$ se simplifie : $T$ \textbf{ne dépend pas de la vitesse} (ni du rayon), mais seulement de $m$, $e$ et $B$.
\[ T = \frac{2\pi\times 9{,}11\times 10^{-31}}{1{,}60\times 10^{-19}\times 1{,}0\times 10^{-3}} \approx \SI{3,6e-8}{s} = \SI{36}{ns}. \]
\textbf{Commentaire :} cette propriété d'\cle{isochronisme} est exploitée dans le \textbf{cyclotron} : un champ électrique alternatif de fréquence fixe $f = 1/T$ accélère les particules à chaque demi-tour, quelle que soit leur énergie (régime non relativiste). Applications : production de radio-isotopes pour l'imagerie médicale et la radiothérapie (cancérologie), physique nucléaire.

\textbf{Points de vigilance (Bac) :}
\begin{itemize}
\item Toujours justifier le sens de $\vec{F}$ en précisant le \textbf{signe de la charge} (ici $q<0$ : force opposée à $\vec{v}\wedge\vec{B}$).
\item Citer explicitement le \textbf{théorème de l'énergie cinétique} pour conclure que $v$ est constante — ne pas l'affirmer sans preuve.
\item Dans le repère de Frenet, écrire les deux composantes de $\vec{a}$ avant d'exploiter $\dfrac{dv}{dt}=0$.
\item Conditions d'application : champ uniforme, $\vec{v}_0 \perp \vec{B}$, poids négligé (à justifier numériquement), régime non relativiste.
\end{itemize}
\end{corrige}

\begin{corrige}[Exercice 10 — Type Bac : cadre dans un champ magnétique — le moteur électrique]
\textbf{Énoncé :} Cadre rectangulaire à $N$ spires dans un champ uniforme, axe de rotation vertical, forces de Laplace, moment du couple, position de couple maximal, rôle du collecteur, applications.
\tcblower
\textbf{1. (2 pts)} Schéma : axe $(\Delta)$ vertical passant par les milieux de $AD$ et $BC$ ; $\vec{B}$ horizontal ; $\vec{n}$ normal au plan du cadre faisant l'angle $\theta$ avec $\vec{B}$ ; courant circulant de $A$ vers $B$, $B$ vers $C$, etc. Voir la figure de l'énoncé.

\textbf{2. (3 pts)} Pour $N$ spires, chaque force est multipliée par $N$.
\begin{itemize}
\item \textbf{Côtés $AB$ et $CD$} (longueur $AB = 0{,}10$ m) : ces côtés sont horizontaux. L'angle entre $\vec{\ell}$ (selon $AB$ ou $CD$) et $\vec{B}$ (horizontal) varie avec $\theta$. Les forces $\vec{F}_{AB}$ et $\vec{F}_{CD}$ sont verticales (perpendiculaires au plan contenant le côté et $\vec{B}$), opposées l'une à l'autre, et leur droite d'action \textbf{passe par l'axe $(\Delta)$} : leur moment par rapport à $(\Delta)$ est nul. Elles ne participent pas à la rotation.
\item \textbf{Côtés $AD$ et $BC$} (longueur $\ell = BC = 0{,}050$ m) : ces côtés restent toujours \textbf{perpendiculaires à $\vec{B}$} (verticaux contre horizontal), quelle que soit la position du cadre. Chacun subit une force de Laplace d'intensité :
\[ F = N I \ell B = 100\times 2{,}0\times 0{,}050\times 0{,}30 = \SI{3,0}{N}. \]
Direction : horizontale, \textbf{perpendiculaire à $\vec{B}$} (donc fixe dans le référentiel du laboratoire), orientée selon $\pm\vec{u}_y$ si $\vec{B}=\vec{u}_x$. Sens : opposés sur $AD$ et sur $BC$ (courants de sens opposés dans les deux côtés).
\end{itemize}

\textbf{3. (1 pt)} Les forces s'exerçant sur $AD$ et $BC$ sont directement opposées ($\vec{F}_{AD} = -\vec{F}_{BC}$) ; celles sur $AB$ et $CD$ également. Donc :
\[ \sum \vec{F} = \vec{0}. \]
La résultante étant nulle, le centre de masse du cadre n'est pas accéléré : le cadre \textbf{ne peut pas être mis en translation}. Il subit un \cle{couple} (résultante nulle mais moment non nul).

\textbf{4. (3 pts)} Seules les forces sur $AD$ et $BC$ ont un moment non nul par rapport à $(\Delta)$ (celles sur $AB$ et $CD$ coupent l'axe). Ces deux forces, d'intensité $F = NI\ell B$, sont horizontales, parallèles et de sens opposés. Leurs droites d'action sont des droites verticales passant par $AD$ et $BC$. La distance entre ces droites (bras de levier du couple), mesurée perpendiculairement à la direction des forces (qui est fixe, perpendiculaire à $\vec{B}$), vaut :
\[ b = AB\,\sin\theta \]
(projection de la largeur $AB$ sur la direction perpendiculaire aux forces). Le moment du couple vaut :
\[ \mathcal{M} = F\times b = NI\ell B \times AB\sin\theta = NI\,(\ell\times AB)\,B\sin\theta. \]
Or $S = \ell\times AB = BC\times AB$ est l'aire d'une spire :
\[ \boxed{\mathcal{M} = N\,I\,S\,B\,\sin\theta}. \]
Analyse dimensionnelle : $[NISB] = \mathrm{A\times m^2\times T} = \mathrm{A\times m^2\times N.A^{-1}.m^{-1}} = \mathrm{N.m}$ : homogène à un moment. 

\textbf{5. (1 pt)} $\mathcal{M}$ est maximal pour $\sin\theta = 1$, c'est-à-dire $\theta = 90^{\circ}$ : le \textbf{plan du cadre est parallèle à $\vec{B}$} (le vecteur normal $\vec{n}$ est perpendiculaire à $\vec{B}$).
\[ \mathcal{M}_{\max} = NISB = 100\times 2{,}0\times (0{,}10\times 0{,}050)\times 0{,}30 = 100\times 2{,}0\times 5{,}0\times 10^{-3}\times 0{,}30 = \SI{0,30}{N.m}. \]

\textbf{6. (2 pts)} \textbf{Fonctionnement du moteur à courant continu :} le couple $\mathcal{M} = NISB\sin\theta$ fait tourner le cadre. Mais si le courant gardait le même sens, après un demi-tour ($\theta$ passant de $90^{\circ}$ à $270^{\circ}$), $\sin\theta$ changerait de signe et le couple s'inverserait : le cadre oscillerait puis s'immobiliserait en position d'équilibre stable ($\theta = 0$, plan du cadre perpendiculaire à $\vec{B}$).
Le \cle{collecteur} (cylindre fendu en deux demi-anneaux solidaires du cadre, frotté par deux balais reliés au générateur) \textbf{inverse le sens du courant à chaque demi-tour}, exactement au passage par la position $\theta = 0$ où le couple s'annule (l'inertie fait franchir cette position). Ainsi le produit $I\sin\theta$ garde toujours le même signe : le couple reste toujours moteur et le cadre tourne en continu.
\textbf{Applications au Cameroun :} moteurs de pompes à eau de forage (électrification rurale, adduction d'eau), ventilateurs, démarreurs de mototaxis et véhicules, moteurs de machines agroalimentaires (moulins à maïs, décortiqueuses), perceuses et outillage des menuiseries de Douala et Yaoundé.

\textbf{Points de vigilance (Bac) :}
\begin{itemize}
\item Ne pas oublier le facteur $N$ (nombre de spires) dans les forces et le moment.
\item Justifier que les forces sur $AB$ et $CD$ ont un moment nul (droites d'action passant par l'axe).
\item Distinguer clairement \emph{résultante nulle} (pas de translation) et \emph{moment non nul} (rotation) : c'est la définition d'un couple.
\item Pour $\mathcal{M}_{\max}$, citer la position : plan du cadre parallèle à $\vec{B}$.
\item \textbf{Attention :} les forces sur les côtés verticaux ($AD$, $BC$) sont \textbf{fixes en direction} (perpendiculaires à $\vec{B}$), elles ne tournent pas avec le cadre ; c'est le bras de levier $b = AB\sin\theta$ qui varie.
\end{itemize}
\end{corrige}

\begin{corrige}[Exercice 11 — Type Bac : balance de Cotton et mesure d'un champ magnétique]
\textbf{Énoncé :} Balance de Cotton : schéma, forces sur le tronçon $MN$, équilibre des moments, expression de $B$, application numérique, proportionnalité $m(I)$, autre méthode de mesure (sonde de Hall).
\tcblower
\textbf{1. (2 pts)} Schéma en vue de côté : axe horizontal $(\Delta)$ ; bras de longueur $d_1$ portant le tronçon $MN$ plongé dans l'entrefer où $\vec{B}$ est horizontal (perpendiculaire au plan de la figure, ou vertical selon la disposition) ; courant $I$ dans $MN$ ; force de Laplace $\vec{F}$ verticale ; de l'autre côté, à la distance $d_2$, la masse marquée $m$ de poids $\vec{P} = m\vec{g}$. Voir la figure de l'énoncé.

\textbf{2. (2 pts)} Le tronçon $MN$ est perpendiculaire à $\vec{B}$ ($\sin\alpha = 1$) :
\[ F = I\,\ell\,B. \]
Les autres tronçons de la tige ne contribuent pas au moment mesuré car :
\begin{itemize}
\item les portions de conducteur situées \textbf{hors de l'entrefer} sont dans un champ quasi nul (le champ de l'électroaimant est confiné dans l'entrefer) ;
\item les tronçons éventuellement \textbf{parallèles à $\vec{B}$} subissent une force nulle ($F = I\ell B\sin 0 = 0$) ;
\item les forces s'exerçant sur les tronçons dont la droite d'action \textbf{passe par l'axe $(\Delta)$} ont un bras de levier nul, donc un moment nul.
\end{itemize}
Seul le moment de la force de Laplace sur $MN$ est mesuré.

\textbf{3. (3 pts)} À l'équilibre de la balance (solide en rotation autour d'un axe fixe), le théorème des moments s'écrit : la somme algébrique des moments par rapport à $(\Delta)$ est nulle.
\begin{itemize}
\item Moment de la force de Laplace (bras de levier $d_1$, force perpendiculaire au bras) : $\mathcal{M}_F = F\,d_1 = I\ell B\,d_1$.
\item Moment du poids de la masse marquée (bras de levier $d_2$) : $\mathcal{M}_P = m g\,d_2$, de sens opposé.
\end{itemize}
Condition d'équilibre :
\[ I\ell B\,d_1 = m g\,d_2 \quad\Rightarrow\quad \boxed{B = \frac{m\,g\,d_2}{I\,\ell\,d_1}}. \]
Analyse dimensionnelle : $\left[\dfrac{mgd_2}{I\ell d_1}\right] = \dfrac{\mathrm{N\times m}}{\mathrm{A\times m\times m}} = \dfrac{\mathrm{N}}{\mathrm{A.m}} = \mathrm{T}$. Correct.

\textbf{4. (1 pt)} Application numérique avec $m = 1{,}63\times 10^{-3}$ kg, $I = 5{,}0$ A, $\ell = 0{,}020$ m, $d_1 = 0{,}10$ m, $d_2 = 0{,}15$ m :
\[ B = \frac{1{,}63\times 10^{-3}\times 9{,}8\times 0{,}15}{5{,}0\times 0{,}020\times 0{,}10} = \frac{2{,}396\times 10^{-3}}{1{,}0\times 10^{-2}} \approx \SI{0,24}{T}. \]
Valeur typique du champ dans l'entrefer d'un électroaimant de laboratoire.

\textbf{5. (1 pt)} À l'équilibre, $B$ étant fixé par l'électroaimant, la relation $m = \dfrac{B I \ell d_1}{g d_2}$ montre que $m$ est \textbf{proportionnelle à $I$}. Si le technicien double le courant ($I' = 2I$), il doit donc \textbf{doubler la masse marquée} : $m' = 2m = \SI{3,26}{g}$.

\textbf{6. (1 pt)} Autre méthode de mesure : la \cle{sonde de Hall}, associée à un \textbf{teslamètre}. Une plaquette semiconductrice parcourue par un courant et placée perpendiculairement au champ développe une tension de Hall $U_H$ proportionnelle à $B$ ; l'appareil affiche directement $B$ en tesla. (On peut aussi citer la méthode de la bobine exploratrice associée à un fluxmètre.)

\textbf{Points de vigilance (Bac) :}
\begin{itemize}
\item Bien préciser que la condition d'équilibre utilisée est le \textbf{théorème des moments} (somme des moments nulle), et non $\sum\vec{F} = \vec{0}$ seul.
\item Convertir systématiquement les longueurs en mètres et la masse en kilogrammes avant le calcul.
\item Citer la condition d'application de $F = I\ell B$ : tronçon \textbf{perpendiculaire} à $\vec{B}$ (sinon $F = I\ell B\sin\alpha$).
\end{itemize}
\end{corrige}

\newpage

\coursec{Fiche bilan}

\begin{popfigure}
\begin{center}
\begin{tikzpicture}[mindmap, grow cyclic, every node/.style={concept, font=\small},
  root concept/.append style={concept color=popblue, font=\bfseries},
  level 1/.append style={level distance=3.4cm, sibling angle=60, font=\scriptsize},
  level 2/.append style={level distance=2.2cm, sibling angle=45, font=\tiny}]
\node[root concept]{Champ magnétique, forces de Lorentz et de Laplace}
  child[concept color=poporange]{ node {Aimants et champ $\vec{B}$}
    child { node {lignes de champ} }
    child { node {champ terrestre} }
  }
  child[concept color=popgreen]{ node {Champ créé par un courant}
    child { node {fil, spire, solénoïde} }
    child { node {$B=\mu_0 n I$} }
  }
  child[concept color=poppurple]{ node {Force de Lorentz}
    child { node {$\vec{F}=q\,\vec{v}\wedge\vec{B}$} }
    child { node {trajectoire circulaire} }
  }
  child[concept color=poppink]{ node {Force de Laplace}
    child { node {$\vec{F}=I\,\vec{\ell}\wedge\vec{B}$} }
    child { node {règle de la main droite} }
  }
  child[concept color=popgold]{ node {Applications}
    child { node {haut-parleur} }
    child { node {moteur électrique} }
  }
  child[concept color=popblue!60]{ node {Deux fils parallèles}
    child { node {définition de l'ampère} }
  };
\end{tikzpicture}
\end{center}
\legende{Carte mentale de la leçon : les six grandes branches à maîtriser.}
\end{popfigure}

\begin{bilan}
\textbf{1. Définitions clés}
\begin{itemize}
  \item \cle{Champ magnétique} $\vec{B}$ : grandeur vectorielle caractérisant l'espace où s'exercent des actions magnétiques ; unité : le \cle{tesla} (T). Il est orienté du pôle \textbf{N}ord vers le pôle \textbf{S}ud d'un aimant (à l'extérieur).
  \item \cle{Lignes de champ} : courbes tangentes en chaque point à $\vec{B}$, orientées dans le sens de $\vec{B}$ ; elles ne se croisent jamais et sortent par le nord d'un aimant.
  \item \cle{Champ magnétique terrestre} : au Cameroun (proche de l'équateur), sa composante horizontale vaut environ $\SI{3e-5}{T}$.
  \item \cle{Solénoïde} : bobine longue dont le champ intérieur est quasi uniforme.
\end{itemize}

\textbf{2. Formules à connaître par c\oe ur}
\begin{center}
\begin{tabularx}{\linewidth}{|l|X|X|}
\hline
\rowcolor{popblueL} \textbf{Grandeur} & \textbf{Formule} & \textbf{Unités / remarques} \\
\hline
Champ d'un solénoïde & $B=\mu_0\,n\,I$ & $\mu_0=4\pi\times10^{-7}$ S.I. ; $n$ : spires par mètre ($\mathrm{m^{-1}}$) \\
\hline
Force de Lorentz & $\vec{F}=q\,\vec{v}\wedge\vec{B}$ & $F=|q|\,v\,B\,|\sin\alpha|$ ; en newtons (N) \\
\hline
Force de Laplace & $\vec{F}=I\,\vec{\ell}\wedge\vec{B}$ & $F=I\,\ell\,B\,|\sin\alpha|$ ; $\vec{\ell}$ orienté dans le sens de $I$ \\
\hline
Cas $\vec{v}\perp\vec{B}$ (ou $\vec{\ell}\perp\vec{B}$) & $F=|q|vB$ \quad ; \quad $F=I\ell B$ & $\alpha=90^{\circ}$, $\sin\alpha=1$ \\
\hline
Force entre deux fils parallèles & $F=\dfrac{\mu_0\,I_1 I_2\,\ell}{2\pi d}$ & attraction si courants de même sens, répulsion sinon \\
\hline
Mouvement d'une charge $\perp\vec{B}$ & $R=\dfrac{mv}{|q|B}$ & trajectoire circulaire uniforme \\
\hline
\end{tabularx}
\end{center}

\textbf{3. Méthodes express}
\begin{itemize}
  \item \cle{Sens de $\vec{B}$ créé par un courant} : règle du tire-bouchon (ou du bonhomme d'Ampère) — le courant entre par les pieds, sort par la tête, $\vec{B}$ sort par la face nord.
  \item \cle{Sens de la force} : règle des trois doigts de la main droite (ou du bonhomme d'Ampère) ; pour $q<0$, la force de Lorentz est \textbf{opposée} à $\vec{v}\wedge\vec{B}$.
  \item \cle{Analyse dimensionnelle} : $1~\mathrm{T}=1~\mathrm{N\,A^{-1}\,m^{-1}}=1~\mathrm{kg\,s^{-1}\,C^{-1}}$.
  \item \cle{Définition de l'ampère} : courant qui, dans deux fils parallèles distants de $\SI{1}{m}$, produit une force de $\num{2e-7}~\mathrm{N}$ par mètre.
\end{itemize}
\end{bilan}

\begin{aretenir}[Checklist avant le Bac]
\begin{itemize}
  \item[$\square$] Je sais orienter $\vec{B}$ autour d'un aimant (du N vers le S à l'extérieur) et tracer des lignes de champ.
  \item[$\square$] Je sais déterminer le sens de $\vec{B}$ créé par un fil, une spire ou un solénoïde (règle du tire-bouchon).
  \item[$\square$] Je connais par c\oe ur $B=\mu_0 n I$ pour un solénoïde et je sais convertir $n$ en $\mathrm{m^{-1}}$.
  \item[$\square$] Je sais écrire $\vec{F}=q\,\vec{v}\wedge\vec{B}$ et calculer $F=|q|vB|\sin\alpha|$ avec les unités SI.
  \item[$\square$] Je n'oublie pas d'inverser le sens de la force de Lorentz quand la charge est \textbf{négative}.
  \item[$\square$] Je sais écrire $\vec{F}=I\,\vec{\ell}\wedge\vec{B}$ et appliquer la règle de la main droite.
  \item[$\square$] Je vérifie que les longueurs sont en mètres, les courants en ampères, les champs en teslas avant tout calcul.
  \item[$\square$] Je sais que deux courants de même sens s'attirent, de sens opposés se repoussent.
  \item[$\square$] Je sais retrouver $R=\dfrac{mv}{|q|B}$ pour une charge perpendiculaire à $\vec{B}$ (force centripète).
  \item[$\square$] Je sais citer une application : haut-parleur, moteur électrique, définition de l'ampère.
  \item[$\square$] Je représente toujours les vecteurs $\vec{F}$, $\vec{v}$, $\vec{B}$ sur un schéma clair avant de conclure.
  \item[$\square$] Je donne chaque résultat numérique avec son unité et un nombre raisonnable de chiffres significatifs.
\end{itemize}
\end{aretenir}

\findecours

\end{document}
